\documentclass[11pt]{article}
\usepackage[a4paper,margin=24mm]{geometry}
\usepackage[T1]{fontenc}
\usepackage{lmodern,microtype}
\usepackage{amsmath,amssymb,amsfonts,amsthm,mathtools,bm,mathrsfs}
\usepackage{esint}
\usepackage{enumitem}
\usepackage{array,booktabs,longtable}
\usepackage{tikz,float}
\usetikzlibrary{arrows.meta,calc,patterns}
\usepackage{hyperref}
\usepackage{aliascnt}
\usepackage[nameinlink,capitalise]{cleveref}
\allowdisplaybreaks
\setlength{\parskip}{1pt}
\setlength{\parindent}{15pt}

% Theorem-like environments share one counter, numbered within sections.
\newtheorem{theorem}{Theorem}[section]
\newaliascnt{proposition}{theorem}
\newtheorem{proposition}[proposition]{Proposition}
\aliascntresetthe{proposition}
\newaliascnt{lemma}{theorem}
\newtheorem{lemma}[lemma]{Lemma}
\aliascntresetthe{lemma}
\newaliascnt{corollary}{theorem}
\newtheorem{corollary}[corollary]{Corollary}
\aliascntresetthe{corollary}
\newaliascnt{claim}{theorem}
\newtheorem{claim}[claim]{Claim}
\aliascntresetthe{claim}
\theoremstyle{definition}
\newaliascnt{definition}{theorem}
\newtheorem{definition}[definition]{Definition}
\aliascntresetthe{definition}
\theoremstyle{remark}
\newaliascnt{remark}{theorem}
\newtheorem{remark}[remark]{Remark}
\aliascntresetthe{remark}
\numberwithin{equation}{section}
\crefname{theorem}{Theorem}{Theorems}
\Crefname{theorem}{Theorem}{Theorems}
\crefname{proposition}{Proposition}{Propositions}
\Crefname{proposition}{Proposition}{Propositions}
\crefname{lemma}{Lemma}{Lemmas}
\Crefname{lemma}{Lemma}{Lemmas}
\crefname{corollary}{Corollary}{Corollaries}
\Crefname{corollary}{Corollary}{Corollaries}
\crefname{definition}{Definition}{Definitions}
\Crefname{definition}{Definition}{Definitions}
\crefname{claim}{Claim}{Claims}
\crefname{remark}{Remark}{Remarks}

\newcommand{\T}{\mathbb T}
\newcommand{\R}{\mathbb R}
\newcommand{\C}{\mathbb C}
\newcommand{\Z}{\mathbb Z}
\newcommand{\E}{\mathbb E}
\newcommand{\Pp}{\mathbb P}
\newcommand{\F}{\mathcal F}
\newcommand{\Q}{\mathcal Q}
\newcommand{\pr}{\mathcal P}
\newcommand{\tr}{\operatorname{tr}}
\newcommand{\divv}{\operatorname{div}}
\newcommand{\dist}{\operatorname{dist}}
\newcommand{\diam}{\operatorname{diam}}
\newcommand{\supp}{\operatorname{supp}}
\newcommand{\Lip}{\operatorname{Lip}}
\newcommand{\one}{\mathbf 1}
\newcommand{\pre}{\preceq}
\newcommand{\ps}{\psi}
\newcommand{\norm}[1]{\left\|#1\right\|}
\newcommand{\abs}[1]{\left|#1\right|}
\newcommand{\dd}{\,\mathrm d}
\newcommand{\wideD}{\widehat D}
\newcommand{\Id}{\mathrm{Id}}
\newcommand{\Var}{\operatorname{Var}}
\newcommand{\HS}{\mathrm{HS}}
\newcommand{\op}{\mathrm{op}}
\newcommand{\cF}{\mathcal F}
\newcommand{\cT}{\mathcal T}
\newcommand{\cM}{\mathcal M}
\newcommand{\cP}{\mathcal P}
\newcommand{\cB}{\mathcal B}
\newcommand{\set}[1]{\left\{#1\right\}}
\newcommand{\pa}[1]{\left(#1\right)}
\newcommand{\br}[1]{\left[#1\right]}
\newcommand{\ip}[2]{\left\langle #1,#2\right\rangle}
\renewcommand{\fint}{\mathop{\rlap{-}\!\int}}
\newcommand{\circleavg}[1]{\frac1{2\pi}\int_0^{2\pi}#1\,\mathrm d\theta}

\title{The Ultimate Matching Theorem}
\hypersetup{pdfauthor={Michael Goldman, Francesco Grotto, Dario Trevisan}}
\author{%
Michael Goldman\textsuperscript{1}\qquad
Francesco Grotto\textsuperscript{2}\qquad
Dario Trevisan\textsuperscript{2}\\[0.7em]
\small\textsuperscript{1}LAMA, CNRS, Universit\'e Gustave Eiffel, Universit\'e Paris-Est Cr\'eteil,\\
\small 5 boulevard Descartes, 77454 Marne-la-Vall\'ee, France\\[-0.1em]
\small\texttt{michael.goldman@cnrs.fr}\\[0.65em]
\small\textsuperscript{2}Universit\`a di Pisa, Dipartimento di Matematica,\\
\small Largo Bruno Pontecorvo 5, 56127 Pisa, Italia\\[-0.1em]
\small\texttt{francesco.grotto@unipi.it}\quad
\texttt{dario.trevisan@unipi.it}
}
\date{\today}
\begin{document}
\maketitle
\begin{abstract}
Let $X_1,\dots,X_N$ be independent uniform points in the unit square,
matched either to an evenly spread deterministic configuration or to a
second independent sample. For exponents $\alpha_1,\alpha_2>2$ with
$1/\alpha_1+1/\alpha_2=1/2$, we prove that with probability tending to
one a single permutation controls the $j$th coordinate displacements in
the exponential Orlicz norm of exponent $\alpha_j$, at the scale
$\sqrt{\log N/N}$, for both $j=1,2$. This settles Talagrand's Ultimate
Matching Conjecture for all pairs of finite critical exponents. The
proof constructs, on the flat torus, a current whose divergence is
Gaussian white noise projected on a dyadic grid and whose coordinates
have the critical anisotropic exponential Orlicz bounds, and transfers
it to matching through a primal optimal-transport reduction. At the
hard endpoints $(2,\infty)$ and $(\infty,2)$ the conjecture fails: a
Brownian local-time construction shows that the normalized optimal
matching cost diverges in probability. This file was mostly generated by LLMs as a result of  interactions with us
 to try and understand the 'monolithic' versions of the proofs which may be found in the
 \textbf{associated GitHub repository}~\cite{umc:repository}.
A digested version of the proofs will be made available soon. 
\end{abstract}

\setcounter{tocdepth}{1}
\tableofcontents

\section{Introduction}\label{umc:integrated-overview}

\subsection{The matching problem}
Let $X_1,\ldots,X_N$ be independent uniform points in the unit square
$D=[0,1]^2$. Talagrand's \emph{Ultimate Matching Conjecture}, posed as
Research Problem~17.1.2 (equation~(17.3)) of
\cite{umc:Talagrand2021}, asks for a permutation matching this random
sample to \emph{evenly spread} deterministic points $Y_1,\ldots,Y_N$
with simultaneous exponential tail bounds in both coordinates.
Here ``evenly spread'' has Talagrand's meaning: $D$ admits
a partition into $N$ rectangles of area $1/N$, with one point $Y_i$
in each rectangle $R_i$, such that
$R_i\subset B(Y_i,10/\sqrt N)$. We also consider the matching of
two independent uniform samples, which satisfies the same conclusion.
Our first result, which we call the Ultimate Matching Theorem, is the
following.

\begin{theorem}\label{umc:thm:main}
Let $\alpha_1,\alpha_2$ be fixed finite exponents with
\[
  \alpha_1,\alpha_2>2,
  \qquad \frac{1}{\alpha_1}+\frac{1}{\alpha_2}=\frac12.
\]
There exists a constant $L=L(\alpha_1,\alpha_2)>0$, independent
of $N$ and of the evenly spread configuration, such that the following
holds with probability tending to one as $N\to\infty$: for
\emph{every fixed evenly spread deterministic} $(Y_i)_{i=1}^N$, or
when $(Y_i)_{i=1}^N$ is a second independent uniform sample,
\emph{one and the same} permutation $\pi\in S_N$ satisfies, for both
coordinates $j=1,2$,
\begin{equation}\label{umc:eq:combined-main}
  \sum_{i=1}^N
  \exp\!\left[
    \left(
      \frac{\sqrt N\,\bigl|X_i^j-Y_{\pi(i)}^j\bigr|}
           {L\sqrt{\log N}}
    \right)^{\alpha_j}
  \right]
  \le 2N,\qquad j=1,2.
\end{equation}
\end{theorem}

The symmetric quartic case $(4,4)$ is included. The proof is completed
in Corollaries~\ref{cor:big-picture} and~\ref{red:cor:evenly-spread},
which combine the Gaussian flow estimate of
Theorem~\ref{thm:gauss} with the reduction of
Part~\ref{umc:partII}.

\paragraph{Finite subcritical exponents.}
The theorem also gives \eqref{umc:eq:combined-main} for every fixed
pair $0<\alpha_1,\alpha_2<\infty$ with
$\alpha_1^{-1}+\alpha_2^{-1}>1/2$, after changing the constant $L$.
Indeed, set
\[
 c=2(\alpha_1^{-1}+\alpha_2^{-1})>1,\qquad r_i=c\alpha_i.
\]
Then $r_i>\alpha_i$, $r_i>2$, and $r_1^{-1}+r_2^{-1}=1/2$,
so the critical theorem applies with exponents $r_1,r_2$.
For the resulting permutation, write
$t_{i,j}=\sqrt N\,|X_i^j-Y_{\pi(i)}^j|/(L\sqrt{\log N})$.
Since $t^{\alpha_j}\le1+t^{r_j}$ for $t\ge0$, Jensen's inequality
with a fixed $A\ge1$ gives
\[
 \frac1N\sum_{i=1}^N e^{(t_{i,j}/A)^{\alpha_j}}
 \le\left(\frac1N\sum_{i=1}^N e^{t_{i,j}^{\alpha_j}}\right)^{A^{-\alpha_j}}
 \le(2e)^{A^{-\alpha_j}}\le2
\]
when $A^{\alpha_j}\ge\log(2e)/\log2$ for both $j=1,2$.
Thus the same permutation works at the same scale, with $L$
replaced by $AL$; the probability of success is unchanged.

\paragraph{The hard endpoints.}
The coordinate endpoints $(2,\infty)$ and $(\infty,2)$ behave
differently. Theorem~\ref{neg:thm:main} in Part~\ref{umc:partIII}
shows that there the optimal matching cost, divided by
$\sqrt{\log N/N}$, diverges in probability, uniformly over deterministic
reference configurations. The Gaussian construction of
Part~\ref{umc:partI} applies only to finite exponents, and the
endpoint obstruction is of a different nature. These coordinatewise
endpoints should not be confused with the usual scalar $L^\infty$
matching problem or with an endpoint of a power-cost problem.
In the endpoint cases it is in fact possible to determine the exact first
order asymptotic (in probability) behavior of the cost: for its (undigested,
AI assisted) derivation we refer once again to the associated GitHub
repository~\cite{umc:repository}, since here we are mainly interested in the
failure of Talagrand's conjecture.

\subsection{Previous work}

\paragraph{The dual route.}
The foundations of planar random matching go back to Ajtai,
Koml\'os and Tusn\'ady \cite{umc:AKT84}. Talagrand's work
\cite{umc:Talagrand92,umc:Talagrand2005,umc:Talagrand2021}
identified the multiscale nature of the problem and led
to Research Problem~17.1.2 in the 2021 monograph.
The guiding method is \emph{dual}: one estimates transport or
discrepancy functionals against classes of test functions, and
controls the resulting random suprema through chaining, majorizing
measures and geometric decompositions. This approach has produced
sharp matching estimates in several regimes.

In dimension three, Chapter~18 of \cite{umc:Talagrand2021} proves a
sharp matching theorem between an independent uniform sample and an
evenly spread deterministic configuration, including
coordinate-dependent exponential costs and hard coordinate
constraints (Theorem~18.1.1 and Exercise~18.1.4); Talagrand observes
that the extension to dimensions $d>3$ requires no new ideas. Above
dimension two, irregularities at different scales cannot accumulate
in the same way, whereas in the plane their cumulative effect produces the critical
logarithmic factor. The remaining difficulty is therefore specifically
two-dimensional.

The dual direction continues to develop: Cheng, Otto, Palmieri and
Wachtel \cite{umc:COPW26} study min--max matching through a
geometrically linear approximation of a dual formulation, and Ouimet
\cite{umc:Ouimet26} proposes a solution to Problem~5.3 in Talagrand's
\emph{My favorite problems} \cite{TalagrandProblems}, controlling
simultaneously total distance and maximal displacement through
electrical flows and planar Lipschitz truncation of lattice
circulations.

\paragraph{The primal route.}
Our proof follows a complementary tradition. Rather than estimating the
dual supremum, we construct an actual \emph{primal} transport current
carrying the discrepancy of the random measures. This approach was
brought into focus by the statistical-physics prediction of
Caracciolo, Lucibello, Parisi and Sicuro \cite{umc:CLPS14}: at leading
order in quadratic bipartite matching, the displacement should behave
like the gradient of a Poisson potential driven by empirical-density
fluctuations, which is the linearization of the Monge--Amp\`ere
equation. Ambrosio, Stra and Trevisan \cite{umc:AST19} established the
leading two-dimensional quadratic asymptotics rigorously, and finer
results on the cost and the optimal map followed
\cite{umc:AG19,umc:AGTmap19}. In parallel, the variational regularity
theory of optimal transport, through the harmonic approximation of
Goldman and Otto \cite{umc:GO20} and its quantitative versions by
Goldman, Huesmann and Otto \cite{umc:GHO21,umc:GHO23}, quantified the
approximation of displacements by Poisson fluxes; related methods
treat nonuniform planar domains \cite{umc:AGT22}. These results make
it natural to study random matching through the Eulerian continuity
equation and a compatible vector current.

For nonquadratic costs the corresponding equation is nonlinear.
Ambrosio, Vitillaro and Trevisan \cite{umc:AVT24} obtained sharp
estimates for power-law costs, and Armegioiu, Goldman, Grotto and
Trevisan \cite{umc:AGGT26} related finite power-cost asymptotics on the
torus to a Gaussian variational-martingale problem and a parabolic
Monge--Amp\`ere Bellman equation. In this viewpoint successive
logarithmic scales are accumulated through stochastic control; it
directly motivates the construction of Part~\ref{umc:partI}. The
power-cost results do not, by themselves, give the coordinatewise
exponential Orlicz bounds proved here.

\subsection{Strategy of the proof}

\paragraph{Critical Gaussian currents.}
The main analytic task is to solve a divergence equation for
mean-zero Gaussian white noise projected onto dyadic squares of
mesh $h=2^{-m}$ on the flat two-dimensional torus. With
$\psi_r(t)=e^{|t|^r}-1$ and $\|\cdot\|_{\psi_r}$ the associated
Luxemburg norm, Part~\ref{umc:partI} constructs, with probability
bounded below uniformly in $m$, a periodic current $J$ with
\[
 \operatorname{div}J=W_h,
 \qquad \max_{i=1,2}\|J_i\|_{\psi_{r_i}}\le C\sqrt m,
\]
for every finite critical pair $(r_1,r_2)=(\alpha_1,\alpha_2)$
(Theorem~\ref{thm:gauss}). The difficulty is not the existence of an
inverse of the divergence, but the selection of one whose two
coordinate tails are simultaneously of the critical size.

At a single point, a Bellman equation suggests an inverse-Hessian
covariance for a Gaussian martingale, and a supersolution evaluated
backward in Bellman time formally gives a supermartingale. The
pointwise choice, however, need not be compatible across space.
Part~\ref{umc:partI} therefore reveals the Gaussian increments through
predictable local divergence corrections on anisotropic Whitney cells.
An adaptive metric follows the locally preferred covariance, and a
discounted memory records spatial oscillations of the accumulated
current. A simultaneous induction on the metric and the memory yields a
contracting Lyapunov quantity, whose negative drift pays for the
variance of the critical exponential tilt. A terminal Gaussian
correction then transfers the current from the adaptive partition to
the fixed dyadic grid.

\paragraph{From currents to matching.}
Part~\ref{umc:partII} supplies the reductions, independently of the
adaptive construction. The minimum-flow gauge is uniformly Lipschitz in
the dyadic Gaussian coefficients, so a fixed positive success
probability upgrades to high probability by Gaussian concentration. An
exact four-sheet folding transfers the torus estimate to the square.
Dyadic regularization, a comparison between transport and static
currents, and an interval-encoded KMT coupling relate the Gaussian
source to one empirical histogram against Lebesgue measure, and then to
two independent histograms. Talagrand's evenly spread points are a
further $O(N^{-1/2})$ transport from Lebesgue measure, and a Birkhoff
decomposition yields one permutation with both coordinate bounds. These
steps give \eqref{umc:eq:combined-main} in both formulations.

\paragraph{The endpoint obstruction.}
At the endpoint $(2,\infty)$ the reduction of Part~\ref{umc:partII}
still applies, so it suffices to show that the square Gaussian
minimum-flow gauge is not of order $\sqrt m$. Part~\ref{umc:partIII}
builds random dual test potentials of bounded cost whose expected
pairing with the projected noise is too large. The input is a planar
Brownian martingale with bounded terminal anisotropic Orlicz cost but
arbitrarily large expected integrated square-root determinant,
obtained from short Brownian local-time stages. An adaptive tree of
piecewise affine Gaussian bubbles converts this determinant into a
positive Gaussian pairing.

\subsection{Organization}
Part~\ref{umc:partI} (Sections~\ref{sec:introduction}--\ref{sec:terminal})
proves the Gaussian flow estimate on the torus; the geometric estimates
for its adaptive construction are collected in
Appendix~\ref{app:geometry}. Part~\ref{umc:partII}
(Sections~\ref{sec:interface}--\ref{red:sec:consequences}) reduces
matching to the square Gaussian flow problem, upgrades fixed-probability
bounds to high probability, transfers the torus estimate to the square,
and completes the proof of Theorem~\ref{umc:thm:main}.
Part~\ref{umc:partIII}
(Sections~\ref{neg:sec:introduction}--\ref{neg:sec:conclusion}) proves
the failure at the hard endpoints.

\clearpage
\part{Critical Gaussian Orlicz flows on the torus}\label{umc:partI}

In this part, $W_h$ is the projection of mean-zero Gaussian white
noise on the flat torus $\T^2$ onto the dyadic square grid of mesh
$h=2^{-m}$. For fixed $2<r_1,r_2<\infty$ with $1/r_1+1/r_2=1/2$, we
bound the minimum coordinatewise Orlicz norm of a periodic current
with divergence $W_h$ by a multiple of $\sqrt m$, with probability
bounded below independently of $m$ (\cref{thm:gauss}). The construction combines predictable Gaussian increments,
local inverses of the divergence in anisotropic Whitney charts, and a
contracting metric--memory functional. A terminal Gaussian correction
transfers the adaptive current to the fixed dyadic source.

\section{The Gaussian flow problem and the strategy}\label{sec:introduction}

\subsection{The Gaussian minimum-flow problem}

Let $W$ be mean-zero real white noise on the flat torus $\T^2$:
\begin{equation}\label{eq:white}
 \E W(f)W(g)=\int_{\T^2}fg-\int_{\T^2}f\int_{\T^2}g.
\end{equation}
\phantomsection\label{not:fixedgrid}
For $h=2^{-m}$, let $P_h$ denote averaging on the fixed periodic
$h$-square grid and write $W_h=P_hW$; thus
$W_h|_Q=h^{-2}W(Q)$ and $\int_{\T^2}W_h=0$.
\phantomsection\label{not:orlicz}
We use the Luxemburg norm $\|\cdot\|_{\psi_r}$ associated with
$\psi_r(t)=e^{|t|^r}-1$. Fix finite exponents $r_i>2$ on the
critical line
\begin{equation}\label{eq:crit}
 a_i=\frac1{r_i},\qquad a_1+a_2=\frac12.
\end{equation}
\phantomsection\label{not:astar}
We also set $a_*:=\max_i a_i<1/2$. For every mean-zero source
$f$, define the minimum-flow gauge
\begin{equation}\label{eq:gauge}
 \mathfrak F_{\boldsymbol r}(f)
 =\inf_{\divv J=f}\max_{i=1,2}\|J_i\|_{\psi_{r_i}},
\end{equation}
where $\boldsymbol r=(r_1,r_2)$ and the divergence constraint is
understood distributionally for periodic currents. The problem is
to control this gauge for $W_h$.

\begin{theorem}\label{thm:gauss}
For every fixed finite $r_1,r_2>2$ with
$1/r_1+1/r_2=1/2$, there exist constants $C<\infty$, $\varepsilon_0>0$ and
$m_0<\infty$, depending only on the exponents, such that for
$m\ge m_0$,
\begin{equation}\label{eq:target}
 \Pp\left(\mathfrak F_{\boldsymbol r}(W_h)\le C\sqrt m\right)
 >\varepsilon_0.
\end{equation}
\end{theorem}

\begin{remark}\label{rem:concentration-upgrade}
The fixed-positive-probability estimate is the conclusion proved
in this part. As shown in Part~\ref{umc:partII}
(Subsection~\ref{red:subsec:gaussian-concentration}),
the fixed-grid minimum-flow gauge is Lipschitz in the Gaussian
coefficients with a constant independent of $h$. Gaussian
concentration therefore upgrades \eqref{eq:target} to
$\Pp(\mathfrak F_{\boldsymbol r}(W_h)>C'\sqrt m)
\le 2e^{-cm}$ for suitable $C',c>0$. This improvement requires
no further adaptive-flow estimates and is not used below.
\end{remark}

\subsection{The quartic case}

Consider first $r_1=r_2=4$. The proof controls the exponential
cost of the accumulated current through the phase
\[
 H(y)=(1+y_1^2)^2+(1+y_2^2)^2.
\]
Here $y\in\R^2$ is a current value: at a spatial point $x$ and
reveal generation $k$, it will be $y=U_k(x)$. It is not a spatial
coordinate. The quadratic terms in $H$ give nondegenerate curvature
near the origin, while its growth at infinity is quartic.

The choice of the next covariance is suggested by the curvature
of $e^H$. Since $D^2H\preceq C H^{1/2}I$, the matrix
\[
 B(y)=H(y)^{1/2}I+\nabla H(y)\nabla H(y)^T
\]
controls both the Hessian term and the square of the linear
variation of the phase. Among positive definite matrices of
determinant one, the choice
$C(y)=\sqrt{\det B(y)}B(y)^{-1}$ minimizes $\tr(B(y)G)$;
the calculation is given in Section~\ref{sec:selector}.
For example, at $y=(R,0)$ with $R\to\infty$,
\[
 B(y)\asymp\operatorname{diag}(R^6,R^2),\qquad
 C(y)\asymp\operatorname{diag}(R^{-2},R^2).
\]
The corresponding phase cost is $\tr(BC)\asymp R^4\asymp H$,
whereas an isotropic choice gives $\tr B\asymp R^6$.
Thus the preferred covariance reduces the variance in the direction
where the current is already large.

This pointwise choice cannot be imposed independently at every
spatial point. The local divergence equation must be solved on a
whole cell, in a chart chosen before the next noise is revealed.
The construction therefore uses one covariance-control matrix $G$
per cell and moves it towards $C$ by bounded updates. A memory of
old spatial oscillations controls the error between the state at
the cell representative and the state at other points. The next
subsection explains this simultaneous metric--memory induction.

For general fixed finite critical exponents, the same construction
uses
\begin{equation}\label{eq:phaseH}
 F_i(t)=(1+t^2)^{r_i/2},\qquad
 H(y)=F_1(y_1)+F_2(y_2),\qquad v(y)=\nabla H(y).
\end{equation}

\subsection{Structure of the proof}

The noise is revealed through a finite predictable sequence of
Gaussian increments. At each stage a local divergence correction adds
the newly revealed noise to the current, with a covariance chosen from
the old history on each active cell. The aim is to control a spatial
exponential modular of the accumulated current. The phase is $H$ from
\eqref{eq:phaseH}, and the preferred covariance is selected by the
regularized Hessian proxy in \eqref{eq:B}. No continuous-time Bellman
equation is solved in the construction; its formal relation to the
discrete argument is explained in
Subsection~\ref{subsec:bellman-motivation}.

The inverse-Hessian rule does not by itself solve the spatial
problem. The preferred covariance $C(U_k(x))$ depends on $x$, whereas
an inverse of the divergence must be constructed coherently on a
whole region. Moreover, the covariance must be chosen \emph{before}
the next Gaussian variables are revealed. The construction therefore
follows the pointwise Bellman covariance with a predictable metric and
records the price of not following it exactly. This is done by a
simultaneous induction on the routing metric and on a memory of
spatial oscillations.

Section~\ref{sec:innovation} reveals the source through a finite
predictable tree. Section~\ref{sec:geom} provides a local divergence
inverse in intrinsic Whitney coordinates and a normalized Gaussian
increment with uniform conditional variance and H\"older envelopes
$E_{k+1}$. These bounds do not depend on the eccentricity of the
physical chart. We use the scales
\begin{equation}\label{eq:parameters}
 s=hm^3,\qquad \delta=\frac{\varepsilon^2}{m},\qquad
 \lambda_k=e^{-\Lambda k\delta},
\end{equation}
where $\varepsilon>0$ is a small fixed step parameter and $\Lambda>0$
is a sufficiently large fixed damping constant. The cumulative
normalized current satisfies $\divv U_k=\sqrt\delta\,P_{\pr_k}W$;
its unscaled counterpart is $U_k/\sqrt\delta$. The grid parameter
$m$ is fixed throughout the construction; $k$ counts reveals
and does not specify a common spatial size of the cells.
\phantomsection\label{not:depthbound}
The depth of any actual revealed chain is at most
$J_0\le C_{\rm dep}m$, for a fixed geometric constant
$C_{\rm dep}$.

\emph{Metric lag.}
On a parent $Q$, choose an old-measurable representative $x_Q$.
Section~\ref{sec:selector} selects an updated determinant-one
metric $G_Q^+$ by clipping the old metric $G_k$ towards
$C(U_k(x_Q))$. The clipping stays within a fixed factor of $G_k$
and contracts, at $x_Q$, the lag
\[
 L_k(x)^2=\operatorname{tr}\!\left(G_k(x)
                      C(U_k(x))^{-1}\right).
\]
If $U_k$ were constant on $Q$, this would be a contraction
throughout $Q$. But $U_k$ varies in space while the metric is chosen
once on $Q$, and the difference $U_k(x)-U_k(x_Q)$ is exactly what
prevents applying the pointwise Bellman optimization independently
at every location.

\emph{Memory.}
Each old increment is H\"older in the intrinsic metric in which it
was created, and descendants shrink in these birth metrics, even
after the bounded changes of metric between refinements. Hence its
oscillation across a much later descendant is reduced by a factor
$r_{\rm osc}^{k-j}$. The old diagonal metric caps and relative
comparisons of $H(U_j)$ then give, schematically for $x,z$ in the
current parent,
\[
 |U_{k,i}(x)-U_{k,i}(z)|
 \lesssim\sqrt{\delta M}\,H(U_k(x))^{a_i}
       \sum_{j\le k}r_{\rm osc}^{k-j}E_j(x).
\]
Along the actual reveals of each branch we retain the discounted
square function
\[
 T_k=\sum_{j\le k}\rho_1^{k-j}E_j^2,\qquad
 S_k=L_k^2+T_k,\qquad r_{\rm osc}^2<\rho_1<1.
\]
Weighted Cauchy--Schwarz bounds the ancestor sum by $C\sqrt{T_k}$.
Thus $T_k$ pays for the oscillation of the old field between the
representative, where the metric is selected, and the point where
the estimate is used.

\emph{Simultaneous induction.}
Section~\ref{sec:bootstrap} closes three estimates together: the
coordinatewise caps $G_{k,ii}\le M H(U_k)^{2a_i}$, with a fixed
large constant $M$; small relative changes of the phase $H(U_k)$;
and the contraction of $S_k$. The \emph{old} memory estimate controls
$U_k(x)-U_k(x_Q)$, so the representative-based metric update
preserves the diagonal caps with a strict margin. The \emph{fresh}
Gaussian envelope gives
$|V_{k+1,i}|\lesssim\sqrt\delta\,H(U_k)^{a_i}E_{k+1}$, which
preserves the caps and the phase comparison at the new state. Only
then do we use the old stopped exponential modular, together with
the lower area bound for active cells, to obtain $H(U_k)\lesssim m$
pointwise. Consequently $\sqrt{\delta H(U_k)}\lesssim\varepsilon$.

This small factor makes the metric update stable under spatial
variation of the state. After the representative contraction and a
comparison of the lag along the old-to-new state segment, which uses
only old-state information, one obtains
\[
 L_{k+1}^2\le q_L L_k^2+C+C\varepsilon^2T_k
                       +CE_{k+1}^2,\qquad q_L<1.
\]
On an actual reveal, $T_{k+1}=\rho_1T_k+E_{k+1}^2$. Choosing
$\rho_1$ and $\varepsilon$ with strict margins, the two inequalities
combine into the renewal estimate
\begin{equation}\label{eq:intro-renewal}
 \boxed{\quad S_{k+1}\le q_0S_k+C(1+E_{k+1}^2),
                   \qquad q_0<1.\quad}
\end{equation}
The old memory justifies the new metric, and the contraction of the
new metric pays for retaining the memory. This is why the two are
propagated simultaneously.

\emph{Exponential estimate.}
The renewal bound \eqref{eq:intro-renewal} supplies the negative
drift that turns the exponential Bellman trial into a spatial
supermartingale. Write $H_k=H(U_k)$. The phase used in
Section~\ref{sec:phase} is
\phantomsection\label{not:lyapunovfirst}
\[
 \Phi_k=\lambda_k H_k(1+\kappa\delta S_k),
\]
where $\kappa>0$ is fixed and the clock $\lambda_k$ decreases
slightly with each reveal. Under a fresh reveal, the linear variation
of $H_k$ is Gaussian. At criticality its conditional variance is
bounded, up to constants, by $\lambda_k\delta H_kS_k$: the factor
$S_k$ measures how far the actual covariance is from the pointwise
Bellman covariance. Exponentiating this linear term produces a
positive drift of that size, and the term $\kappa\delta H_kS_k$ is
added to the exponent to compensate it. On a good reveal,
$H_{k+1}\le(1+\xi)H_k$ for a small fixed $\xi$, and
\eqref{eq:intro-renewal} gives
\[
 H_{k+1}S_{k+1}-H_kS_k
 \le -d_0 H_kS_k+C H_k(1+E_{k+1}^2),
 \qquad d_0=1-q_0(1+\xi)>0.
\]
Thus the additional phase contributes
$-\kappa d_0\lambda_{k+1}\delta H_kS_k$, which dominates the
conditional Gaussian variance cost once $\kappa$ is large. The
remaining errors have size $\lambda_{k+1}\delta H_k(1+E_{k+1}^2)$.
The square of the envelope has a uniform conditional exponential
moment, and its coefficient satisfies $\delta H_k\lesssim\varepsilon^2$
by the old modular bound $H_k\lesssim m$. The constant part is paid
for by the decrease $\lambda_k-\lambda_{k+1}\asymp\lambda_{k+1}\delta$.

The order of these estimates matters. We bound the product
$H_{k+1}S_{k+1}$ by relative phase control, rather than expanding the
new lag inside the phase. Consequently the only linear Gaussian term
is the variation of $H_k$, whose coefficient depends on the old state
and on the metric selected \emph{before} the noise is revealed. Its
conditional moment-generating function can therefore be evaluated
directly, and no independence between the Gaussian increment and its
H\"older envelope is needed. This gives a one-step conditional
estimate for the spatial exponential modular.

To iterate it, Section~\ref{sec:phase} uses a nonnegative
killed-and-frozen numerical supermartingale. On histories that pass
every envelope and modular test it is the spatial integral of
$e^{\Phi_k}$; on an envelope failure it is set to zero. Overshoot stopping
and the separately controlled probability of an envelope failure
then give a bounded spatial modular with a fixed positive
probability. Terminal branches are frozen without artificial memory
discounts. In this sense the construction plays the role of the
formal Bellman supermartingale $u(T-t,X_t)$ from
\eqref{eq:intro-bellman-supermart}, with the metric--memory mechanism
paying for spatial incompatibility.

\emph{Terminal correction.}
Section~\ref{sec:terminal} transfers the divergence from the adaptive
partition to the fixed-grid source $W_h$ in physical coordinates.
Combining that correction with the successful leaf current gives
\eqref{eq:target}. The geometric estimates used in the construction
are collected in Appendix~\ref{app:geometry}.

\subsection{Bellman motivation}\label{subsec:bellman-motivation}

The starting analogy is the power-energy problem
\[
 \inf_{\divv J=W_h}\int_{\T^2}|J(x)|^p\,\dd x,
 \qquad 1<p<\infty.
\]
For fixed $p$, its logarithmic-scale asymptotics lead to a
stochastic-control problem; see \cite{umc:AGGT26}. Fix a
\emph{spatial point} $x\in\T^2$ and reveal the Gaussian noise
through successive logarithmic Fourier shells. The divergence
constraint fixes the longitudinal contribution to the current at
$x$, but transverse corrections can be optimized. In the
continuous-scale limit, the accumulated two-component current at
this fixed point is modeled by a controlled Gaussian martingale.
Its state belongs to $\R^2$; we denote a generic state by $y$.
The corresponding power-cost Bellman equation is
\begin{equation}\label{eq:intro-ma}
 \partial_tu_p(t,y)=\frac12\sqrt{\det D_y^2u_p(t,y)},
 \qquad u_p(0,y)=|y|^p.
\end{equation}
Where its Hessian is positive definite, the optimal
one-point covariance has the determinant-one inverse-Hessian form
\[
 C_p(t,y)=\sqrt{\det D_y^2u_p(t,y)}\,
           (D_y^2u_p(t,y))^{-1}.
\]
This is the continuous-time antecedent of the selector $C(y)$
used later in \eqref{eq:B}.

Here is the supermartingale interpretation, which will guide the
critical construction. For any smooth, strictly convex candidate
$u$ with positive Hessian, denote by $C_u$ the same normalized
inverse-Hessian matrix. Fix a terminal Bellman time $T>0$ and an
initial current state $X_0=y$. Formally let the one-point
current evolve as
\[
 \dd X_t=\frac1{\sqrt2}C_u(T-t,X_t)^{1/2}
                  \,\dd\mathbf B_t,\qquad 0\le t\le T,
\]
where $\mathbf B_t$ is a standard planar Brownian motion modeling
the newly revealed Gaussian innovations after the transverse
control. Thus $X_t$ is \emph{the current accumulated at the fixed
spatial point $x$}, not the position of a particle or a spatial
current on the whole torus. Its conditional covariance per unit
time is $\frac12C_u(T-t,X_t)$. The factor $1/2$ is exactly the
time normalization in \eqref{eq:intro-ma}. Since
$\frac14\operatorname{tr}(C_uD_y^2u)
 =\frac12\sqrt{\det D_y^2u}$, It\^o's formula gives
\begin{equation}\label{eq:intro-bellman-supermart}
 \dd\,u(T-t,X_t)
 =\left[\frac14\operatorname{tr}(C_uD_y^2u)
              -\partial_tu\right](T-t,X_t)\dd t
   +\dd M_t\ \le\ \dd M_t
\end{equation}
whenever $u$ is a Bellman supersolution. With equality in the
Bellman equation the drift vanishes. In the power case,
$u_p(0,X_T)=|X_T|^p$, so the supermartingale property yields the
one-point terminal-cost bound $\E|X_T|^p\le u_p(T,y)$ (after
localization and the necessary integrability justification).
The reversal of Bellman time is essential: the PDE evolves in the remaining-time variable,
while $X_t$ accumulates newly revealed noise. This calculation is
formal (it requires localization and a well-defined feedback
martingale), but it identifies the goal: a supermartingale built
from an exponential trial evaluated along the evolving current.

For the Orlicz problem, use the phase $H$ from \eqref{eq:phaseH}
and the critical weights $a_i=1/r_i$.
\phantomsection\label{not:bellmanclock}
Try $u(t,y)=\exp\{\zeta_{\rm B}(t)H(y)\}$, with
$\zeta_{\rm B}>0$ an auxiliary increasing Bellman-time coefficient.
\phantomsection\label{not:bellmansigma}
To compare with the subcritical exponents, put
$\sigma=a_1+a_2-\frac12\ge0$, so $\sigma=0$ on the critical
line. Differentiation and the $2\times2$ rank-one determinant
identity yield
\begin{equation}\label{eq:intro-bellman-determinant}
 e^{-\zeta_{\rm B}H}\sqrt{\det D_y^2u}
 \le C(\zeta_{\rm B}+\zeta_{\rm B}^{3/2})H^{1-\sigma}.
\end{equation}
As $\partial_tu=\zeta_{\rm B}'Hu$, choosing
$\zeta_{\rm B}'\ge C(\zeta_{\rm B}+\zeta_{\rm B}^{3/2})$
produces a supersolution for a prescribed finite interval,
provided the initial coefficient is sufficiently small.
In particular, the exponential Bellman supersolution persists
at criticality, including the symmetric quartic example
$r_1=r_2=4$. The subcritical advantage is instead the
large-amplitude power reserve $H^{-\sigma}$, which disappears
when $\sigma=0$.

The one-point limiting cost is heuristically finite, but this
discards all the geometry. A pointwise covariance optimization
does not construct a \emph{spatially compatible} current
$x\mapsto J(x)$ driven by the same white noise and satisfying
$\divv J=W_h$. At polynomial order the reconstruction errors
can be controlled by Wiener-chaos estimates as in
\cite{umc:AGGT26}. At critical exponential growth the missing
power reserve means that the one-point supermartingale cannot
by itself absorb the errors of spatial reconstruction.
The solution sought in this part is therefore a \emph{spatial}
analogue of \eqref{eq:intro-bellman-supermart}, coupled to an
adaptive routing metric and a memory of spatial oscillations.

\subsection{Notation}\label{subsec:notation}
The notation used throughout Part~\ref{umc:partI} is collected in Table~\ref{tab:notation}.
All references point to the first definition of the corresponding object.

{\small
\setlength{\tabcolsep}{4pt}
\renewcommand{\arraystretch}{1.12}
\begin{longtable}{@{}p{0.285\linewidth}p{0.535\linewidth}p{0.14\linewidth}@{}}
\caption{Principal notation and location of first definitions.}\label{tab:notation}\\
\toprule
Symbol & Meaning & Definition \\
\midrule
\endfirsthead
\toprule
Symbol & Meaning & Definition \\
\midrule
\endhead
\bottomrule
\endfoot
$W$ & Mean-zero white noise on $\T^2$. & \eqref{eq:white}\\
$P_h,\,W_h$ & Averaging on the fixed dyadic $h$-grid; $W_h=P_hW$. & p.~\pageref{not:fixedgrid}\\
$\boldsymbol r=(r_1,r_2),\,a_i$ & Orlicz exponents and critical weights $a_i=1/r_i$. & \eqref{eq:crit}\\
$a_*$ & Larger critical weight $\max_i a_i<1/2$. & p.~\pageref{not:astar}\\
$\psi_r,\,\|\cdot\|_{\psi_r}$ & Exponential Young function and Luxemburg norm. & p.~\pageref{not:orlicz}\\
$\mathfrak F_{\boldsymbol r}$ & Minimum coordinatewise Orlicz flow gauge. & \eqref{eq:gauge}\\
$h,\,s,\,\delta,\,\varepsilon$ & Fixed-grid mesh, terminal intrinsic scale, time step $\delta=\varepsilon^2/m$, and fixed small $\varepsilon$. & \eqref{eq:parameters}\\
$J_0,\,C_{\rm dep}$ & Deterministic bound $J_0\le C_{\rm dep}m$ for actual reveal depth. & p.~\pageref{not:depthbound}\\
$\lambda_k,\,\Lambda$ & Deterministic damping clock and its fixed rate. & \eqref{eq:parameters}\\
$\pr_k$ & Nested partition after the $k$th reveal. & p.~\pageref{not:predictabletree}\\
$\F_k$ & Sigma-field of geometry and masses revealed through stage $k$. & p.~\pageref{not:revealedfiltration}\\
$f_k,\,\Delta f_k$ & Adaptive projection of $W$ and its next Gaussian increment. & p.~\pageref{not:adaptiveprojection}\\
$Q_k(x),\,M_k(x),\,v_{k+1}(x)$ & Cell containing $x$, its Gaussian average, and next conditional variance. & \eqref{eq:cellmartnotation}\\
$G_Q^+,\,A_Q,\,\widehat Q$ & Predictable routing metric, its square root, and the intrinsic parent chart. & \eqref{eq:intrinsicchart}\\
$K,\,c_0,\,D$ & Fixed adjacent-metric distortion, inball ratio, and current intrinsic parent diameter. & p.~\pageref{not:Kconstant}\\
$\theta,\,D_0$ & Whitney size cap and intrinsic terminal stopping threshold $D\le D_0s$. & p.~\pageref{not:whitneyparameters}\\
$\gamma,\,A,\,C_A$ & Whitney boundary-clearance, inball-displacement and receiver-size constants. & Lemma~\ref{lem:whitney}\\
$\vartheta$ & Child size contraction across an actual birth-chart refinement, distinct from $\theta$. & p.~\pageref{not:whitneyparameters}\\
$z,\,\widehat S$ & Intrinsic coordinates and transformed child. & \eqref{eq:intrinsicchart}\\
$d_{G_Q^+}$ & Distance in the predictable intrinsic metric. & p.~\pageref{not:intrinsicdistance}\\
$\ell_S$ & Intrinsic scale of a Whitney child. & \eqref{eq:whitney}\\
$\alpha,\,\beta$ & H\"older exponents: kernel $\alpha$ and field $\beta$. & p.~\pageref{not:holderparameters}\\
$c_*$ & Fixed positive floor $|S|\ge c_*s^2$ for revealed-cell areas. & Lemma~\ref{lem:randommax}\\
$\widehat K_Q,\,\widehat{\mathcal B}_Q$ & Ray kernel and local divergence inverse in intrinsic coordinates. & \eqref{eq:ray}\\
$K_Q^{\rm phys},\,\mathcal B_Q^{\rm phys}$ & Corresponding physical Piola kernel and inverse. & \eqref{eq:piolakernel}\\
$\widehat P$ & Averaging onto the intrinsic Whitney child partition. & p.~\pageref{not:childprojector}\\
$H(y),\,v(y)$ & Anisotropic exponential phase and its gradient. & \eqref{eq:phaseH}\\
$\zeta_{\rm B}(t),\,\sigma$ & Bellman time coefficient (distinct from $\lambda_k$) and subcritical reserve $\sigma=a_1+a_2-1/2$. & pp.~\pageref{not:bellmanclock}--\pageref{not:bellmansigma}\\
$B(y),\,C(y)$ & Bellman matrix and determinant-one target metric. & \eqref{eq:B}\\
$\mathcal T_C$ & Spectral metric clipping toward $C$, with relative distortion bounded by $K$. & Lemma~\ref{lem:selector}\\
$L(y,G)$ & Lag between actual and target metrics. & \eqref{eq:lag}\\
$G_k,\,G_k^+$ & Birth metric and predictable updated metric at a split. & p.~\pageref{not:birthmetric}\\
$M$ & Fixed coefficient in the diagonal cap $G_{ii}\le M H(U)^{2a_i}$. & Lemma~\ref{lem:selector}\\
$\eta$ & Relative change allowed for $H(U_k)$ at an actual reveal. & p.~\pageref{not:relativephase}\\
$r_{\rm osc}$ & Effective ancestor spatial-oscillation decay factor. & Lemma~\ref{lem:actualphase}\\
$q_d,\,q_L,\,q_0$ & Diagonal-cap, lag and combined-memory contraction factors. & pp.~\pageref{not:qd},\pageref{not:qL}\\
$\xi$ & Realized relative phase-change bound $|H'-H|\le\xi H$, controlled by $\varepsilon$. & p.~\pageref{not:xi}\\
$Z_{k+1}$ & Intrinsic normalized Gaussian increment field. & p.~\pageref{not:intrinsicgauss}\\
$V_{k+1},\,U_k$ & Physical increment and cumulative divergence current. & \eqref{eq:update}\\
$E_{k+1,S}$ & Gaussian supremum--Hölder envelope on a child. & \eqref{eq:envelope}\\
$T_k,\,S_k$ & Discounted spatial memory and combined lag--memory energy. & \eqref{eq:memory}\\
$\rho_1$ & Fixed memory discount; satisfies $r_{\rm osc}^2<\rho_1<q_0$. & \eqref{eq:memory}, \eqref{eq:strictschedule}\\
$R$ & Global good-event cutoff $E_{k,S}\le R\sqrt m$ for Gaussian envelopes. & \eqref{eq:globalcut}\\
$T_*$ & Fixed stopping threshold for the exponential modular. & p.~\pageref{not:modularfirst}\\
$\Phi_k,\,\kappa$ & Exponential Lyapunov phase and its positive memory penalty coefficient. & p.~\pageref{not:lyapunovfirst}\\
$\mathcal Z_k$ & Spatial exponential modular. & p.~\pageref{not:modularfirst}\\
$\pr_*$ & Final predictable stopping partition. & \eqref{eq:bridgevar}\\
$g^{(1)},\,g^{(2)}$ & Revealed and unrevealed pieces of the terminal divergence correction. & \eqref{eq:decomp}\\
$J^{(1)},\,J^{(2)}$ & Corresponding physical correction currents. & pp.~\pageref{not:terminalcurrents}--\pageref{not:periodicgreen}\\
$D_*$ & Uniform bound on physical terminal-leaf diameters. & \eqref{eq:Dstar}\\
$K_{\T},\,d_{\T}$ & Periodic Green vector kernel and fixed physical torus distance. & \eqref{eq:periodicK}\\
$K_{x,h},\,\sigma_*^2(x)$ & Fixed-grid averaged kernel and conditional variance of the terminal bridge current. & \eqref{eq:terminalvar}\\
\end{longtable}
}

\section{Gaussian increments along predictable partitions}\label{sec:innovation}

This section describes how the noise is revealed. The partition is
refined predictably, and Lemma~\ref{lem:bridge} gives the conditional
Gaussian law of the new masses at each split. Lemma~\ref{lem:randommax}
then bounds all revealed cell averages simultaneously, although the
refinement tree depends on $W$.

\phantomsection\label{not:predictabletree}
Start with the full fundamental square, understood modulo opposite sides, as a
single root $Q_0$; the mean-zero convention gives $W(Q_0)=0$. The identified
sides form a coordinate seam, not a physical boundary. A local ray current on $Q_0$ has zero
trace across that seam and thus defines a periodic current. At generation $k$ a
finite partition $\pr_k$ and the masses $\{W(Q):Q\in\pr_k\}$ have been revealed.
The next subdivision of each active $Q$ is selected using only this history,
then its child masses are revealed. On inactive cells nothing is revealed.

\phantomsection\label{not:adaptiveprojection}
Set
\[
 f_k=P_{\pr_k}W=\sum_{Q\in\pr_k}\frac{W(Q)}{|Q|}\one_Q,
 \qquad \Delta f_k=f_{k+1}-f_k.
\]
For a child $S\subset Q$, define the new mass
\[
 \zeta_S=W(S)-\frac{|S|}{|Q|}W(Q).
\]

\begin{lemma}\label{lem:bridge}
Conditionally on the old history and chosen subdivision,
\begin{equation}\label{eq:covsplit}
 \E_k\zeta_S=0,\qquad
 \E_k(\zeta_S\zeta_T)=|S\cap T|-\frac{|S||T|}{|Q|},
\end{equation}
with conditional independence between distinct active parents. If $\pr_*$ is a finite predictable stopping partition and
$g\in L^2(\T^2)$ is deterministic, then
\begin{equation}\label{eq:bridgevar}
 \operatorname{Var}\left((W-P_{\pr_*}W)(g)\mid\F_*\right)
 =\sum_{Q\in\pr_*}\int_Q|g-g_Q|^2,
 \qquad g_Q=\fint_Qg.
\end{equation}
The conditional mean of this residual is zero.
\end{lemma}
\begin{proof}
Inside a parent $Q$, subtract from $\one_S$ its projection onto the constants on
$Q$. The resulting test function is
$\one_S-(|S|/|Q|)\one_Q$, which is orthogonal to the previously revealed
parent mass. Isonormality gives \eqref{eq:covsplit}. The new test functions are
chosen predictably, hence their coordinates are Gaussian conditional on the old
sigma-field. Apply this decomposition successively. At a stopping time pad the
tree to a deterministic depth, freezing all inactive leaves; conditional Gaussian
orthogonal projection on the resulting, random but predictably generated
finite-dimensional subspace proves \eqref{eq:bridgevar}. Here the test function
$g$ is deterministic; for random tests measurable with respect to the final
reveal history, the pairing with the isonormal residual would first have to be
defined conditionally. The identity would fail if a future subdivision were
selected \emph{after} inspecting unrevealed bridge coefficients.
\end{proof}

If the new density is sent through any predictable linear kernel $K_x$,
then $Z(x)=\Delta f_k(K_x)$ satisfies
\begin{equation}\label{eq:mgflinear}
 \E_k e^{u\cdot Z(x)}
 =\exp\bigl(\tfrac12u^T\operatorname{Cov}_k(Z(x))u\bigr),
 \qquad u\in\R^2.
\end{equation}
There is no Bernstein restriction on the tilt $u$. The adaptivity is handled
through predictability, not by declaring the random tree independent of $W$.


\subsection{Maxima of adaptive Gaussian cell averages}
\phantomsection\label{not:revealedfiltration}
Let $\pr_k$ be the partition after the $k$th reveal, with
$\pr_0=\{\T^2\}$, and let $\F_k$ contain the partitions and masses
revealed through that time. The refinement $\pr_{k+1}$ is selected
$\F_k$-measurably, before its new cell masses are observed.
For a deterministic point $x$ outside the cell boundaries, let
$Q_k(x)$ be the cell of $\pr_k$ containing $x$, and put
\begin{equation}\label{eq:cellmartnotation}
 M_k(x)=\frac{W(Q_k(x))}{|Q_k(x)|},\qquad
 v_{k+1}(x)=\frac1{|Q_{k+1}(x)|}-\frac1{|Q_k(x)|}.
\end{equation}
An inactive cell is retained in subsequent partitions, so that
$v_{k+1}(x)=0$ on an inactive branch. The following estimate applies
uniformly to this random refinement tree.

\begin{lemma}\label{lem:randommax}
Let $\pr_0,\ldots,\pr_J$ be a finite predictably refined partition sequence
of $\T^2$, with inactive cells frozen and every nonempty cell of area at
least $c_*s^2$. For $Q\in\pr_k$ put $M(Q)=W(Q)/|Q|$. There are
constants $C,c>0$ depending only on $c_*$ such that for every $b>0$,
\begin{equation}\label{eq:randommax}
 \Pp\!\left(\max_{0\le k\le J}\max_{Q\in\pr_k}
       |M(Q)|>\frac b s\right)
 \le Cs^{-2}\exp(-cb^2).
\end{equation}
In particular, if $s=2^{-m}m^3$, then
\begin{equation}\label{eq:randommaxwhp}
 \max_{k,Q}\frac{|W(Q)|}{|Q|}\le C\frac{\sqrt m}{s}
 \quad\text{with probability }1-o(1).
\end{equation}
\end{lemma}
\begin{proof}
Fix a deterministic $x$ for which \eqref{eq:cellmartnotation}
is well defined at every stage. Conditional on $\F_k$, the chosen
subdivision of $Q_k(x)$ is fixed, and the centered child bridge is
Gaussian. Hence $M_k(x)$ is a martingale with conditional increment
variance $v_{k+1}(x)\ge0$. The predictable
variance telescopes \emph{pathwise}:
\[
 \sum_{j=0}^{k-1}v_{j+1}(x)
 =|Q_k(x)|^{-1}-|Q_0(x)|^{-1}\le(c_*s^2)^{-1}.
\]
For each fixed $u\in\R$, the exponential
\[
 \exp\Big\{u(M_k(x)-M_0(x))-
              \tfrac12u^2\sum_{j<k}v_{j+1}(x)\Big\}
\]
is a nonnegative martingale. Apply the exponential maximal
inequality, first with $u>0$ and then with $-u$, and optimize
at $u\asymp bs$ to obtain
\begin{equation}\label{eq:pointxmax}
 \Pp\bigl(\max_{k\le J}|M_k(x)|>b/s\bigr)
 \le 2\exp(-c b^2).
\end{equation}
At the root $M_0(x)=W(\T^2)=0$. If \emph{any}
revealed cell at any time has a large average, then the spatial
set of $x$ for which the running maximum is large contains the
entire cell and hence has measure at least $c_*s^2$.
Consequently, for the event $B$ on the left of \eqref{eq:randommax},
\[
 \one_B\le (c_*s^2)^{-1}\int_{\T^2}
       \one_{\{\max_k|M_k(x)|>b/s\}}\dd x.
\]
Take expectations and use \eqref{eq:pointxmax}. For $b=A\sqrt m$,
$s^{-2}=2^{2m}m^{-6}$, and $A$ large enough that
$cA^2>2\log 2$, the upper bound is $o(1)$.
\end{proof}
No independence between the random partition and $W$ is assumed.


\section{Local divergence inverses and Whitney partitions}\label{sec:geom}

At each refinement step the new current is obtained by inverting the
divergence on an active parent cell and projecting onto the children of
that cell. This section constructs the two geometric ingredients. The
first is a ray inverse of the divergence, defined in an intrinsic chart
attached to the predictable routing metric and transported to physical
coordinates by a Piola transform. The second is a finite Whitney
partition of the parent in the same chart, in which clipped boundary
pieces are attached to full receiver squares. The section ends with
the projected-kernel bounds of Lemma~\ref{lem:pk}, which are the
geometric input for the covariance estimates of
Section~\ref{sec:selector}. Proofs of some of the geometric statements
are collected in Appendix~\ref{app:geometry}.

\paragraph{Physical and intrinsic coordinates.}
The current $U_k$, the two coordinate directions, and the phase $H(U_k)$
always live in the original physical coordinates on $\T^2$.
The local inverse of divergence, however, is constructed in a chart
chosen using the \emph{old} Gaussian history. Fix an active physical parent
$Q$, regarded as a convex set in one chosen lift of the torus, and let
$G_Q^+$ be its new predictable routing metric. Set
\begin{equation}\label{eq:intrinsicchart}
 A_Q=(G_Q^+)^{1/2}=A_Q^T,\quad \det A_Q=1,\quad
 z=A_Q^{-1}(x-b_Q),\quad
 \widehat Q=A_Q^{-1}(Q-b_Q),\quad
 \widehat S=A_Q^{-1}(S-b_Q).
\end{equation}
Here $b_Q$ is an old-measurable choice of origin, and all displayed
quantities are local to this parent. In particular, $A_Q$ is generally
not an automorphism of the torus. \phantomsection\label{not:intrinsicdistance}
The intrinsic distance is
\[
 d_{G_Q^+}(x,y)=|A_Q^{-1}(x-y)|,\qquad
 d_{G_Q^+}(x,\partial Q)=\dist(A_Q^{-1}(x-b_Q),\partial\widehat Q).
\]
The physical child $S=b_Q+A_Q\widehat S$ has $|S|=|\widehat S|$,
but need not be a Euclidean square (Figure~\ref{fig:intrinsic-physical}). All local radii, Whitney lengths,
H\"older seminorms and point-source kernel estimates below are in the
\emph{intrinsic} $z$-chart. There is no uniform Euclidean physical
kernel bound independent of the eccentricity of $A_Q$.

The transformed parent $\widehat Q$ is a uniformly shape-regular
convex parallelogram: before updating the metric the active parent is a
square in its birth chart, and successive metrics have a uniformly
bounded relative Loewner distortion. Let $D=\diam(\widehat Q)$ and
choose an inball $B(z_0,c_0D)\subset\widehat Q$.
\phantomsection\label{not:Kconstant}
Here $K>1$ denotes the fixed bound on distortion between a cell's
birth metric $G$ and the updated metric $G^+$:
$K^{-1}G\preceq G^+\preceq KG$ (see \eqref{eq:clipping}).
The inball ratio $c_0=c_0(K)>0$ is independent of the physical
eccentricity of the chart.

\begin{figure}[tbp]
\centering
\begin{tikzpicture}[>=Latex, font=\footnotesize, line join=round]
  % intrinsic chart
  \begin{scope}[shift={(0.5,0)}]
    \fill[black!3] (0,0.35)--(3.0,0.35)--(3.38,2.58)--(0.38,2.58)--cycle;
    \draw[line width=.65pt] (0,0.35)--(3.0,0.35)--(3.38,2.58)--(0.38,2.58)--cycle;
    \draw[densely dashed,gray] (2.28,1.60) circle (.45);
    \draw[fill=black!14, line width=.7pt] (.70,.70) rectangle (1.39,1.39);
    \node[anchor=north] at (1.65,.15) {$\widehat Q\subset\R^2_z$};
    \node at (1.045,1.045) {$\widehat S$};
    \node at (2.28,1.60) {$\scriptstyle z_0$};
  \end{scope}
  % affine arrow between panels
  \draw[->,line width=.7pt] (4.25,1.45)--(5.55,1.45);
  \node[align=center] at (4.9,1.92) {$x=b_Q+A_Qz$};
  \node[align=center] at (4.9,.91) {$\det A_Q=1$};
  % physical chart, affine image under an area preserving shear/stretch
  \begin{scope}[shift={(6.0,.1)},x={(1.18cm,0cm)},y={(.44cm,.84746cm)}]
    \fill[black!3] (0,.35)--(3,.35)--(3.38,2.58)--(.38,2.58)--cycle;
    \draw[line width=.65pt] (0,.35)--(3,.35)--(3.38,2.58)--(.38,2.58)--cycle;
    \draw[fill=black!14,line width=.7pt] (.70,.70)--(1.39,.70)--(1.39,1.39)--(.70,1.39)--cycle;
    \draw[densely dashed,gray] (2.28,1.60) circle (.45);
    \node at (1.045,1.045) {$S$};
  \end{scope}
  \node at (8.47,.24) {$Q\subset\R_x^2$};


\end{tikzpicture}
\caption{The predictable intrinsic chart. Whitney squares and kernel estimates are defined on $\widehat Q$; their physical images are generally parallelograms. The local affine map preserves area, but not Euclidean lengths. The dashed intrinsic ball becomes an ellipse in physical coordinates.}
\label{fig:intrinsic-physical}
\end{figure}
Take a fixed nonnegative smooth probability bump $\omega$ supported in
$B(0,c_0/4)$ and define
\[
 \widehat\omega_Q(w)=D^{-2}\omega((w-z_0)/D),\qquad
 \int_{\widehat Q}\widehat\omega_Q=1,\quad
 \|D^j\widehat\omega_Q\|_\infty\le C_jD^{-2-j}\quad(j=0,1).
\]
Define the \emph{intrinsic} ray inverse by
\begin{equation}\label{eq:ray}
 \widehat K_Q(z,w)=(z-w)\int_1^\infty t\widehat\omega_Q(w+t(z-w))\,dt,
 \qquad
 \widehat{\mathcal B}_Q\widehat f(z)
     =\int_{\widehat Q}\widehat K_Q(z,w)\widehat f(w)\,dw.
\end{equation}
For mean-zero $\widehat f$, integration along the ray proves
$\divv_z\widehat{\mathcal B}_Q\widehat f=\widehat f$ and
zero normal trace on $\partial\widehat Q$.
The \emph{physical} Piola kernel and inverse are
\begin{equation}\label{eq:piolakernel}
 K_Q^{\rm phys}(x,y)=A_Q\widehat K_Q(z,w),\qquad
 \mathcal B_Q^{\rm phys}f(x)=A_Q\widehat{\mathcal B}_Q\widehat f(z),
 \quad \widehat f(w)=f(b_Q+A_Qw),
\end{equation}
where $z=A_Q^{-1}(x-b_Q)$ and $w=A_Q^{-1}(y-b_Q)$.
Because $\det A_Q=1$, no Jacobian factor is missing. A change of
variables and the Piola identity give
\begin{equation}\label{eq:weakray}
 \int_Q\mathcal B_Q^{\rm phys}f\cdot\nabla_x\phi\,dx
 =\int_{\widehat Q}\widehat{\mathcal B}_Q\widehat f
       \cdot\nabla_z\widehat\phi\,dz
 =-\int_Qf\phi\,dx,\qquad \int_Q f=0.
\end{equation}
The zero extension across the physical boundary has no flux. The inverse
acts on the convex parent, not on its potentially disconnected terminal
children.

\paragraph{Intrinsic kernel bounds.}
The scaled bump gives, for $z\ne w$,
\begin{equation}\label{eq:singular}
 |\widehat K_Q(z,w)|\le\frac{C}{|z-w|},\qquad
 |D_z\widehat K_Q(z,w)|+|D_w\widehat K_Q(z,w)|
      \le\frac{C}{|z-w|^2}.
\end{equation}
Furthermore,
\begin{equation}\label{eq:support}
 \widehat K_Q(z,w)=0\quad\hbox{if}\quad
 |z-w|>C\dist(z,\partial\widehat Q).
\end{equation}
Both properties are proved in Lemma~\ref{lem:app-ray}
(Appendix~\ref{app:kernelchain}): \eqref{eq:singular} follows by
differentiating the normalized bump under the integral, and
\eqref{eq:support} from convexity, since a ray point
$w+t(z-w)$ in the central bump forces
$\dist(z,\partial\widehat Q)\ge cD/t$ while $|z-w|\le D/t$.
In physical variables the uniform bounds read
\[
 |A_Q^{-1}K_Q^{\rm phys}(x,y)|\le\frac{C}{d_{G_Q^+}(x,y)},
 \quad
 |A_Q^{-1}D_xK_Q^{\rm phys}(x,y)A_Q|
 +|A_Q^{-1}D_yK_Q^{\rm phys}(x,y)A_Q|
 \le\frac{C}{d_{G_Q^+}(x,y)^2}.
\]
The support condition is likewise expressed in $d_{G_Q^+}$.
For example $A_Q=\operatorname{diag}(L,L^{-1})$ makes the raw physical
kernel as large as $L$ times its intrinsic counterpart, and the bound
$|K_Q^{\rm phys}(x,y)|\lesssim|x-y|^{-1}$ can fail by a factor
of order $L^2$. Determinant one preserves area, not Euclidean lengths.


\subsection{Whitney partitions in intrinsic coordinates}
Truncating a Whitney decomposition and completing it by boundary
sets of diameter $O(s)$ and area comparable to $s^2$ is the
geometric construction of \cite[Lemma~5.1 and Appendix~A]{umc:AGT22},
also recalled in \cite[Section~2.3, Lemma~2.1]{GT-combinatorial}.
Here it is carried out inside each predictably chosen intrinsic
parent. The explicit receiver construction below also gives
the child-size cap and the constants, uniform under successive
changes of chart, which the adaptive construction requires.

\phantomsection\label{not:whitneyparameters}
The constants $\theta\in(0,1)$ and $D_0>0$ denote, respectively,
the upper fraction of a parent's intrinsic diameter allowed for a
Whitney square and the terminal stopping threshold $D\le D_0s$.
The distinct contraction parameter $\vartheta\in(0,1)$ bounds the
\emph{birth-chart} size ratio at an actual refinement:
$\ell_{j+1}\le\vartheta\ell_j$, with one corresponding
diameter contraction for a newly terminal union.  It is achieved by
choosing $\theta$ small and $D_0$ large after $K$ is fixed;
Lemma~\ref{lem:birthchart} gives the precise constants.

\begin{lemma}\label{lem:whitney}
Fix a shape parameter $c_0>0$ and a sufficiently small fixed $\theta>0$.
There are $\gamma,D_0,C$ depending only on $c_0,\theta$ so that, whenever
$D>D_0s$, the \emph{intrinsic} parallelogram $\widehat Q$ admits a disjoint
measurable partition into children $\widehat S$ with scales $\ell_S\ge s$:
\begin{equation}\label{eq:whitney}
 |\widehat S|\asymp \ell_S^2,\quad
 \diam(\widehat S)\le C\ell_S,\quad
 \ell_S\asymp\max\{s,\dist(z,\partial\widehat Q)\}
       \quad(z\in\widehat S).
\end{equation}
Nonterminal children are \emph{intrinsic} dyadic squares with
$\ell_S\le\theta D$; terminal children have $\ell_S=s$, area at least
$s^2$ and intrinsic diameter at most $Cs$.
After choosing $D_0$ sufficiently large, the maximal intrinsic child
diameter is a fixed small fraction of $D$. The physical partition
consists of $S=b_Q+A_Q\widehat S$.
If $D\le D_0s$, the parent is left terminal.
\end{lemma}

Since $\det A_Q=1$, the physical children have the same areas as the
intrinsic ones, but not in general the same Euclidean diameters.

\begin{proof}
Throughout this proof write $Q,S,R,x$ for \emph{intrinsic} subsets and
points of $\widehat Q$; they are not physical coordinates.
Use a dyadic lattice in this predictable intrinsic chart, with smallest
squares of side $s$.
Take \emph{all inclusion-maximal} dyadic squares $R\subset Q$ with side
$\ell_R\in[2s,\theta D]$ and
$\dist(R,\partial Q)\ge\gamma\ell_R$.
Since dyadic squares are nested or disjoint, these maximal squares have
disjoint interiors. A maximal square's dyadic parent fails admissibility:
it either is not contained in $Q$, violates the distance condition, or
exceeds $\theta D$. These three alternatives imply
$\dist(x,\partial Q)\lesssim_{\gamma,\theta}\ell_R$ for
$x\in R$. The reverse inequality follows from admissibility. More explicitly,
for a square of side $\ell$, if its $2\ell$ parent is not contained
in $Q$, every point of the child is within $2\sqrt2\ell$ of the
boundary; if the parent fails the distance requirement, the same
bound is $(2\gamma+2\sqrt2)\ell$; if it exceeds the cap, then
$D<2\ell/\theta$. This proves the upper Whitney comparison, with
constants fixed before the Gaussian experiment, so the
nonterminal squares satisfy \eqref{eq:whitney}.

Every remaining point of $Q$ belongs to an uncovered $s$-lattice square.
If this $s$-square is wholly inside $Q$, retain it as a full uncovered
square. Its $2s$-ancestor cannot be admissible: otherwise that ancestor would
lie in some inclusion-maximal admissible square and would cover the
$s$-square. The cap cannot fail since $\theta D>4s$.
If the ancestor fails containment or boundary clearance, every point of
the uncovered $s$-square lies within $(2\gamma+3\sqrt2)s$ of
$\partial Q$. In particular every full uncovered square is
within $Cs$ of the boundary of $Q$.

It remains to assign a \emph{clipped} piece $P=Q\cap R$, where an
$s$-square $R$ meets $\partial Q$, to a full uncovered $s$-square.
Choose $x\in P$. Then $\dist(x,\partial Q)\le\sqrt2s$.
\phantomsection\label{not:receiverconstants}
Let $z_0$ be the center of an inball $B(z_0,c_0D)\subset Q$ and set
$x_t=(1-t)x+tz_0$ with $t=A s/D$. Choose a constant $A=A(c_0)$ large
enough that $c_0As>5s$, and choose $D_0>2A$ to ensure $t<1/2$.
Convexity and the supporting-line characterization of boundary distance give
\[
 5s<\dist(x_t,\partial Q)\le(\sqrt2+A)s,
 \qquad |x_t-x|\le As.
\]
The $s$-lattice square $R'$ containing $x_t$ is therefore fully inside $Q$.
Take $\gamma$ larger than $(A+\sqrt2)/2$ (and enlarge it by an absolute
constant if necessary). Then $R'$ cannot lie in an admissible
dyadic square of side at least $2s$, because any such square would put
$x_t$ at boundary distance at least $2\gamma s$. Hence $R'$ is an
uncovered full receiver within $O(s)$ of $P$.

Choose one receiver by a deterministic tie-break rule for every clipped
$s$-square and join its clipped piece to that receiver. A receiver can get
only $O_A(1)$ pieces because their originating lattice squares lie within
$O(As)$; the union has area at least $s^2$, diameter $O_A(s)$, and all
points have boundary distance $O_A(s)$. Any full uncovered square not
receiving a clipped piece is a terminal child by itself. The maximal
squares and these terminal unions partition $Q$ modulo a fixed
measure-zero boundary convention; see Figure~\ref{fig:whitney-receivers}.

Choose $D_0$ only after $A,\gamma,\theta$ and the terminal diameter
constant $C_A$ have been fixed, so that $C_As\le\theta D$ whenever
$D>D_0s$. The exact old/new-chart side ratio and the metric transport factor
are recorded in Lemmas~\ref{lem:birthchart} and
\ref{lem:ancestorindexed}. In particular, choose $\theta$ and then
$D_0$ so that the birth-side shrink factor $\vartheta$ in those
lemmas satisfies $\vartheta\sqrt K<1$.
All selections can be made by a deterministic lexicographic rule.
The construction terminates in at most $C_K\log(1/s)=O(m)$
steps: every nonterminal square has intrinsic area at most a fixed
fraction of its parent's area, and every nonempty cell has area
at least $c s^2$. No lower bound on the area of an individual
clipped fragment is needed.
\end{proof}

\begin{figure}[tbp]
\centering
\begin{tikzpicture}[>=Latex,font=\footnotesize,line join=round]
  \def\Pleft{1.0}
  \def\Pright{9.1}
  % convex oblique parent; draw a clipped intrinsic s grid
  \begin{scope}
    \clip (1.0,.55)--(9.05,.55)--(7.67,4.45)--(-.38,4.45)--cycle;
    \fill[black!3] (-.5,.5) rectangle (9.2,4.6);
    \draw[step=.46cm,very thin,black!22] (-.5,.3) grid (9.35,4.6);
    \fill[black!12] (2.26,2.14) rectangle (3.18,3.06);
    \draw[line width=.7pt] (2.26,2.14) rectangle (3.18,3.06);
    % clipped boundary piece of a small lattice square
    \fill[black!37] (8.70,.76) rectangle (9.16,1.22);
    % an uncovered full receiver (schematic size s)
    \fill[black!22] (7.32,1.22) rectangle (7.78,1.68);
    \draw[line width=.85pt] (7.32,1.22) rectangle (7.78,1.68);
  \end{scope}
  \draw[line width=.85pt] (1.0,.55)--(9.05,.55)--(7.67,4.45)--(-.38,4.45)--cycle;
  \draw[dashed,black!65] (4.2,2.45) circle (.85);
  \fill (4.2,2.45) circle (.035);
  \node[anchor=south west] at (4.25,2.49) {$z_0$};
  % extrapolation of the ray from clipped piece towards inball center
  \coordinate (x) at (8.84,.93);
  \coordinate (xt) at (7.68,1.31);
  \draw[->,line width=.85pt] (x)--(xt);
  \draw[densely dashed,gray] (xt)--(4.2,2.45);
  \fill (x) circle (.042);
  \fill (xt) circle (.042);
  \node[anchor=south west] at (8.88,.95) {$x$};
  \node[anchor=south] at (7.55,1.75) {$x_t$};
  \draw[->,thin] (2.72,3.57)--(2.72,3.09);
  \node[anchor=south] at (2.72,3.59) {interior dyadic square};
  \draw[->,thin] (9.15,1.92)--(8.91,1.17);
  \node[anchor=west,align=left] at (9.13,2.01) {clipped\\piece};
  \draw[->,thin] (6.59,.32)--(7.35,1.22);
  \node[anchor=north,align=center] at (6.20,.27) {full $s$-square receiver};
\end{tikzpicture}
\caption{Clipped-square receivers in an intrinsic chart (schematic, not to scale). A boundary $s$-square may have an arbitrarily small intersection with $\widehat Q$. Moving $x$ toward the inball center gives $x_t$, contained in an uncovered full $s$-square. The clipped fragment is assigned to that receiver; only the resulting union is a terminal child and has area at least $s^2$.}
\label{fig:whitney-receivers}
\end{figure}

The half-open boundary convention and the finite receiver assignment are
spelled out in Lemma~\ref{lem:exactcover}.


\subsection{Boundary receivers and the choice of constants}
We fix the constants in the following order.
Fix an inball ratio $c_0$. First take $A>5/c_0$ and
$\gamma>(A+\sqrt2)/2$. Next choose $\theta>0$ small enough for
subsequent metric contraction. Finally take $D_0$ so large that
$D_0>2A$, $\theta D_0>4$, and every $O_A(s)$ terminal union has
diameter at most a prescribed small fraction of $D$ in the current chart. The constants
in the Whitney bounds may depend on $c_0,A,\gamma,\theta$ but
\emph{not} on $D/s$, the lattice translation, or the physical
eccentricity of a determinant-one chart.

The two kinds of children behave differently. A maximal admissible
square has its full area $\ell^2$, whereas a clipped mesh piece may
have arbitrarily small area. The receiving full $s$-square provides
the area floor for the union, and only this union is a child of
the partition. At an acute corner the receiving square may lie
along a different face, but the inball displacement gives
bounded distance and bounded multiplicity without any
angle-dependent connectivity hypothesis. The operator
$\mathcal B_Q$ is applied to the convex parent, not to a terminal
union. Only the area, diameter and Whitney bounds of the children
enter the projected kernel estimates, so neither connectedness nor
a John constant for terminal unions is used.

\paragraph{Lattice conventions.}
All squares are half-open after applying one fixed lexicographic
boundary convention. The $s$-lattice belongs to the predictable
intrinsic chart, not to the fixed physical grid; it is translated
according to a preassigned measurable rule. At exact lattice contacts we choose
one closed representative for the support test and assign
measure-zero interfaces arbitrarily. No Gaussian sample is
consulted when making any of these choices. A parent with
$D\le D_0s$ is terminal, and no inverse is applied on its
children. Lemma~\ref{lem:whitney} is thus a finite, measurable
construction rather than a limiting Whitney decomposition.


\paragraph{Notation for the projected kernels.}
\phantomsection\label{not:holderparameters}
We use $0<\alpha<1$ for the projected-kernel H\"older exponent.
The Gaussian field and its spatial memory use a fixed smaller exponent
$0<\beta<\alpha$; constants may depend on both exponents.
Here $G_Q^+$ is the updated physical covariance metric, $A_Q=(G_Q^+)^{1/2}$,
$z=A_Q^{-1}(x-b_Q)$ and $\widehat S=A_Q^{-1}(S-b_Q)$.  The length
$\ell_S$ and every distance in Lemma~\ref{lem:pk} are measured in
this updated intrinsic chart.  The averaging operator $\widehat P$
is intrinsic; a physical vector is recovered only after multiplying
by $A_Q$.  The fixed physical distance of the torus occurs only in
Section~\ref{sec:terminal}.

The cell-averaged singular-kernel estimates are related to the
near- and far-cell bounds in
\cite[Section~2.4, Lemma~2.6]{GT-combinatorial} and to the
covariance calculation for the histogram correction in
\cite[Section~4, Step~3]{umc:AGT22}. In the present argument,
the boundary-distance support \eqref{eq:support} and the exact
Gaussian covariance identity \eqref{eq:gramZ} give a pointwise
conditional variance bound uniform in the intrinsic scale.
The projected-kernel bounds needed for this are the following.

\begin{lemma}\label{lem:pk}
\phantomsection\label{not:childprojector}
Let $\widehat P$ average over the children $\widehat S$ in the
intrinsic chart $\widehat Q$, and let $z,z'$ belong to one $\widehat S$
of intrinsic scale $\ell$. For every fixed $0<\alpha<1$,
\begin{align}\label{eq:pk}
 \|\widehat P\widehat K_Q(z,\cdot)\|_2&\le C,&
 \|\widehat P\widehat K_Q(z,\cdot)\|_\infty&\le C\ell^{-1},\\
 \|\widehat P(\widehat K_Q(z,\cdot)-\widehat K_Q(z',\cdot))\|_2
  &\le C(|z-z'|/\ell)^\alpha,&
 \|\widehat P(\widehat K_Q(z,\cdot)-\widehat K_Q(z',\cdot))\|_\infty
  &\le C\ell^{-1}(|z-z'|/\ell)^\alpha.\nonumber
\end{align}
The same inequalities hold after subtracting the parent-mean
projection. They are uniform in the intrinsic chart.
\end{lemma}

Lemma~\ref{lem:pk} concerns the intrinsic kernel only; the passage
to physical vectors is discussed after \eqref{eq:projectedpiola}.

\begin{proof}
In this proof, for compactness, write $Q,S,T,x,y$ and $K_Q,P$
for the intrinsic objects $\widehat Q,\widehat S,\widehat T,z,w$
and $\widehat K_Q,\widehat P$. Set $d_Q(x)=\dist(x,\partial Q)$.
The Whitney bounds give $d_Q(x)\le C_W\ell$ for $x\in S$;
by the support restriction, $K_Q(x,\cdot)$ is supported in
$B(x,C\ell)$. If a child $T$ meets this ball, choose
$y\in T\cap B(x,C\ell)$. The distance-to-boundary function is
$1$-Lipschitz, so
$\ell_T\le C_W\max\{s,d_Q(y)\}\le C\ell$.
Thus $\diam T\le C\ell$ and the union of projected-support
children lies in $B(x,C'\ell)$ and has area at most $C\ell^2$.

For the one-point $L^\infty$ bound we separate near and far
children. Choose $c_1>0$ smaller than
the inverse of $10$ times every fixed Whitney comparison constant.
For any child $T$ meeting $B(x,c_1\ell)$,
\begin{equation}\label{eq:nearchild}
 |T|\ge c\ell^2.
\end{equation}
Indeed, if $\ell=s$ this is the global area floor $|T|\ge s^2$.
If $\ell>s$, then $S$ is a nonterminal square and admissibility
implies $d_Q(x)\ge\gamma\ell$; the $1$-Lipschitz property gives
$d_Q(y)\ge\gamma\ell/2$ for $y\in B(x,c_1\ell)$.
The pointwise Whitney comparison on $T$ implies $\ell_T\ge c\ell$,
and $|T|\ge c\ell_T^2$.
For these near children,
\[
 \left|\fint_T K_Q(x,y)\,dy\right|
 \le \frac{C}{|T|}\int_{B(x,C'\ell)}\frac{dy}{|x-y|}
 \le \frac{C\ell}{\ell^2}=\frac{C}{\ell}.
\]
On each other projected-support child $T$, all $y\in T$ satisfy
$|x-y|\ge c_1\ell$, so the same bound follows pointwise from
$|K_Q(x,y)|\le C/|x-y|$. Multiplication by the square root of
the $O(\ell^2)$ projected-support area proves the $L^2$ estimate.

For the two-point estimates put $d=|x-x'|$. If $d\ge c_1\ell/4$,
apply the two one-point estimates and enlarge the constant.
Suppose $0<d<c_1\ell/4$. The combined supports lie in
$B(x,C''\ell)$. In the two disks of radius $2d$ about the two
poles, the two singular kernels have $L^1$ norm $O(d)$.
Outside the disks, the mean-value estimate for
$\nabla_xK_Q$ gives
\begin{equation}\label{eq:twopole}
 \int_Q|K_Q(x,y)-K_Q(x',y)|\,dy
 \le Cd\bigl(1+\log(\ell/d)\bigr).
\end{equation}
For a projected child meeting $B(x,c_1\ell/2)$, use
\eqref{eq:nearchild} and divide \eqref{eq:twopole} by $|T|$.
For any other child the segment joining $x,x'$ remains at
distance at least $c_1\ell/4$ from that child; integrate the
pointwise gradient bound to get $Cd/\ell^2$.
Hence
\[
 \|P(K_Q(x,\cdot)-K_Q(x',\cdot))\|_\infty
 \le C\ell^{-1}\frac d\ell
       \bigl(1+\log(\ell/d)\bigr)
 \le C_\alpha\ell^{-1}(d/\ell)^\alpha.
\]
The projected support of the difference has area $O(\ell^2)$,
which gives its $L^2$ bound. Subtracting the parent-constant
projection changes the bounds by at most a fixed factor, since
$P_QP=P_Q$ and a constant projection is an $L^2$ contraction;
its $L^\infty$ norm is bounded by that of the projected kernel.
The same two-pole calculation on a whole parent gives the
translation estimate
\begin{equation}\label{eq:L1mod}
 \int_Q|K_Q(x,y)-K_Q(x',y)|\,dy
 \le Cd\bigl(1+\log(D/d)\bigr),\qquad 0<d\le D/2.
\end{equation}
\end{proof}

\paragraph{Projected Piola evaluations.}
Let $P^{\rm phys}$ average over the physical children
$S=b_Q+A_Q\widehat S$. Since $\det A_Q=1$, projection commutes
with pullback and multiplication by the old-measurable matrix:
\begin{equation}\label{eq:projectedpiola}
 A_Q^{-1}P^{\rm phys}K_Q^{\rm phys}(x,\cdot)
 =\bigl(\widehat P\widehat K_Q(z,\cdot)\bigr)
       \circ A_Q^{-1}(\cdot-b_Q),\qquad z=A_Q^{-1}(x-b_Q).
\end{equation}
The $L^2$ and $L^\infty$ norms of the normalized vector
$A_Q^{-1}P^{\rm phys}K_Q^{\rm phys}$ are therefore controlled by
\eqref{eq:pk}. The corresponding raw physical vector can cost a factor
$\|A_Q\|_{\rm op}$; this factor must remain in physical covariance
and componentwise bounds. The parent mean subtraction enjoys the same
commutation property.

\section{Anisotropic covariance selection}\label{sec:selector}

This section defines the target metric $C(y)$ toward which the
covariance of each fresh Gaussian increment is steered, and the lag
$L(y,G)$ measuring how far a current metric $G$ is from it.
Lemma~\ref{lem:selector} collects the bounds on $C$, on $L$ and on
the spectral clipping $\mathcal T_C$ that the bootstrap of
Section~\ref{sec:bootstrap} uses.

Recall the phase $H$ and its gradient $v$ from
\eqref{eq:phaseH}.
The continuous-time exponential Bellman trial has normalized Hessian
$\zeta_{\rm B}D^2H+\zeta_{\rm B}^2vv^T$, as recalled in
\eqref{eq:intro-bellman-determinant}. For the discrete construction
we use instead a regularized Hessian proxy, whose diagonal
part also dominates the curvature of $H$. On the critical line its
determinant is comparable to $H^2$, so there is no large-amplitude
factor $H^{-\sigma}$ to spare. The metric must therefore be
adjusted dynamically.

Define
\begin{equation}\label{eq:B}
 B(y)=\operatorname{diag}\bigl(H^{1-2a_1},H^{1-2a_2}\bigr)+v(y)v(y)^T,
 \qquad C(y)=\sqrt{\det B(y)}B(y)^{-1}.
\end{equation}
In two dimensions, $\det C=1$. To interpret this choice, for any
positive definite $G$ with $\det G=1$, the arithmetic--geometric
mean inequality applied to $B^{1/2}GB^{1/2}$ gives
\[
 \tr(BG)\ge2\sqrt{\det B}=\tr(BC),
\]
with equality precisely when $G=C$. We measure the mismatch of an old
metric $G$ from the target by
\begin{equation}\label{eq:lag}
 L(y,G)=\sqrt{\tr\bigl(GC(y)^{-1}\bigr)}.
\end{equation}
Since $C^{-1}=B/\sqrt{\det B}$, the lag is exactly the relative
covariance cost:
\begin{equation}\label{eq:lag-cost}
 \frac{\tr(B(y)G)}{\tr(B(y)C(y))}=\frac{L(y,G)^2}{2}.
\end{equation}
Thus $L\ge\sqrt2$, with equality precisely at $G=C(y)$.
The matrix $G$ will bound the fresh physical covariance from above
(see \eqref{eq:covphysical}), whereas the spatial distance is defined
by $G^{-1}$, as in \eqref{eq:intrinsicchart}.

\begin{lemma}\label{lem:selector}
On \eqref{eq:crit}, with constants depending on the fixed exponents,
\begin{equation}\label{eq:selectors}
 \det B\asymp H^2,\qquad C_{ii}\lesssim H^{2a_i},\qquad
 \frac{vv^T}{H}\preceq C_0C(y)^{-1}.
\end{equation}
If $G_{ii}\le M H(y)^{2a_i}$ and $\det G=1$, then
\begin{equation}\label{eq:lagder}
 L(y,G)^2\le C_M H(y),\qquad
 |\partial_{y_i}L(y,G)|\le C_M H(y)^{1/2-a_i}.
\end{equation}
For a fixed $K>1$, the determinant-one spectral clipping
\[
 \mathcal T_C(G)=G^{1/2}
  \operatorname{clip}_{[K^{-1},K]}\bigl(G^{-1/2}CG^{-1/2}\bigr)G^{1/2}
\]
satisfies, with $G^+=\mathcal T_C(G)$,
\begin{equation}\label{eq:clipping}
 K^{-1}G\preceq G^+\preceq KG,\quad
 G^+\preceq K^{-1}G+C,\quad
 L(y,G^+)\le K^{-1/2}L(y,G)+C_K.
\end{equation}
\end{lemma}
\begin{proof}
Write $D_i=H^{1-2a_i}$. The rank-one determinant formula gives
\[
 \det B=D_1D_2\left(1+\sum_i\frac{v_i^2}{D_i}\right)
 =H\left(1+\sum_i v_i^2H^{-1+2a_i}\right)\asymp H^2,
\]
where the last comparison follows by examining the dominant summand
$F_i(y_i)\ge H/2$. The formula for the inverse of a $2\times2$
positive matrix gives $C_{ii}=B_{jj}/\sqrt{\det B}\lesssim H^{2a_i}$.
Also
\[
 C^{-1}=\frac{B}{\sqrt{\det B}}
 =\chi\operatorname{diag}(H^{-2a_i})+ww^T,
 \quad \chi=\frac H{\sqrt{\det B}}\asymp1,
 \quad w=\frac v{(\det B)^{1/4}}.
\]
One has $|\partial_i\chi|\lesssim H^{-a_i}$ and
$|\partial_i w_j|\lesssim H^{1/2-a_i-a_j}$.
Represent $L$ as the Euclidean norm of the direct sum of
$\sqrt\chi\,G_{ii}^{1/2}H^{-a_i}$ and $G^{1/2}w$.
Under the diagonal caps, differentiation of this vector proves
\eqref{eq:lagder} without dividing by a possibly small $L$.

The two eigenvalues of $G^{-1/2}CG^{-1/2}$ are reciprocal, and
clipping both to $[K^{-1},K]$ preserves their
product. The inequality $\operatorname{clip}(t)\le K^{-1}+t$ yields the
Loewner bound. Along the common eigendirections the logarithmic
mismatch is shortened by $\log K$. Separating the regimes of large and
bounded log-mismatch, this gives
$L(y,G^+)\le K^{-1/2}L(y,G)+C_K$.
\end{proof}

\section{The adaptive construction and the bootstrap}\label{sec:bootstrap}

We now build the adaptive field $U_k$ generation by generation and
prove that it stays in the regime where the estimates of
Sections~\ref{sec:geom} and~\ref{sec:selector} apply. At each active
parent the metric of the next increment is obtained by clipping the
current metric toward the target $C$ of Section~\ref{sec:selector};
the fresh Gaussian increment is then pushed forward by this metric.
The main result, Proposition~\ref{prop:bootstrap}, shows by induction
that, as long as no envelope is too large and the modular has not
overshot, the diagonal caps, the relative phase comparisons and a
contraction of the lag-plus-memory functional $S_k$ hold
simultaneously. These are the inputs of the exponential moment
estimates of Section~\ref{sec:phase}.

Throughout, every stage $k\to k+1$ has the same order. First the
geometry is chosen from the old history: the representative, the
metric $G_Q^+$, the intrinsic parent and its Whitney partition, and
the stopping decisions are $\F_k$-measurable. Only then are the new
Gaussian child masses revealed. Every quantity called old,
predictable or old-measurable below is fixed before the noise of the
current stage.

\phantomsection\label{not:birthmetric}
Initialize $U_0=0,G_0=I,\pr_0=\{Q_0\}$. At an active
physical parent $Q$ in a chosen lift, choose a representative $x_Q$
by a fixed old-measurable rule, and set
$C_Q=C(U_k(x_Q))$, $G_Q^+=\mathcal T_{C_Q}(G_k)$ and
$A_Q=(G_Q^+)^{1/2}$. This matrix, the intrinsic parent
$\widehat Q$, its bump and its Whitney partition are all fixed
before the new child masses are revealed. After revelation let
$\Delta f_k=(P_{\pr_{k+1}}-P_{\pr_k})W$ be the physical,
parentwise mean-zero step density and put
\phantomsection\label{not:intrinsicgauss}
\[
 \widehat{\Delta f}_k(w)=\Delta f_k(b_Q+A_Qw),\qquad
 Z_{k+1}(z)=\int_{\widehat Q}\widehat K_Q(z,w)
                        \widehat{\Delta f}_k(w)\,dw.
\]
Thus $Z_{k+1}$ is an intrinsically normalized Gaussian
vector field. The physical increment is its Piola push-forward:
\begin{equation}\label{eq:update}
 V_{k+1}(x)=\sqrt\delta\,A_Q
       Z_{k+1}\bigl(A_Q^{-1}(x-b_Q)\bigr),\qquad
 U_{k+1}=U_k+V_{k+1},\qquad
 \divv_x U_{k+1}=\sqrt\delta\,f_{k+1}.
\end{equation}
Every such increment has zero normal trace on its physical parent,
so its zero extension is globally periodic. In particular the local
matrix $A_Q\in\mathrm{SL}(2,\mathbb R)$ is never used as a global
coordinate change on $\T^2$.

\begin{lemma}
\label{lem:piolaGaussian}
Conditional on the old filtration and on the predictably selected
chart, $Z_{k+1}(z)$ is a centered Gaussian vector with covariance
matrix
\begin{equation}\label{eq:gramZ}
 \operatorname{Cov}_k\bigl(Z_{k+1}(z)\bigr)_{ab}
 =\int_{\widehat Q}
 \bigl[(\widehat P-\widehat P_{\widehat Q})
       \widehat K_{Q,a}(z,\cdot)\bigr](w)
 \bigl[(\widehat P-\widehat P_{\widehat Q})
       \widehat K_{Q,b}(z,\cdot)\bigr](w)\,dw.
\end{equation}
In particular,
\begin{equation}\label{eq:covphysical}
 \operatorname{Cov}_k Z_{k+1}(z)\preceq CI,
 \qquad
 \operatorname{Cov}_k V_{k+1}(x)\preceq C\delta G_Q^+,
 \qquad
 \divv_x V_{k+1}=\sqrt\delta\,\Delta f_k.
\end{equation}
\end{lemma}
\begin{proof}
Fix the old history. The chart $A_Q$ and the partition are
then deterministic. Pulling back a white-noise test $g$ by
$x=b_Q+A_Qw$ preserves its $L^2$ inner product: the Jacobian is
$\det A_Q=1$. The conditional increment $\Delta f_k$ is
precisely the orthogonal Gaussian projection onto the span of
child-constant functions with their parent constant removed.
Pairing it with the two components of $\widehat K_Q(z,\cdot)$
gives the Gram matrix \eqref{eq:gramZ}. Lemma~\ref{lem:pk}
provides the intrinsic $L^2$ bound, so the Gram matrix is
$\preceq CI$. The covariance identity
$\operatorname{Cov}_k(V)=\delta A_Q\operatorname{Cov}_k(Z)A_Q^T$
then proves the physical bound, since $A_QA_Q^T=G_Q^+$.
For the divergence, $D_x V=\sqrt\delta\,A_Q(D_zZ)A_Q^{-1}$,
so $\operatorname{tr}D_xV=\sqrt\delta\operatorname{tr}D_zZ
=\sqrt\delta\widehat{\Delta f}_k$.
Equivalently, the weak identity \eqref{eq:weakray} proves the
result without pointwise differentiation. The Piola normal-flux
identity preserves the no-flux boundary condition.
\end{proof}

The physical field does not, in general, have covariance bounded
by $C\delta I$; its covariance is controlled only through the
anisotropic metric $G_Q^+$. The next lemma gives uniform
sub-Gaussian control of the normalized field $Z_{k+1}$ and of its
H\"older seminorm on each child.

\begin{lemma}\label{lem:envelope}
Fix $0<\beta<\alpha<1$. At generation $k$, condition on
$\F_k$. Let $Q$ be an active parent, let its intrinsic chart and
Whitney children be selected $\F_k$-measurably as above, and let
$Z_{k+1}$ be the centered intrinsic field in
\eqref{eq:update}. For each newly revealed child
$S=b_Q+A_Q\widehat S$ of intrinsic scale $\ell_S$, choose the
continuous version of $Z_{k+1}$ on $\overline{\widehat S}$ and set
\begin{equation}\label{eq:envelope}
 E_{k+1,S}:=
 \sup_{z\in\overline{\widehat S}}|Z_{k+1}(z)|
 +\ell_S^\beta[Z_{k+1}]_{C^\beta(\overline{\widehat S})}.
\end{equation}
There are deterministic $c,C>0$, depending only on
$\alpha,\beta$ and the uniform intrinsic Whitney and kernel
constants, such that, for every $t\ge0$,
\begin{equation}\label{eq:envelopetail}
 \Pp(E_{k+1,S}>t\mid\F_k)\le Ce^{-ct^2},
 \qquad \E_k e^{cE_{k+1,S}^2}\le C,
\end{equation}
possibly after reducing $c$. These estimates hold for each
$\F_k$-measurably selected child and are uniform in $k,m$,
its intrinsic scale, and the eccentricity of $A_Q$.
The H\"older seminorm in \eqref{eq:envelope} is computed with the
intrinsic $z$-distance; its physical pullback uses
$d_{G_Q^+}(x,x')=|A_Q^{-1}(x-x')|$.

If $J_0\le C_{\rm dep}m$ is the deterministic reveal-depth bound,
$s=2^{-m}m^3$, and each child has area at least $c_*s^2$,
then for a sufficiently large fixed \emph{global envelope cutoff}
$R>0$,
\begin{equation}\label{eq:globalcut}
 \Pp\!\left(\exists\text{ an actually attempted reveal }(k+1,S):
                 E_{k+1,S}>R\sqrt m\right)
 \le C J_0s^{-2}e^{-cR^2m}
 \le e^{-c_Rm}
\end{equation}
for all sufficiently large $m$, with $c_R>0$.
\end{lemma}
\begin{proof}
We work conditionally on $\F_k$ and write
$\Pp_k(\cdot)=\Pp(\cdot\mid\F_k)$; in particular the chart,
$\widehat Q$, the averaging projector and the child
$\widehat S$ are deterministic. By the Gram formula
\eqref{eq:gramZ} and the projected evaluations \eqref{eq:pk},
\begin{equation}\label{eq:envelopecanonical}
 \sup_{z\in\widehat S}\E_k|Z_{k+1}(z)|^2\le C_0,
 \qquad
 \E_k|Z_{k+1}(z)-Z_{k+1}(z')|^2
 \le C_0\left(\frac{|z-z'|}{\ell_S}\right)^{2\alpha}
 \quad(z,z'\in\widehat S).
\end{equation}
The same bounds extend to the closure by continuity. Both
components are centered Gaussian. The constants do not see the
Piola matrix, because \eqref{eq:envelopecanonical} concerns the
normalized field $Z_{k+1}$.

We bound $E_{k+1,S}$ by a chaining argument.
Set $\ell=\ell_S$ and fix an exponent
$\gamma$ with $\beta<\gamma<\alpha$. The intrinsic diameter
bound $\operatorname{diam}(\widehat S)\le C_W\ell$ allows us to
construct, for every integer $j\ge0$, a finite
$2^{-j}\ell$-net $N_j\subset\overline{\widehat S}$ with
$|N_j|\le C2^{2j}$: take a fixed countable dense set in the
child, run a deterministic greedy procedure with a fixed order,
and retain the finitely many selected points. Since the intrinsic
chart is predictable, the nets are measurable with respect to the
old history. Each point $z\in\overline{\widehat S}$ has a
representative $z_j\in N_j$ at distance at most $C2^{-j}\ell$.
The nets need not be nested, nor need $\widehat S$ be connected.

At level zero test the two Gaussian coordinates at all points of
$N_0$. At every level $j\ge0$ test all coordinate differences
between pairs of points in $N_j\cup N_{j+1}$ whose separation is
at most $C_1 2^{-j}\ell$, where $C_1$ is chosen to cover both
consecutive-level approximations and pairs of nearest
representatives of nearby points. There are at most $C2^{4j}$
such tests. By \eqref{eq:envelopecanonical}, each tested
difference has conditional standard deviation at most
$C2^{-\alpha j}$. For $A\ge1$ and $u\ge0$, let $\mathcal G_u$ be the event
on which every level-zero coordinate has absolute value at most
$A(1+u)$ and every tested difference at level $j$ has absolute
value at most $A(1+u)2^{-\gamma j}$. The scalar Gaussian tail
bound and a union bound give
\begin{align}\label{eq:chainingfailure}
 \Pp_k(\mathcal G_u^c)
 &\le Ce^{-cA^2(1+u)^2}
 +C\sum_{j\ge0}2^{4j}
  \exp\!\left[-cA^2(1+u)^2
               2^{2(\alpha-\gamma)j}\right]
 \le C'e^{-c'u^2}.
\end{align}
Here $A$ is fixed sufficiently large. The series converges
uniformly: since $\alpha-\gamma>0$, its exponential factor
dominates $2^{4j}$, and $(1+u)^2\ge1+u^2$ separates the
$u$-tail. The increment threshold is $2^{-\gamma j}$, without a
factor growing with $j$; this is what yields a genuine $C^\beta$
bound.

On $\mathcal G_u$, telescope the values at the
representatives $z_j$ of a point $z$ from level zero to
arbitrarily fine levels. Consecutive representatives are
$O(2^{-j}\ell)$ apart, so every increment occurs among the
tested pairs. The convergent geometric series
$\sum_{j\ge0}2^{-\gamma j}$ gives
\[
 \sup_{\widehat S}|Z_{k+1}|\le C A(1+u).
\]
For $z,z'$ with $d=|z-z'|\le\ell$, choose $j$ so that
$2^{-j}\ell\asymp d$. Their level-$j$ representatives are
$O(2^{-j}\ell)$ apart, hence their difference is one of the
tested pairs; the two finer telescopes cost at most
$CA(1+u)2^{-\gamma j}$. Consequently,
\[
 |Z_{k+1}(z)-Z_{k+1}(z')|
 \le CA(1+u)(d/\ell)^\gamma
 \le CA(1+u)(d/\ell)^\beta.
\]
For $d>\ell$ the supremum bound gives the same
$C^\beta$ estimate. Taking suprema over a fixed countable
dense set and extending to the continuous version makes
$E_{k+1,S}$ measurable. We conclude from
\eqref{eq:chainingfailure} that
$\Pp_k(E_{k+1,S}>C A(1+u))\le C'e^{-c'u^2}$,
which, after adjusting constants, proves the first inequality
in \eqref{eq:envelopetail}. Tail integration gives, for
$c$ smaller than the tail exponent,
\[
 \E_k e^{cE_{k+1,S}^2}
 =1+\int_0^\infty 2ct e^{ct^2}
       \Pp_k(E_{k+1,S}>t)\,dt\le C.
\]

Finally, the area floor implies that at each generation there
are at most $(c_*s^2)^{-1}$ children, including terminal unions.
Because the next subdivision is $\F_k$-measurable, the
conditional bound in \eqref{eq:envelopetail} can be summed over
all children selected for the next actual reveal, and
then averaged over the past. A further union bound over the
$J_0$ generations gives the first inequality of
\eqref{eq:globalcut}. The envelopes on distinct children need not
be independent for this argument. Since $s=2^{-m}m^3$ and $J_0\le C_{\rm dep}m$,
\[
 C J_0s^{-2}e^{-cR^2m}
 \le C' m^{-5}\exp\bigl[-(cR^2-2\log2)m\bigr].
\]
Choose fixed $R$ with $cR^2>2\log2$. This proves the second
inequality in \eqref{eq:globalcut}.
\end{proof}

\paragraph{Local estimates used in the induction.}
Conditionally on its old history, every active parent satisfies the
following, by Lemmas~\ref{lem:whitney}, \ref{lem:pk},
\ref{lem:piolaGaussian} and~\ref{lem:envelope}. The chart and the
entire child partition are predictable, and every child has area at
least $c_*s^2$. Active children contract in birth charts, while
terminal children have intrinsic diameter $O(s)$; the contraction is
quantified in Lemma~\ref{lem:birthchart}, and its use along an entire
ancestry is recorded in Lemma~\ref{lem:actualphase}. The increment
$V_{k+1}=\sqrt\delta A_QZ_{k+1}$ has zero normal trace on its parent
and divergence $\sqrt\delta\,\Delta f_k$. For $x=b_Q+A_Qz$ in a
child $S$,
\[
 \operatorname{Cov}_k V_{k+1}(x)\preceq C\delta G_Q^+,
 \qquad
 |V_{k+1,i}(x)|\le
 \sqrt\delta\sqrt{(G_Q^+)_{ii}}\,E_{k+1,S},
 \qquad \E_k e^{cE_{k+1,S}^2}\le C,
\]
and the same envelope controls the intrinsic $C^\beta$ seminorm.
The constants are independent of the physical eccentricity, whose
effect is confined to the factors involving $G_Q^+$.

\paragraph{Lag and memory.}
For each $x$, let $E_j(x)$ be the envelope associated to its stage-$j$
child, or zero on an inactive generation. Fix a deterministic memory
discount $0<\rho_1<1$; it is chosen in
Section~\ref{subsec:parameters} so that the square of the geometric
ancestor-decay factor ($r_{\rm osc}$ in Lemma~\ref{lem:actualphase})
is strictly smaller than $\rho_1$. The subscript in $\rho_1$ is part
of its name, not a generation index. Set
\begin{equation}\label{eq:memory}
 T_0=0,\quad T_{k+1}=\rho_1T_k+E_{k+1}^2,\quad
 L_k(x)=L(U_k(x),G_k(x)),\quad S_k=L_k^2+T_k.
\end{equation}
Both terms of $S_k$ are nonnegative. The memory $T_k$ records old
oscillations with an exponentially discounted weight, and $L_k^2$
measures how far the current routing metric lags behind the
instantaneous Bellman metric. The recurrence for $T_k$ applies only
at actual reveals: on a frozen leaf, $T_k$, $L_k$ and $S_k$ remain
unchanged, so padding the filtration does not discount the past.

\paragraph{Stopping rules.}
Fix the deterministic modular stop threshold $T_*>0$. The construction
uses exactly two stopping rules, both checked before the next
subdivision is selected. If some previously attempted child envelope
exceeded $R\sqrt m$, the history has been killed. If the old modular
\phantomsection\label{not:modularfirst}
$\mathcal Z_k:=\int_{\T^2}e^{\Phi_k}$ (see \eqref{eq:Phi}) satisfies
$\mathcal Z_k\ge T_*$, the history has overshot. Otherwise the next
subdivision and metric update are chosen, the new child masses are
revealed, and their envelopes are computed. If a new envelope
exceeds $R\sqrt m$, we kill the scalar stopped supermartingale. If an overshoot
occurs, we retain the attained modular value and freeze further
geometric refinement. The diagonal caps and the relative phase
comparisons are not tested by the construction. On every
history that has been neither killed nor stopped by overshoot they
are consequences of Proposition~\ref{prop:bootstrap}. This separates
a random failure event from the field estimates valid on successful
histories.

\paragraph{Birth and reveal convention.}
A stage-$j$ square $Q_j$ is born with metric $G_j$, which is also the
physical covariance metric of the increment $V_j=\sqrt\delta\,G_j^{1/2}Z_j$.
On an active parent $Q_j$ the predictable next metric is $G_j^+$,
and a newly created child is born with $G_{j+1}=G_j^+$. A child created
as a clipped terminal union receives that increment once and then
freezes. A square stopped directly at the threshold is not updated
and reveals no further Gaussian increment. At padded times, the
exponent of the ancestor contraction is the number of genuine
reveals, not the calendar difference. We use essential suprema across
cell boundaries and pointwise continuous versions on each cell
interior.


\subsection{Auxiliary estimates}
\phantomsection\label{not:relativephase}
Throughout the induction, $0<\eta<1$ is the fixed relative
phase tolerance per actual Gaussian reveal:
$(1-\eta)H(U_{j-1})\le H(U_j)\le(1+\eta)H(U_{j-1})$.
It is chosen small before the Gaussian step size $\varepsilon$.
There is no new phase comparison on a frozen stage.

The two lemmas below are the deterministic inputs of the induction.
Lemma~\ref{lem:actualphase} bounds the oscillation of $U_b$ on a cell
by the memory $T_b$, summing only over actual reveals and not over
padded calendar times. Lemma~\ref{lem:lagsegments} identifies short
anisotropic segments along which the lag may be differentiated using
only the old diagonal cap. Neither follows from the local ray-kernel
estimate alone.

\begin{lemma}\label{lem:actualphase}
Let $Q_0\supset\cdots\supset Q_b$ be the actual ancestral chain of a
point, with exactly one new Gaussian increment $V_j$ revealed when
$Q_j$ is born, $1\le j\le b$.  The birth metric $G_j$ and intrinsic
length $\ell_j$ belong to $Q_j$, and $E_j(Q_j)$ is the intrinsic $C^\beta$ envelope of the
normalized field $Z_j$, before the Piola factor $G_j^{1/2}$. Suppose the already certified caps and
phase inequalities hold at each reveal:
\[
 (G_j)_{ii}\le M H(U_j)^{2a_i},\qquad
 (1-\eta)H(U_{j-1})\le H(U_j)\le(1+\eta)H(U_{j-1}).
\]
If the actual child contractions and metric comparisons of
Lemma~\ref{lem:ancestorindexed} hold, define the effective
ancestor-oscillation ratio $r_{\rm osc}$ by
\[
 r_{\rm osc}=(\vartheta\sqrt K)^\beta(1-\eta)^{-a_*},
 \qquad r_{\rm osc}^2<\rho_1<1.
\]
Then for $x,z$ in the same stage-$b$ cell (whether this cell is
active, newly terminal, or directly stopped),
\begin{equation}\label{eq:actualphasetransfer}
 |U_{b,i}(x)-U_{b,i}(z)|
 \le \frac{C\sqrt{\delta M}}{\sqrt{1-r_{\rm osc}^2/\rho_1}}
 H(U_b(x))^{a_i}\sqrt{T_b(x)},
 \qquad T_b(x)=\sum_{j=1}^{b}\rho_1^{b-j}E_j(Q_j(x))^2.
\end{equation}
The sum is over actual revelations only.  After a cell becomes terminal,
$U$, $G$ and $T$ are frozen there and no extra contraction factor is
generated by padding the filtration to the deterministic horizon.
\end{lemma}
\begin{proof}
We bound each revealed increment separately and sum.
For $1\le j\le b$, the pair $x,z$ lies in the actual stage-$j$
child on which $V_j=\sqrt\delta G_j^{1/2}Z_j$ was constructed.
The row bound for the Piola transform and the intrinsic envelope give
\[
 |V_{j,i}(x)-V_{j,i}(z)|
 \le\sqrt\delta\sqrt{(G_j)_{ii}}\,E_j(Q_j)
 \left(\frac{d_{G_j}(x,z)}{\ell_j}\right)^\beta.
\]
The $j$-th cap and the inequalities along the actual chain imply
\[
 \sqrt{(G_j)_{ii}}
 \le \sqrt M(1-\eta)^{-a_i(b-j)}H(U_b(x))^{a_i}.
\]
By Lemma~\ref{lem:ancestorindexed}, with a final terminal transition
counted only when it was actually revealed,
$d_{G_j}(x,z)/\ell_j\le C(\vartheta\sqrt K)^{b-j}$.
Since $U_0$ is spatially constant, summing the revealed increments
and using $a_i\le a_*$ yields
\[
 |U_{b,i}(x)-U_{b,i}(z)|\le
 C\sqrt{\delta M}H(U_b(x))^{a_i}
 \sum_{j=1}^{b}r_{\rm osc}^{b-j}E_j(Q_j(x)).
\]
The weighted Cauchy--Schwarz inequality, applied over actual
reveals only, gives
\[
 \sum_{j=1}^{b}r_{\rm osc}^{b-j}E_j
 \le\left(\sum_{n\ge0}(r_{\rm osc}^2/\rho_1)^n\right)^{1/2}
 \left(\sum_{j=1}^{b}\rho_1^{b-j}E_j^2\right)^{1/2}.
\]
For a directly stopped cell $Q_b$ there is no subsequent increment;
for a newly terminal union the $b$-th increment is included once.
After termination the field does not change and the inequality
is retained without further discounting.
\end{proof}

\begin{lemma}\label{lem:lagsegments}
Fix the finite critical exponents and $K,M<\infty$.  There are
$b_0>0$ and $C=C(K,M,\boldsymbol r)$ such that the following holds.
Let $y\in\R^2$, let $G$ be positive definite with $\det G=1$ and
$G_{ii}\le M H(y)^{2a_i}$, and let $\widetilde G$ satisfy
$0\prec\widetilde G\preceq K G$ and $\det\widetilde G=1$.  Suppose $y^{(1)},y^{(2)}$
belong to the anisotropic rectangle
\[
 \mathcal R_b(y)=\{w:|w_i-y_i|\le bH(y)^{a_i},\ i=1,2\},
 \qquad 0<b\le b_0.
\]
At every $w$ in the rectangle, and hence along every segment
joining any two of its points,
\begin{equation}\label{eq:segmentcaps}
 \tfrac12H(y)\le H(w)\le 2H(y),\quad
 G_{ii}\le C_M H(w)^{2a_i},\quad
 \widetilde G_{ii}\le C_{K,M}H(w)^{2a_i}.
\end{equation}
Consequently, for $G'\in\{G,\widetilde G\}$,
\begin{equation}\label{eq:segmentlag}
 |L(y^{(1)},G')-L(y^{(2)},G')|
 \le C_{K,M}\sqrt{H(y)}\sum_{i=1}^2
 \frac{|y_i^{(1)}-y_i^{(2)}|}{H(y)^{a_i}}.
\end{equation}
\end{lemma}
\begin{proof}
For $F_i(t)=(1+t^2)^{r_i/2}$ one has
$|F_i'(t)|\le r_i(1+t^2)^{(r_i-1)/2}$. For
$w\in\mathcal R_1(y)$, the bound $(1+y_i^2)^{1/2}\le H(y)^{a_i}$
gives $(1+w_i^2)^{1/2}\le 2H(y)^{a_i}$, hence
$|\partial_iH(w)|\le C H(y)^{1-a_i}$.  Integrating along coordinate segments,
$|H(w)-H(y)|\le CbH(y)$ on $\mathcal R_b(y)$.
Taking $b_0$ sufficiently small gives the first comparison in
\eqref{eq:segmentcaps}; the remaining two follow directly from
the old diagonal cap and $\widetilde G\preceq KG$.
Lemma~\ref{lem:selector} applies to $G'$ at each point of the
segment, with the enlarged fixed cap constant $C_{K,M}$, and gives
$|\partial_iL(w,G')|\le C_{K,M}H(w)^{1/2-a_i}$.
Integrating this derivative along the straight segment and using
$H(w)\asymp H(y)$ proves \eqref{eq:segmentlag}.
\end{proof}

The proof uses only old-state information and the rectangle bound,
not any contraction of the lag. Hence
Lemma~\ref{lem:lagsegments} applies to the segment from the old
representative to a point and to the segment of a new increment,
without assuming the new-state cap in advance.

\subsection{The induction}
The induction is carried out on the following events. For $k\ge0$
let $\mathcal H_k\in\F_k$ be the old good event at
stage $k$: the history has been neither killed nor stopped by
overshoot before stage $k$, and
\begin{enumerate}[label=(H\arabic*),leftmargin=3em]
\item for every $j\le k$ and almost every $x$ in an active or terminal
  cell, $(G_j)_{ii}(x)\le M H(U_j(x))^{2a_i}$ for $i=1,2$, and, if a
  Gaussian increment was actually revealed at $x$ at stage $j\ge1$,
  $(1-\eta)H(U_{j-1}(x))\le H(U_j(x))\le(1+\eta)H(U_{j-1}(x))$;
\item every envelope revealed up to stage $k$ satisfies
  $E_{j,S}\le R\sqrt m$;
\item $\mathcal Z_k<T_*$.
\end{enumerate}
Let $\mathcal G_{k+1}\in\F_{k+1}$ be the event that every child
revealed at stage $k+1$ has $E_{k+1,S}\le R\sqrt m$.
Part~(i) of the next proposition concerns $\F_k$-measurable
quantities only; part~(ii) adds the fresh increment.

\begin{proposition}\label{prop:bootstrap}
Fix $K>1$ and choose the constants in the order of
Section~\ref{subsec:parameters}. Then the following holds for every
$k\ge0$.
\begin{enumerate}[label=(\roman*),leftmargin=2.7em]
\item On $\mathcal H_k$, before the stage-$(k+1)$ revelation, every
  active parent $Q$ and almost every $x\in Q$ satisfy
  \begin{equation}\label{eq:Hcap}
   H(U_k(x))\le C_{T_*}m,
  \end{equation}
  and the predictable metric $G^+=G_Q^+$ satisfies
  \eqref{eq:margin} below.
\item On $\mathcal H_k\cap\mathcal G_{k+1}$, simultaneously for all
  new children and almost every point (and for the chosen continuous
  representatives inside each child),
  \begin{align}\label{eq:threecaps}
    (G_{k+1})_{ii}(x)&\le M H(U_{k+1}(x))^{2a_i},\\
    (1-\eta)H(U_k(x))&\le H(U_{k+1}(x))
                      \le(1+\eta)H(U_k(x)),\nonumber\\
    S_{k+1}(x)&\le q_0S_k(x)+C(1+E_{k+1}(x)^2).\nonumber
  \end{align}
  On an inactive cell, $U$, $G$ and $S$ are unchanged, so the first
  two bounds persist and $S_{k+1}=S_k$.
\end{enumerate}
In particular,
$\mathcal H_k\cap\mathcal G_{k+1}\cap\{\mathcal Z_{k+1}<T_*\}
\subset\mathcal H_{k+1}$.
\end{proposition}
\begin{proof}
Fix $k\ge0$ and work on $\mathcal H_k$, simultaneously on the whole
predictable tree. Steps~1, 2 and~4 involve only $\F_k$-measurable
quantities and prove part~(i); Steps~3 and~5 add the fresh increment
and prove part~(ii).

\textbf{Step 1: the old oscillation estimate.}
For $x$ in an active cell of stage $k$, the
caps and phase comparisons in (H1) are exactly the hypotheses of
Lemma~\ref{lem:actualphase} along the actual ancestral chain of $x$,
with $b=k$; the child contractions of
Lemma~\ref{lem:ancestorindexed} hold by construction. Since
$\delta=\varepsilon^2/m$ and, by (H2),
$T_k\le R^2m/(1-\rho_1)$, the estimate
\eqref{eq:actualphasetransfer} gives, for $x,z$ in the same current
active cell,
\begin{equation}\label{eq:osc}
 |U_{k,i}(x)-U_{k,i}(z)|
 \le C\sqrt{\delta M}\,H(U_k(x))^{a_i}\sqrt{T_k(x)}
 \le C\varepsilon R\sqrt M\,H(U_k(x))^{a_i}.
\end{equation}
The same estimate holds on a new terminal child, by the terminal case
of Lemma~\ref{lem:ancestorindexed}. Since terminal children are
frozen, it is never invoked for a later refinement of such a child.

\textbf{Step 2: the old representative and a strict cap margin.}
The derivative bounds $|\partial_iH(y)|\le CH(y)^{1-a_i}$ imply that
if $|z_i-y_i|\le bH(y)^{a_i}$ with $b$ small, then
$H(z)=(1+O(b))H(y)$. Thus \eqref{eq:osc} makes
$H(U_k(x_Q))$ and $H(U_k(x))$ relatively close.
The clipping inequality \eqref{eq:clipping} gives
\[
 (G^+)_{ii}\le K^{-1}(G_k)_{ii}+C(U_k(x_Q))_{ii}
 \le\bigl(K^{-1}M+C(1+O(\varepsilon R\sqrt M))\bigr)
            H(U_k(x))^{2a_i}.
\]
\phantomsection\label{not:qd}
Choose the strict diagonal-cap contraction factor $q_d\in(K^{-1},1)$,
then take $M$ large enough that
$K^{-1}M+C\le q_dM$ with slack, and finally take $\varepsilon$
small enough to absorb the representative error. Thus
\begin{equation}\label{eq:margin}
 (G^+)_{ii}\le q_dM H(U_k(x))^{2a_i}.
\end{equation}

\textbf{Step 3: the fresh increment and the relative phase.}
On the new intrinsic envelope cutoff, the physical row estimate
$|e_i^TA_QZ|\le\sqrt{(G_Q^+)_{ii}}\,E_{k+1}$ and
\eqref{eq:margin} yield
\[
 |V_{k+1,i}(x)|\le
 C\sqrt{\delta M}\,H(U_k(x))^{a_i}E_{k+1}(x)
 \le C\varepsilon R\sqrt M\,H(U_k(x))^{a_i}.
\]
\phantomsection\label{not:xi}
Consequently $|H(U_{k+1})-H(U_k)|\le\xi H(U_k)$,
where $\xi=C\varepsilon R\sqrt M$ can be made smaller than
any prescribed fixed $\eta>0$. Moreover, with
$q_d<(1-\eta)^{2a_*}$ (Section~\ref{subsec:parameters}),
\eqref{eq:margin} and this phase comparison preserve the diagonal cap
$(G_{k+1})_{ii}\le M H(U_{k+1})^{2a_i}$.
This proves the first two lines of \eqref{eq:threecaps}
for all new children simultaneously.

\textbf{Step 4: the old modular gives a pointwise phase cap.}
On an active parent $Q$, the oscillation bound \eqref{eq:osc} of
Step~1 makes $H(U_k)$ comparable throughout $Q$.
Every active $Q$ has area at least $cs^2$. By (H3),
$\mathcal Z_k\le T_*$, and $\lambda_k\ge\lambda_{J_0}\ge1/2$ by the
choice of $\varepsilon$ after $\Lambda$ and $C_{\rm dep}$
(Section~\ref{subsec:parameters}). Integration over $Q$ gives
\[
 T_*\ge\int_Qe^{\lambda_kH(U_k)}
 \ge c|Q|e^{cH(U_k(x))},
\]
so
$H(U_k(x))\le C\bigl(1+\log(T_*/s^2)\bigr)\le C_{T_*}m$.
This proves \eqref{eq:Hcap}, which is $\F_k$-measurable.

\textbf{Step 5: contraction of the lag and memory.}
Set $y=U_k(x)$, $y_Q=U_k(x_Q)$ and
$y'=U_{k+1}(x)$.  Steps~1 and~3 establish
$|y_{Q,i}-y_i|+|y'_i-y_i|\le
bH(y)^{a_i}$ with $b\le C\varepsilon R\sqrt M$.
Both metrics $G_k$ and $G^+$ satisfy the old diagonal cap,
with $G^+\preceq KG_k$ and $\det G^+=1$. Once $\varepsilon$ is
chosen so that $b\le b_0$, Lemma~\ref{lem:lagsegments} applies with
$y^{(1)}=y_Q$ and $y^{(2)}=y'$. It gives the cap throughout
the segments $[y,y_Q]$ and $[y,y']$, which permits integration of
\eqref{eq:lagder} along them; no conclusion of the lag
update is used. At the representative the clipping contracts the
lag, $L(y_Q,G^+)\le K^{-1/2}L(y_Q,G_k)+C_K$ by \eqref{eq:clipping}.
Together with the segment estimates, this gives
\[
 L_{k+1}(x)
 \le K^{-1/2}L_k(x)+C_K
    +C_M\sqrt{\delta H(U_k(x))}
             \bigl(\sqrt{T_k(x)}+E_{k+1}(x)\bigr).
\]
By Step~4, $\sqrt{\delta H}\le C_{T_*}\varepsilon$.
\phantomsection\label{not:qL}
Fix the lag-contraction factor $q_L>K^{-1}$ and a Young slack
$\tau>0$ with $(1+\tau)K^{-1}\le q_L$. Squaring with this slack gives
\[
 L_{k+1}^2\le q_LL_k^2+C_\tau+C_\tau\varepsilon^2T_k
                  +C_\tau\varepsilon^2E_{k+1}^2.
\]
Since $T_{k+1}=\rho_1T_k+E_{k+1}^2$, choose
$\max\{q_L,\rho_1\}<q_0<1$ (the combined energy contraction)
and shrink $\varepsilon$
further, obtaining
\[
 S_{k+1}\le q_0S_k+C(1+E_{k+1}^2).
\]
This proves part~(ii) simultaneously on all active children.
On inactive branches nothing is renewed: the old bounds remain
valid, and the clock is the only quantity that changes.

Finally, (H1) at stage $k+1$ is \eqref{eq:threecaps} together with
the persistence on inactive cells, (H2) at stage $k+1$ is
$\mathcal G_{k+1}$, and (H3) is the remaining condition
$\mathcal Z_{k+1}<T_*$. All constants in Steps~1--5 are the fixed
constants of Section~\ref{subsec:parameters}; none depends on the
parent.
\end{proof}

Since $U_0=0$, $G_0=I$, $H(0)=2$ and $M\ge1$, the event
$\mathcal H_0$ is the event $\{\mathcal Z_0<T_*\}$, which is the
whole space by the choice of $T_*$ in Section~\ref{sec:phase}.
By induction on $k$, using the last inclusion of
Proposition~\ref{prop:bootstrap}, $\mathcal H_k$ contains every
history that has been neither killed nor stopped by overshoot before
stage $k$.


\subsection{Choice of parameters}\label{subsec:parameters}
All constants of Part~\ref{umc:partI} are fixed, with strictly positive margins, in
the order of Table~\ref{tab:parameters}. Each constant depends only on
the fixed exponents $\boldsymbol r$ and on the constants listed above
it. None of them depends on the grid scale $m$ or on the freshly
revealed Gaussian data, and no constant is constrained by a constant
chosen after it.

\begin{table}[ht]
\centering
\small
\renewcommand{\arraystretch}{1.15}
\begin{tabular}{@{}p{0.25\textwidth}p{0.69\textwidth}@{}}
\toprule
Constant & Requirement and source\\
\midrule
$K>1$, $0<\beta<\alpha<1$
 & Metric distortion and H\"older exponents; arbitrary.\\
$c_0$; $A$, $\gamma$
 & Inball ratio of intrinsic parents, depending on $K$; then the
   Whitney displacement and clearance constants, depending only on
   $c_0$ (Lemma~\ref{lem:whitney}).\\
$q_d,\,q_L\in(K^{-1},1)$; $\tau$
 & Cap and lag contraction factors; Young slack $\tau>0$ with
   $(1+\tau)K^{-1}\le q_L$ (Step~5 of
   Proposition~\ref{prop:bootstrap}).\\
$\eta$, $q_0$
 & $q_d<(1-\eta)^{2a_*}$, and $q_0$ with $q_L<q_0<1/(1+\eta)$.\\
$\vartheta$; $\theta$, $D_0$; $C_{\rm dep}$
 & $\vartheta\sqrt K<1$ and
   $r_{\rm osc}:=(\vartheta\sqrt K)^\beta(1-\eta)^{-a_*}$ with
   $r_{\rm osc}^2<q_0$; implemented by the Whitney cap $\theta$ and
   the stopping threshold $D_0$ (Lemma~\ref{lem:birthchart}); this
   fixes the depth constant $C_{\rm dep}$.\\
$\rho_1$
 & $r_{\rm osc}^2<\rho_1<q_0$.\\
$M$
 & $K^{-1}M+C\le q_dM$ with slack, where $C$ comes from
   \eqref{eq:selectors} and \eqref{eq:clipping} (Step~2).\\
$\kappa$, $\Lambda$
 & Drift balance in the one-step contraction
   (Section~\ref{sec:phase}).\\
$T_*$
 & $T_*>\mathcal Z_0$, large enough for a positive success
   probability (Proposition~\ref{prop:stopped}).\\
$R$
 & $cR^2>2\log2$, with $c$ the envelope tail constant
   (Lemma~\ref{lem:envelope}).\\
$\varepsilon$
 & Last: $b\le b_0$ in Lemma~\ref{lem:lagsegments}; $\xi<\eta$;
   absorption of the representative error in \eqref{eq:margin} and of
   $\varepsilon^2T_k$ in Step~5; $\lambda_{J_0}\ge1/2$, i.e.\
   $\Lambda C_{\rm dep}\varepsilon^2\le\log2$; the smallness conditions
   of Section~\ref{sec:phase}; the terminal corrections of
   Section~\ref{sec:terminal}.\\
\bottomrule
\end{tabular}
\caption{Order of the parameter choices in Part~\ref{umc:partI}.}
\label{tab:parameters}
\end{table}

\noindent In particular,
\begin{equation}\label{eq:strictschedule}
 r_{\rm osc}^2<\rho_1<q_0<1,\qquad q_L<q_0,\qquad
 q_d<(1-\eta)^{2a_*},\qquad q_0(1+\eta)<1.
\end{equation}

\section{Exponential moment estimates}\label{sec:phase}

This section turns the deterministic bounds of
Proposition~\ref{prop:bootstrap} into a probabilistic estimate on the
exponential modular of $U_k$. Proposition~\ref{prop:mgf} shows that,
at a single point and conditionally on the past, one refinement step
does not increase $e^{\Phi_k}$ in expectation on the local envelope
event. Proposition~\ref{prop:stopped} integrates this in space, kills
the process at envelope failures and stops it at overshoot, and
deduces from Ville's inequality that with probability bounded below
all leaves terminate with the Orlicz bound \eqref{eq:leafbound}.

Define the positive phase and its spatial modular by
\begin{equation}\label{eq:Phi}
 \Phi_k(x)=\lambda_kH(U_k(x))\bigl(1+\kappa\delta S_k(x)\bigr),
 \qquad \mathcal Z_k=\int_{\T^2}e^{\Phi_k(x)}\dd x.
\end{equation}
At the critical exponents, the variance associated with $\lambda_kH$
has the same order as the phase itself. The additional term
$\kappa\delta HS$ has a negative drift under active refinement,
which compensates for this covariance cost. The mechanism is that,
after a Taylor expansion, the fresh Gaussian field enters the
increment of $\Phi_k$ only through a linear term $\ell\cdot Z_{k+1}$
whose coefficient $\ell$ is old-measurable
(Lemma~\ref{lem:path}). Its exponential moment is then computed
exactly, and the envelope contributes only a quadratic correction
(Lemma~\ref{lem:joint}).

\begin{lemma}\label{lem:path}
On an old good active cell, put
$y=U_k(x)$, $H=H(y)$, $H'=H(U_{k+1}(x))$,
$S=S_k(x)$, $S'=S_{k+1}(x)$, $E=E_{k+1}(x)$, and
$\lambda=\lambda_{k+1}$. On the local envelope cutoff,
\begin{equation}\label{eq:path}
 \Phi_{k+1}-\Phi_k
 \le-\Lambda\lambda\delta H
       -\kappa d_0\lambda\delta HS
       +\ell\cdot Z_{k+1}
       +A_\kappa\lambda\delta H(1+E^2),
\end{equation}
where
\[
 d_0=1-q_0(1+\xi)>0,
 \qquad \ell=\lambda\sqrt\delta\,A_Q^Tv(y)
       =\lambda\sqrt\delta\,(G_Q^+)^{1/2}v(y).
\]
In particular, $\ell$ is measurable \emph{before} the fresh Gaussian
variables are revealed.
\end{lemma}
\begin{proof}
Taylor's formula and $F_i''(t)\lesssim H^{1-2a_i}$ along the
short segment give
\begin{equation}\label{eq:taylor}
 H'-H\le v(y)\cdot V_{k+1}+A\delta H E^2,
 \qquad |H'-H|\le\xi H.
\end{equation}
By the memory inequality $S'\le q_0S+A(1+E^2)$ and positivity,
\[
 H'S'-HS
 \le (1+\xi)H\bigl(q_0S+A(1+E^2)\bigr)-HS
 =-d_0HS+A(1+\xi)H(1+E^2).
\]
Choose $\xi$ so that $q_0(1+\xi)<1$. Now use the exact
algebraic identity
\[
 \Phi_{k+1}-\Phi_k
 =-(\lambda_k-\lambda)H(1+\kappa\delta S)
   +\lambda(H'-H)+\kappa\lambda\delta(H'S'-HS).
\]
The clock gives $\lambda_k-\lambda\ge\Lambda\lambda\delta$.
Insert the two preceding bounds and
$V_{k+1}=\sqrt\delta A_QZ_{k+1}$.
This proves \eqref{eq:path}.
\end{proof}

\begin{lemma}\label{lem:joint}
Suppose, conditionally on $\F_k$, $Z$ is a centered Gaussian
vector with covariance at most $\sigma^2I$ and an accompanying
(not necessarily independent) envelope satisfies $\E_ke^{cE^2}\le C$.
For $0\le u\le u_0<c/4$ and any old-measurable $\ell\in\R^2$,
\begin{equation}\label{eq:joint}
 \log\E_ke^{\ell\cdot Z+uE^2}
       \le \sigma^2|\ell|^2+C_*u.
\end{equation}
\end{lemma}
\begin{proof}
Since $E$ and $Z$ may be dependent, we separate them by
Cauchy--Schwarz:
\[
 \E_ke^{\ell\cdot Z+uE^2}
 \le (\E_ke^{2\ell\cdot Z})^{1/2}
      (\E_ke^{2uE^2})^{1/2}.
\]
The first factor is evaluated exactly using Gaussianity.
For the second, $\E_k(E^2e^{2u_0E^2})$ is uniformly finite
by the assumed exponential-square moment. Integrating the
$u$-derivative of its logarithm from $0$ to $u$ gives $C_*u$.
\end{proof}

\begin{proposition}\label{prop:mgf}
For $\kappa$ sufficiently large, $\Lambda$ sufficiently large
after $\kappa$, and $\varepsilon$ sufficiently small after all fixed
constants, on any old good active history,
\begin{equation}\label{eq:onestep}
 \E_k\left[\one_{\{E_{k+1}(x)\le R\sqrt m\}}
          e^{\Phi_{k+1}(x)}\right]
 \le e^{\Phi_k(x)}.
\end{equation}
\end{proposition}
\begin{proof}
We exponentiate \eqref{eq:path} and apply Lemma~\ref{lem:joint}.
First, from $vv^T\preceq CH C(y)^{-1}$ and $G^+\preceq KG_k$,
\[
 |\ell|^2=\lambda^2\delta v^TA_QA_Q^Tv
 =\lambda^2\delta v^TG_Q^+v
 \le C_K\lambda\delta H\tr(G_kC(y)^{-1})
 \le C_K\lambda\delta HS.
\]
This is the only place where the lag enters.
Set $u=A_\kappa\lambda\delta H$. By the old phase cap
$H\le C_{T_*}m$, we have $u\le C\varepsilon^2$,
so Lemma~\ref{lem:joint} applies.
Exponentiating \eqref{eq:path} on the local good event,
then dropping its indicator in the upper bound, gives
\[
 \log\E_k\bigl[
  \one_{\{E\le R\sqrt m\}}e^{\Phi_{k+1}-\Phi_k}\bigr]
 \le -\Lambda\lambda\delta H
       -\kappa d_0\lambda\delta HS
       +C_K\lambda\delta HS+C_\kappa\lambda\delta H.
\]
Choosing $\kappa d_0>C_K$ and then $\Lambda>C_\kappa$ makes the
right-hand side nonpositive.
\end{proof}

Because the tilt $\ell$ is predictable and the fresh field enters
linearly, the estimate involves no $\delta H^2$ term and no fourth
cumulant.

\subsection{Stopping rule}
The one-step estimate is applied up to a stopping time. We define
the numerical process on all outcomes of each attempted reveal.

Let $J_0$ be a fixed upper bound on active depth, and let $\F_k$
be the completed sigma-field generated by the geometry and all
Gaussian cell masses revealed through stage $k$; as described in
Section~\ref{sec:bootstrap}, the entire next subdivision is
$\F_k$-measurable. We pad terminal branches to the deterministic
horizon $J_0$ by doing nothing. The old good event is
$A_k:=\mathcal H_k\in\F_k$, defined before
Proposition~\ref{prop:bootstrap}: previous envelopes were acceptable
and the last attained modular is strictly below $T_*$; the caps and
relative-phase bounds then hold by that proposition. At stage $k+1$
write
\[
 B_{k+1}=\mathcal G_{k+1}
 =\{\text{every newly revealed child envelope is }\le R\sqrt m\}
 \in\F_{k+1},
\]
where no new child is assigned an envelope on a frozen branch.

Define a nonnegative adapted scalar process $\widehat{\mathcal Z}_k$,
the killed-and-frozen modular, as follows:
\begin{enumerate}[label=(\roman*),leftmargin=2.7em]
\item Initially $\widehat{\mathcal Z}_0=\mathcal Z_0$.
\item Before failure or overshoot, attempt the next geometric
  refinement. If $B_{k+1}$ fails, put
  $\widehat{\mathcal Z}_{k+1}=0$ and thereafter keep the
  numerical process zero.
\item If $B_{k+1}$ holds, use the actual new modular
  $\widehat{\mathcal Z}_{k+1}=\mathcal Z_{k+1}$. If it is at
  least $T_*$, freeze the geometric field and its memory for all
  later stages, allowing only $\lambda$ to decrease.
\item On a branch that became terminal for geometric reasons,
  freeze $(U,S)$ on that branch from then on, while the clock
  continues to decrease.
\end{enumerate}

\begin{proposition}
\label{prop:stopped}
The process $\widehat{\mathcal Z}_k$ is a nonnegative supermartingale
and
\begin{equation}\label{eq:overshoot}
 \Pp\bigl(\exists j\le J_0:\,\widehat{\mathcal Z}_j\ge T_*\bigr)
 \le \frac{\mathcal Z_0}{T_*}.
\end{equation}
The global envelope-failure probability is $o(1)$.
Consequently, for large fixed $T_*$, there is a constant
$p_*>0$ independent of $m$ such that with probability at least
$p_*$ all leaves terminate and the leaf current satisfies
\begin{equation}\label{eq:leafbound}
 \divv J_{\mathrm{leaf}}=P_{\pr_*}W,
 \qquad \max_i\norm{J_{\mathrm{leaf},i}}_{\psi_{r_i}}
 \le C\delta^{-1/2}\le C_{\boldsymbol r}\sqrt m.
\end{equation}
\end{proposition}

\begin{proof}
We first prove the supermartingale property, case by case.
Suppose the past has neither failed nor overshot. For each active
$x$, Propositions~\ref{prop:bootstrap} and~\ref{prop:mgf} give
\[
 \E_k\left[\one_{B_{k+1}}e^{\Phi_{k+1}(x)}\right]
 \le\E_k\left[\one_{\{E_{k+1}(x)\le R\sqrt m\}}
                    e^{\Phi_{k+1}(x)}\right]
 \le e^{\Phi_k(x)}.
\]
The first inequality holds because the global envelope event is
contained in the local one. On an inactive branch the
field and memory are frozen and the clock decreases, so the
integrand can only decrease and the same inequality holds.
Integrating over $x$ and using conditional Tonelli, we obtain
\[
 \E_k[\one_{B_{k+1}}\mathcal Z_{k+1}]\le\mathcal Z_k.
\]
Since the definition sets $\widehat{\mathcal Z}_{k+1}=0$ on
$B_{k+1}^c$, the left side is exactly the conditional expectation of
the numerical process at the new stage, until the first overshoot.
After the overshoot there is no further geometric evolution, and
$\lambda_{k+1}\le\lambda_k$ ensures
$\widehat{\mathcal Z}_{k+1}\le
\widehat{\mathcal Z}_k$ pathwise. After a failure the process is
zero. No conditioning on the success of a future reveal is involved.

The initial values are deterministic:
$U_0=0$, $H(U_0)=2$, $G_0=I$ and $T_0=0$; hence
$\mathcal Z_0=\exp(2(1+\kappa\delta L(0,I)^2))=O(1)$.
For $\tau=\inf\{k:\widehat{\mathcal Z}_k\ge T_*\}$,
Ville's inequality yields \eqref{eq:overshoot} for the killed
process. Independently of this stopping argument,
Lemma~\ref{lem:envelope} bounds the probability of any actual
attempted envelope failure by $o(1)$, using conditional tails and the
deterministic bound on the number of attempted children. Write
$\mathcal E_m$ for the event of any attempted envelope
failure and $\mathcal O_m$ for the first overshoot along an
all-envelope-good history. The successful completion event satisfies
\[
 \mathbb P(A_m)\ge1-\mathbb P(\mathcal O_m)
                   -\mathbb P(\mathcal E_m)
 \ge1-\mathcal Z_0/T_*-o(1).
\]
Thus $T_*$ may be fixed sufficiently large so that this lower
bound exceeds some $p_*>0$.

On this good completion event the final positive modular is at
most $T_*$. Since $\lambda_{J_0}\ge1/2$ (after choosing
$\varepsilon$ small), $\Phi_{J_0}\ge\frac12 H(U_{J_0})$.
The fixed exponential modular bound implies the coordinate
Luxemburg bounds for $U_{J_0}$ after an absolute dilation, and
dividing by $\sqrt\delta$ proves \eqref{eq:leafbound}.
\end{proof}

The threshold in \eqref{eq:overshoot} is tested on the killed
process $\widehat{\mathcal Z}_k$, so Ville's bound concerns a first
overshoot along a history on which all envelopes are good. On a killed
history the process is not the modular of an admissible current. The actual-reveal geometry,
including direct stops, terminal unions and padded times, is recorded
in Lemma~\ref{lem:actualreveal}. The next section corrects the
adaptive current from $P_{\pr_*}W$ to the fixed-grid source $W_h$.
\section{The terminal Gaussian correction}\label{sec:terminal}

The adaptive field from Proposition~\ref{prop:stopped} solves the
cell-averaged source $PW$, where $P$ averages over the terminal leaves
of a random partition. The target source is $W_h$, obtained by averaging
over the fixed $h$-grid. In this section we construct a current with
divergence $W_h-PW$ and show that it is negligible in the Orlicz norms
under consideration; this completes the proof of Theorem~\ref{thm:gauss}.
The difference $W_h-PW$ splits exactly into two pieces. The first is
built from the revealed averages $PW$ and has zero mean on each fixed
$h$-square, so it is inverted square by square. The second is the
fixed-grid average of the unrevealed residual $W-PW$; conditionally on
the terminal history it is Gaussian, and we invert it by convolution
with the periodic Green kernel. On the torus this is a periodic
convolution, with no physical boundary data.

\subsection{The exact splitting}
On a successful history, every terminal leaf $S$ has area at
least $cs^2$ and diameter $O(s)$ in its birth metric $G_S$.
This includes both the clipped unions born at a split and the active
squares stopped without another split; by Lemma~\ref{lem:birthchart},
the latter have intrinsic diameter at most $C D_0s$.
Write $A_S=G_S^{1/2}$ and $S=b_S+A_S\widehat S$ in its chosen
lift. The diagonal cap holds in physical coordinate directions,
so at the final good stage $H(U)\le C_{T_*}m$ gives
\[
 \|A_S\|_{\rm op}^2\le\operatorname{tr}G_S
 \le M\sum_{i=1}^2H(U)^{2a_i}\le C_{M,T_*}m^{2a_*}.
\]
We now pass from the local intrinsic geometry to the fixed physical
flat-torus metric, in which
\begin{equation}\label{eq:Dstar}
 \diam_{\rm physical}(S)\le D_*:=Csm^{a_*},
 \qquad a_*<\frac12,\qquad D_*/h=Cm^{3+a_*}.
\end{equation}
This bound uses the cap: with $Cs$ in place of $D_*$ it would fail to
hold uniformly in the eccentricity.

\begin{lemma}
\label{lem:terminalhistory}
On the successful first-failure event, every terminal leaf $S$,
including a directly stopped square and a clipped union born at the
last reveal, satisfies
\[
 |S|\ge cs^2,\qquad
 \operatorname{diam}_{\rm physical}S\le Csm^{a_*}.
\]
The terminal partition $\pr_*$ and the success event belong to the
finite terminal revealed sigma field $\F_*$.  For every fixed
deterministic $g\in L^2(\T^2)$, conditional on $\F_*$,
\[
 \operatorname{Var}\bigl((W-P_{\pr_*}W)(g)\mid\F_*\bigr)
 =\sum_{S\in\pr_*}\int_S|g-g_S|^2.
\]
\end{lemma}
\begin{proof}
There are two kinds of terminal leaves. A newly terminal union has
$G_S$-intrinsic diameter at most $C_As$, and a square stopped without
another reveal has birth-metric diameter at most $\sqrt K D_0s$. In
both cases the receiver construction gives the area floor.

On successful completion the final modular is at most $T_*$ and
$\lambda_{J_0}\ge1/2$. Lemma~\ref{lem:actualphase}, applied at
the last actual reveal of $S$, gives throughout $S$ the same phase
comparability that holds on an active cell. Since $U$ is frozen
afterwards, the comparison remains valid at time $J_0$. Integrating
the final modular over $S$ then gives $H(U|_S)\le C_{T_*}m$, also
when $S$ was created at the final reveal. The newly inherited diagonal
cap, or the old cap for a directly stopped square, yields
$\|G_S^{1/2}\|_{\rm op}\le C m^{a_*}$. Multiplying by the intrinsic
diameter proves the physical bound.

The algorithm makes all geometric choices before the corresponding
fresh Gaussian child contrasts are revealed, and it is padded to a
finite deterministic horizon. Since stopping and the success event
depend only on revealed contrasts, they are $\F_*$-measurable. The
residual white noise in the orthogonal complement of the finally
revealed cell indicators is conditionally Gaussian. Testing it against
a deterministic $g$ and using the conditional orthogonal projection
gives the displayed variance.
\end{proof}

From now on all distances and derivatives are physical, not
intrinsic; in particular the periodic Green kernel is never pulled
back by an adaptive $A_S$. Write $P=P_{\pr_*}$. The identity
\begin{equation}\label{eq:decomp}
 W_h-PW=(P_h-I)PW+P_h(W-PW)=:g^{(1)}+g^{(2)}
\end{equation}
holds exactly, as an identity of centered distributions.

\subsection{Correction of the revealed part}
\phantomsection\label{not:terminalcurrents}
The first source $g^{(1)}$ has zero integral on each fixed $h$-square.
Applying the ray inverse on these squares gives a periodic current
$J^{(1)}$ with
\[
 \divv J^{(1)}=g^{(1)},\qquad
 \norm{J^{(1)}}_\infty\le Ch\norm{PW}_\infty.
\]
Lemma~\ref{lem:randommax} controls the random final cells directly and
gives, with probability $1-o(1)$,
\begin{equation}\label{eq:smallfirst}
 \norm{PW}_\infty\le C\sqrt m/s,\qquad
 \norm{J^{(1)}}_\infty\le C\sqrt m/m^3=o(1).
\end{equation}

\subsection{Periodic kernel estimates for the unrevealed bridge}
\phantomsection\label{not:periodicgreen}
Let $G_{\T}$ be the mean-zero Green function of the torus
Laplacian in the fixed physical coordinates, and choose
the periodic vector kernel
$K_{\T}(x,y)=\nabla_xG_{\T}(x-y)$ with the sign convention
\[
 \divv_xK_{\T}(x,y)=\delta_y-1.
\]
The standard local expansion of the periodic Green function
and smoothness away from the diagonal give
\begin{equation}\label{eq:periodicK}
 |K_{\T}(x,y)|\le C(1+d_{\T}(x,y)^{-1}),\qquad
 |\nabla_yK_{\T}(x,y)|\le C(1+d_{\T}(x,y)^{-2}).
\end{equation}
In the source variable let
$K_{x,h}=P_h^{(y)}K_{\T}(x,\cdot)$, a piecewise constant
function with jumps along the fixed $h$-grid, and define
\[
 J^{(2)}(x)=(W-PW)(K_{x,h}).
\]
Its divergence is $g^{(2)}$, with the vector notation understood
coordinatewise. Lemma~\ref{lem:terminalhistory}, applied to the
deterministic function $K_{x,h}$, gives the exact conditional variance
\begin{equation}\label{eq:terminalvar}
 \sigma_*^2(x)
 :=\operatorname{Var}\bigl(J^{(2)}(x)\mid\F_*\bigr)
 =\sum_{Q\in\pr_*}\int_Q
       |K_{x,h}-(K_{x,h})_Q|^2\dd y.
\end{equation}
The next lemma collects the kernel bounds used to estimate this
variance and the spatial modulus of $J^{(2)}$.

\begin{lemma}
\label{lem:periodic}
For $h$ sufficiently small, uniformly in $x,y\in\T^2$,
\begin{equation}\label{eq:gridreg}
 |K_{x,h}(y)|\le \frac{C}{h+d_{\T}(x,y)}.
\end{equation}
If $y,z$ lie in a common physical set of Euclidean torus
diameter at most $D_*$,
with $r=d_{\T}(x,y)\ge4(D_*+2h)$, then
\begin{equation}\label{eq:faravg}
 |K_{x,h}(y)-K_{x,h}(z)|
 \le C\frac{D_*+h}{r^2}.
\end{equation}
Finally, for $|u|\le1/4$,
\begin{equation}\label{eq:periodictrans}
 \norm{K_{\T}(\cdot+u,\cdot)-K_{\T}(\cdot,\cdot)}_{L^1(\T^2)}
 \le C|u|(1+|\log|u||).
\end{equation}
\end{lemma}
\begin{proof}
Fix a periodic lift at distance less than $1/4$ from the pole.
When the $h$-square containing $y$ is at distance at least
$2h$ from $x$, average \eqref{eq:periodicK} over that square
and use $|x-z|\asymp d_{\T}(x,y)$. When the square is close
to the pole, the integral of $1/|x-z|$ over an $h$-square is $O(h)$,
and division by the area $h^2$ gives $O(h^{-1})$.
The smooth periodic remainder contributes $O(1)$, which proves
\eqref{eq:gridreg}.

For the second claim, let $R_y,R_z$ be the fixed $h$-squares
containing $y,z$. If $\zeta\in R_y$ and $\zeta'\in R_z$,
then $|\zeta-\zeta'|\le D_*+2\sqrt2h$. Every point of the segment
$[\zeta,\zeta']$ lies within $\sqrt2h+D_*+2\sqrt2h\le\frac{3\sqrt2}{2}(D_*+2h)$
of $y$, and $\frac{3\sqrt2}{2}(D_*+2h)\le\frac{3\sqrt2}{8}r<\frac34r$
because $r\ge4(D_*+2h)$. Hence the whole segment stays at distance
at least $r/4$ from the pole. We apply the derivative bound for the
original kernel $K_{\T}$ in \eqref{eq:periodicK} along this segment
and average twice; this proves \eqref{eq:faravg}. The step function
$K_{x,h}$ itself is never differentiated.

For \eqref{eq:periodictrans}, split the torus into
$B(x,2|u|)$ and its complement. In the ball, integrating
the two $1/r$ singularities gives $C|u|$. Outside,
the mean-value theorem gives the pointwise bound $C|u|/r^2$, whose
integral in two dimensions is $C|u|(1+|\log|u||)$.
\end{proof}

\begin{proposition}
\label{prop:termvar}
Under the leaf diameter bound \eqref{eq:Dstar},
\begin{equation}\label{eq:logvar}
 \sup_{x\in\T^2}\sigma_*^2(x)
 \le C\left[1+\log_+\left(\frac{D_*}{h}\right)\right]
 \le C\log m.
\end{equation}
\end{proposition}
The estimate is deterministic once the successful terminal geometry is
fixed, and it uses nothing about how the tree was selected.

\begin{proof}
Put $R_*=8(D_*+2h)$ and split the leaves into those meeting
$B(x,R_*)$ and the remaining ones. Any near leaf lies inside
$B(x,R_*+D_*)\subset B(x,CR_*)$. Subtracting a leaf average
only decreases the $L^2$ norm, so by \eqref{eq:gridreg}
\[
 \sum_{Q\,\mathrm{near}}\int_Q|K_{x,h}-(K_{x,h})_Q|^2
 \le C\int_{B(x,CR_*)}\frac{\dd y}{(h+d_{\T}(x,y))^2}
 \le C\bigl[1+\log_+(D_*/h)\bigr].
\]
For a far leaf $Q$ and $y,z\in Q$ we apply \eqref{eq:faravg}. Since
the leaf is much smaller than its distance to the pole, this gives
\[
 \int_Q|K_{x,h}-(K_{x,h})_Q|^2
 \le C|Q|\frac{(D_*+h)^2}{r_Q^4},
 \quad r_Q=\dist_{\T}(x,Q).
\]
Summing over far leaves and comparing with the integral in
the two-dimensional radial variable bounds their contribution by
\[
 C(D_*+h)^2\int_{R_*/2}^{1}r^{-3}\dd r\le C.
\]
This proves the first inequality in \eqref{eq:logvar}, and the second
follows from \eqref{eq:Dstar}. Borderline cells are avoided by
increasing the universal radius constants, which does not change the
scaling.
\end{proof}

\subsection{Orlicz and supremum estimates}
We now turn the variance bound into estimates for $J^{(2)}$: a
$\psi_2$ bound with fixed positive conditional probability, and an
$L^\infty$ bound with probability $1-o(1)$. Interpolating between the
two gives the $\psi_{r_i}$ bounds.

Write $A_m\in\mathcal F_*$ for the event that the adaptive construction
completes without numerical killing or modular overshoot, so that
$\mathbb P(A_m)\ge p_*>0$ by Proposition~\ref{prop:stopped}. On $A_m$,
conditionally on $\F_*$, each coordinate $J^{(2)}_i(x)$ is centered
Gaussian by Lemma~\ref{lem:terminalhistory}, and its variance is at most
$C_0\log m$ by \eqref{eq:logvar}, with a fixed $C_0$.

We first bound the quadratic exponential modular. For a centered
Gaussian variable $X$ of variance $\sigma^2$ and $t^2>2\sigma^2$,
\[
 \E e^{X^2/t^2}=\Bigl(1-\frac{2\sigma^2}{t^2}\Bigr)^{-1/2}.
\]
With $A_1^2=8C_0/3$ and $t^2=A_1^2\log m$ one has
$2\sigma^2/t^2\le3/4$, hence
\[
 \E\left[\exp\left(\frac{|J^{(2)}_i(x)|^2}
                    {A_1^2\log m}\right)\Big|\F_*\right]\le2
 \qquad(x\in\T^2,\ i=1,2).
\]
By conditional Tonelli the spatial integral has conditional
expectation at most $2$. Markov's inequality at level $16$, applied
to each of the two coordinates, shows that with conditional
probability at least $1-2\cdot\frac2{16}=\frac34$,
\[
 \int_{\T^2}\exp\left(\frac{|J^{(2)}_i|^2}{A_1^2\log m}\right)\le16,
 \qquad i=1,2.
\]
On this event, Jensen's inequality for the concave map
$t\mapsto t^{1/4}$ on the unit-area torus gives
$\int_{\T^2}\exp\bigl(|J^{(2)}_i|^2/(4A_1^2\log m)\bigr)\le16^{1/4}=2$,
that is, $\int_{\T^2}\psi_2\bigl(J^{(2)}_i/(2A_1\sqrt{\log m})\bigr)\le1$.
Consequently the event
\[
 C_m=\Bigl\{\max_{i=1,2}\|J_i^{(2)}\|_{\psi_2}
             \le 2A_1\sqrt{\log m}\Bigr\}
\]
satisfies
\begin{equation}\label{eq:ps2corr}
 \mathbb P(C_m\mid\mathcal F_*)\ge\tfrac34
 \qquad\text{on }A_m.
\end{equation}

The $L^\infty$ bound requires a separate high-probability argument,
since quadratic exponential integrability in space alone does not
control the supremum. Recall that $g^{(2)}=W_h-P_hPW$, and let $D_m$
be the intersection of the fixed-$h$ Gaussian maximum event and the
predictable adaptive-cell maximum event. Since $s\gg h$, on $D_m$
\[
 \norm{g^{(2)}}_\infty\le C\sqrt m/h.
\]
The fixed-grid event follows from the ordinary union bound over
$h^{-2}$ deterministic Gaussians, and the adaptive event is
\eqref{eq:randommaxwhp}. Hence $\mathbb P(D_m^c)=o(1)$; this bound is unconditional.
Since $J^{(2)}=K_{\T}*g^{(2)}$, the translation estimate
\eqref{eq:periodictrans} gives on $D_m$ the deterministic modulus
\[
 |J^{(2)}(x)-J^{(2)}(z)|
 \le C\frac{\sqrt m}{h}\,d_{\T}(x,z)
             \bigl(1+|\log d_{\T}(x,z)|\bigr).
\]
Fix a deterministic torus net $\mathcal N_m$ of mesh $q_m=2^{-Lm}$,
where $L>3$ is fixed. It has at most $C2^{2Lm}$ points. Since
$h=2^{-m}$, the modulus above is $O(m^{3/2}2^{-(L-1)m})=o(1)$ at
distance $q_m$. Let $N_m$ be the event that
$|J^{(2)}(y)|\le A_0\sqrt{m\log m}$ for every $y\in\mathcal N_m$.
The conditional Gaussian tails at variance level $C_0\log m$ and a
union bound over the net give, on $A_m$,
\[
 \mathbb P(N_m^c\mid\mathcal F_*)
 \le C2^{2Lm}e^{-A_0^2m/(2C_0)}=o(1)
\]
once $A_0^2>4C_0L\log2$. This estimate uses neither the fixed-grid
nor the adaptive-cell maximum event. On $D_m\cap N_m$ the modulus
extends the net bound to
\begin{equation}\label{eq:supcorr}
 \norm{J^{(2)}}_\infty\le C\sqrt{m\log m}.
\end{equation}

For a scalar function $f$ and fixed $r>2$, the elementary
inequality
\[
 \left(\frac{|f|}{A^{1-2/r}B^{2/r}}\right)^r
 \le \left(\frac{|f|}{B}\right)^2
 \quad\text{whenever}\quad \norm f_\infty\le A
\]
shows, after a fixed dilation, that
\[
 \norm f_{\psi_r}\le C_r
 \norm f_\infty^{1-2/r}\norm f_{\psi_2}^{2/r}.
\]
Consequently, on $C_m\cap D_m\cap N_m$,
\begin{equation}\label{eq:interp}
 \norm{J^{(2)}_i}_{\psi_{r_i}}
 \le C_{r_i}m^{1/2-1/r_i}\sqrt{\log m}
 =o(\sqrt m)\qquad (r_i<\infty).
\end{equation}
The gain over $\sqrt m$ comes from the finiteness of $r_i$ and is not
uniform as $r_i\to\infty$.

\subsection{Joint probability estimate}
Since $A_m\in\mathcal F_*$, \eqref{eq:ps2corr} gives
$\mathbb P(A_m\cap C_m)=\E[\one_{A_m}\mathbb P(C_m\mid\F_*)]
\ge\frac34\mathbb P(A_m)$. The remaining two events are removed by
subtracting their complements, which requires no conditional
independence:
\begin{align*}
 \mathbb P(A_m\cap C_m\cap D_m\cap N_m)
 &\ge\mathbb P(A_m\cap C_m)
      -\mathbb P(D_m^c)-\mathbb P(A_m\cap N_m^c)\\
 &\ge\tfrac34\,\mathbb P(A_m)-o(1)
 \ge\tfrac34\,p_*-o(1)>0.
\end{align*}
On this intersection, put
$J=J_{\mathrm{leaf}}+J^{(1)}+J^{(2)}$. The exact splitting
\eqref{eq:decomp} gives $\divv J=W_h$, and the three estimates are
\[
 \max_i\|J_{\mathrm{leaf},i}\|_{\psi_{r_i}}\le C\sqrt m,
 \qquad \|J^{(1)}\|_\infty\le C m^{-5/2},
 \qquad
 \|J^{(2)}_i\|_{\psi_{r_i}}
       \le C_{r_i}m^{1/2-1/r_i}\sqrt{\log m}.
\]
In the quartic case the last bound is
$C m^{1/4}\sqrt{\log m}=o(\sqrt m)$, and for every fixed finite
pair the same strict improvement follows from \eqref{eq:interp}.
The triangle inequality therefore gives, for some fixed $p_{**}>0$,
\begin{equation}\label{eq:fixedprob}
 \mathbb P\bigl(\mathfrak F_{\boldsymbol r}(W_h)
                         \le C\sqrt m\bigr)\ge p_{**}
\end{equation}
for all sufficiently large $m$.

\paragraph{Completion of the proof of Theorem~\ref{thm:gauss}.}
The estimate \eqref{eq:fixedprob} is the required
fixed-positive-probability statement. Taking any
$0<\varepsilon_0<p_{**}$ gives \eqref{eq:target} and
proves Theorem~\ref{thm:gauss}. The upgrade to a high-probability
bound by Gaussian concentration, with constants independent of the
mesh, is proved in Part~\ref{umc:partII}; see
Remark~\ref{rem:concentration-upgrade}.



\clearpage
\part{From matching to dyadic Gaussian Orlicz flows}\label{umc:partII}

This part reduces the matching problem to a Gaussian minimum-flow problem
on the square (\cref{red:thm:main}), for every critical exponent pair
including the hard endpoints, and transfers the torus estimate of
Part~\ref{umc:partI} to the square (\cref{sec:fold}).

\section{Overview of the reduction}\label{sec:interface}

Fix a critical exponent pair $\boldsymbol r=(r_1,r_2)$, with
$1/r_1+1/r_2=1/2$. Let $\mathfrak M_N^{\boldsymbol r}$ denote the minimum,
over permutations of two samples of $N$ independent uniform points in
$[0,1]^2$, of the maximum of their coordinatewise Orlicz matching norms.
The natural scale is $s_N=\sqrt{(\log N)/N}$. We relate a bound at this
scale to the minimum Orlicz norm of a current whose divergence is Gaussian
white noise projected onto the dyadic partition of mesh $h_N\asymp s_N$.
The two statements are equivalent at the level of high-probability bounds
with deterministic constants; the equivalence does not identify the
constants or the minimum costs. The one-cloud coupling also gives
Talagrand's original matching problem with any evenly spread deterministic
reference configuration (\cref{red:cor:evenly-spread}).

For finite critical exponents, \cref{thm:gauss} in Part~\ref{umc:partI}
gives a fixed-positive-probability bound for the Gaussian current on the
flat torus. The concentration principle of
\cref{red:subsec:gaussian-concentration} upgrades it to high probability
before the matching reduction is applied. The reduction also covers the
hard endpoint, for which no Gaussian flow is constructed; see
Part~\ref{umc:partIII}. No Gaussian flow estimate is assumed in the proof
of \cref{red:thm:main}, and the concentration step does not use the flow
construction of Part~\ref{umc:partI}; the proof of \cref{red:thm:main}
does not use Part~\ref{umc:partI}.

The reduction uses the projection onto dyadic cells, rather than a
Fourier or heat regularization. The random source produced by the
reduction is then exactly the cellwise Gaussian source of the flow
construction.

We now outline the argument. Choose $h_N$ so that $Nh_N^2$ is comparable
to $\log N$. With high probability, both empirical histograms then have
densities bounded above and away from zero. Spreading the atoms uniformly
within their cells changes each transport displacement by at most a
multiple of $h_N$. Conversely, a coupling of the histograms gives a doubly
stochastic matrix on the original labels, and a Birkhoff decomposition
yields a single permutation controlling both coordinates. This reduces
matching to transport between the two histograms
(\cref{red:prop:atomic-hist}).

For bounded positive densities, transport with coordinatewise Orlicz
costs is comparable to the minimum Orlicz norm of a static divergence
current (\cref{red:thm:transport-flow-finite,red:thm:transport-flow-endpoint}).
One direction follows by linearly interpolating densities along
a given current. The reverse direction uses straight-segment interpolation
associated with a strictly convex translation cost. The relevant density
bound is obtained from noncrossing of transport rays and a Jacobian
inequality; we give the details needed at a hard coordinate endpoint.

The remaining random step is a coupling of the histogram discrepancy with
white noise on the same dyadic grid (\cref{red:thm:KMT}). Encode the fixed
binary rectangle tree by nested intervals and apply the classical
one-dimensional KMT coupling to a Brownian bridge. Its increments give
exactly the Gaussian cell masses. The rectangle discrepancies have an
explicit Haar divergence inverse whose $L^\infty$ norm is at most $Cs_N$
with probability tending to one. This controls every coordinate
Orlicz norm, including an $L^\infty$ coordinate, and suffices to
transfer bounds up to fixed constants.

Finally, the Gaussian statement on the square follows from the analogous
torus statement (\cref{sec:fold}). Folding the unit torus onto the square
pushes a periodic current to a no-flux current and sends white noise
projected at mesh $h/2$ to square white noise projected at mesh $h$.
The equality of the Gaussian source laws is exact. We prove only the
needed direction, from torus to square; folding does not give an inverse
comparison.

\section{Notation and statement of the reduction}\label{red:sec:matching}

Let $D=[0,1]^2$.  Let $(X_i)_{i\ge1}$ and $(Y_i)_{i\ge1}$ be independent
sequences of independent uniform random variables in $D$, and set
\[
 \mu_N=\frac1N\sum_{i=1}^N\delta_{X_i},
 \qquad
 \nu_N=\frac1N\sum_{i=1}^N\delta_{Y_i},
 \qquad
 s_N=\sqrt{\frac{\log N}{N}}.
\]

\subsection{Matching and Orlicz norms}

Throughout, the exponent pair $\boldsymbol r=(r_1,r_2)$ is fixed and
satisfies the critical relation
\begin{equation}
 2\le r_1,r_2\le\infty,\qquad
 \frac1{r_1}+\frac1{r_2}=\frac12,
 \qquad\frac1\infty:=0.
 \label{red:eq:critical-pair}
\end{equation}
Thus either $2<r_1,r_2<\infty$, or $\boldsymbol r=(2,\infty)$ or
$(\infty,2)$. We put $a_i=1/r_i$ with $a_i=0$ at the infinite endpoint.
For $2\le r<\infty$, put
\[
 \psi_r(t)=e^{|t|^r}-1.
\]
For the endpoint, put
\[
 \psi_\infty(t)=
 \begin{cases}
 0,&|t|\le1,\\
 +\infty,&|t|>1.
 \end{cases}
\]
The Luxemburg norm associated with $\psi_\infty$ is the $L^\infty$ norm.
For the exponent pair $\boldsymbol r$ we use the coordinate Young functions $\psi_{r_1}$ and $\psi_{r_2}$ throughout.

For a finite vector $b=(b_i)_{i=1}^N$, define
\[
 \norm b_{\psi_{r_j},N}
 :=\inf\set{c>0:\frac1N\sum_{i=1}^N\psi_{r_j}(b_i/c)\le1}.
\]
For $r_j=\infty$, this is $\max_i |b_i|$.

For a permutation $\sigma\in S_N$, write
$d_{i,j}^\sigma=\abs{X_i^j-Y_{\sigma(i)}^j}$ and define
\begin{equation}
 \mathfrak M_N^{\boldsymbol r}
 :=\min_{\sigma\in S_N}
 \max_{j=1,2}\norm{(d_{i,j}^\sigma)_{i=1}^N}_{\psi_{r_j},N}.
 \label{red:eq:atomic-aniso}
\end{equation}
The two-sample form of the Ultimate Matching Conjecture is the following.

\begin{definition}
\label{red:def:UMC}
The conjecture $\mathrm{UMC}_{\boldsymbol r}$ is the assertion that
there is a deterministic $C<\infty$ such that
\begin{equation}
 \Pp\pa{\mathfrak M_N^{\boldsymbol r}\le Cs_N}\longrightarrow1.
 \label{red:eq:UMC}
\end{equation}
\end{definition}

Talagrand's Research Problem~17.1.2 \cite{umc:Talagrand2021} states the
same coordinatewise criterion for a random sample and an evenly spread
reference configuration.
The symmetric quartic case is $\boldsymbol r=(4,4)$; the hard coordinate
endpoint is $\boldsymbol r=(2,\infty)$.

\begin{lemma}
\label{red:lem:isotropic-quartic}
For $\boldsymbol r=(4,4)$, the coordinate gauge in
\eqref{red:eq:atomic-aniso} is equivalent, up to a factor $\sqrt2$, to the
isotropic gauge generated by $e^{|z|^4}-1$.
\end{lemma}

\begin{proof}
The lower bound follows from $|z_j|\le|z|$.  Conversely, if both coordinate
Luxemburg norms are at most $a$, then
$|z|^4\le2(|z_1|^4+|z_2|^4)$ and Cauchy--Schwarz yields
\[
 \frac1N\sum_i
 e^{|z_i|^4/(\sqrt2a)^4}
 \le
 \prod_{j=1}^2
 \left(\frac1N\sum_i e^{|z_{i,j}|^4/a^4}\right)^{1/2}
 \le2.
\]
\end{proof}

\subsection{Dyadic Gaussian flows}

Write $\T^2=(\R/\mathbb Z)^2$.  For $\Omega\in\{D,\T^2\}$,
$P_h^{\Omega}$ denotes averaging onto the fixed $h$-square grid
in $\Omega$, acting on mean-zero functions when appropriate.  With
mean-zero white noise $W^{\Omega}$ (the isonormal process on
$L_0^2(\Omega)$), set
\begin{equation}\label{red:eq:compatible-source}
  W_h^{\Omega}=P_h^{\Omega}W^{\Omega}.
\end{equation}
On the torus we also write $W_h=W_h^{\T}$ and
$\mathfrak F_{\boldsymbol r}=\mathfrak F_{\boldsymbol r}^{\T}$ when
there is no ambiguity.

Choose a sufficiently large constant $A$ and a dyadic $h_N$ satisfying
\begin{equation}
 As_N\le h_N<2As_N.
 \label{red:eq:h-choice}
\end{equation}
The Gaussian source on the square used in the reduction is
\begin{equation}
 W_{h_N}^{D}=P_{h_N}^{D}W^{D},
 \label{red:eq:haar-source}
\end{equation}
with the same white-noise normalization as $W_h$ on the torus.

For a vector field $J=(J_1,J_2)$, put
\begin{equation}
 \norm J_{\boldsymbol r}
 :=\max_{j=1,2}\norm{J_j}_{\psi_{r_j}}.
 \label{red:eq:product-flow-norm}
\end{equation}
For a mean-zero distribution $f$, we encode both the divergence equation
and the no-flux boundary condition weakly.  A vector field $J$ is admissible
for $f$ when
\begin{equation}
 -\int_D J\cdot\nabla\varphi\,dx=\langle f,\varphi\rangle
 \qquad\text{for every }\varphi\in C^\infty(\overline D).
 \label{red:eq:weak-div-no-flux}
\end{equation}
For regular $J$, this is exactly $\divv J=f$ in $D$ together with the zero
normal trace $J\cdot n=0$ on $\partial D$.  We then define
\begin{equation}
 \mathfrak F_{\boldsymbol r}^{D}(f)
 :=\inf\set{\norm J_{\boldsymbol r}: J\text{ satisfies }
 \eqref{red:eq:weak-div-no-flux}}.
 \label{red:eq:aniso-flow}
\end{equation}

The corresponding Gaussian minimum-flow problem is the following.

\begin{definition}
\label{red:def:GOMF}
The conjecture $\mathrm{GOMF}_{\boldsymbol r}$ asserts that there is a
deterministic $C<\infty$ such that
\begin{equation}
 \Pp\pa{\mathfrak F_{\boldsymbol r}^{D}(W_{h_N}^{D})
 \le C\sqrt{\log N}}\longrightarrow1.
 \label{red:eq:GOMF}
\end{equation}
\end{definition}

The main result of this part is the following equivalence.

\begin{theorem}
\label{red:thm:main}
For every fixed pair $\boldsymbol r$ satisfying
\eqref{red:eq:critical-pair}, and for $A$ sufficiently large,
$\mathrm{UMC}_{\boldsymbol r}$ and
$\mathrm{GOMF}_{\boldsymbol r}$ are equivalent.  More precisely, the
following statements are equivalent:
\begin{enumerate}[label=\textnormal{(\roman*)}]
 \item $\Pp(\mathfrak M_N^{\boldsymbol r}\le Cs_N)\to1$ for some $C$;
 \item
 \[
 \Pp\pa{\mathfrak F_{\boldsymbol r}^{D}
 (P_{h_N}^D(\mu_N-\nu_N))\le Cs_N}\to1
 \]
 for some $C$;
 \item
 \[
 \Pp\pa{\mathfrak F_{\boldsymbol r}^{D}(W_{h_N}^{D})
 \le C\sqrt{\log N}}\to1
 \]
 for some $C$.
\end{enumerate}
All constants are deterministic and independent of $N$.
\end{theorem}

\section{Preliminary estimates}

We collect elementary facts on Orlicz norms and on the minimum-flow
gauge. We then show that the gauge of projected white noise concentrates
around its median, so that a bound with fixed positive probability
implies a bound with high probability.

\begin{lemma}
\label{red:lem:moment}
For every fixed $2\le r<\infty$ on a probability space,
\begin{equation}
 \norm Z_{\psi_r}
 \asymp_r
 \sup_{s\ge1}s^{-1/r}\norm Z_{L^s}.
 \label{red:eq:moment-characterization}
\end{equation}
\end{lemma}

\begin{proof}
The forward implication follows by integrating the tail furnished by
$\E e^{|Z|^r/a^r}\le2$.  Conversely, expand the exponential in
moments and use $k!\ge(k/e)^k$.
\end{proof}

\begin{lemma}
\label{red:lem:absorb}
Let $\Theta$ be any Young function and $K\ge1$.  If
$\int\Theta(|f|/r)\le K$, then $\norm f_{L^\Theta}\le Kr$.
\end{lemma}

\begin{proof}
Convexity and $\Theta(0)=0$ give $\Theta(t/K)\le K^{-1}\Theta(t)$.
\end{proof}

\begin{lemma}
\label{red:lem:flow-lipschitz}
For mean-zero sources $f,g$,
\[
 \abs{\mathfrak F_{\boldsymbol r}^{D}(f)-
       \mathfrak F_{\boldsymbol r}^{D}(g)}
 \le \mathfrak F_{\boldsymbol r}^{D}(f-g),
 \qquad
 \mathfrak F_{\boldsymbol r}^{D}(af)=|a|\mathfrak F_{\boldsymbol r}^{D}(f).
\]
\end{lemma}

\begin{proof}
Subtract admissible currents and use the triangle inequality in the product
norm \eqref{red:eq:product-flow-norm}.
\end{proof}

\subsection{Gaussian concentration for dyadic minimum-flow gauges}
\label{red:subsec:gaussian-concentration}

The argument of this subsection uses neither the matching reduction,
nor the transport--flow comparison, nor any construction of the
Gaussian current. It applies without change when one of the exponents
equals $\infty$.

For $\Omega\in\{D,\T^2\}$, let $\mathfrak F_{\boldsymbol r}^{\Omega}$
be the minimum-flow gauge of \eqref{red:eq:aniso-flow}, with the
periodic divergence convention in the torus case. If $h=2^{-m}$,
write its projected mean-zero white noise in the form
\begin{equation}\label{red:eq:gaussian-coefficients}
 W_h^{\Omega}(Z)|_Q=h^{-1}(Z_Q-\overline Z),
 \qquad \overline Z=h^2\sum_Q Z_Q,
\end{equation}
where the $h$-grid coefficients $Z_Q$ are independent standard
Gaussian variables.

\begin{lemma}
\label{red:lem:uniform-gaussian-lipschitz}
Fix any $r_i\in[2,\infty]$. There is an absolute constant $C$,
independent of $h$, the number of cells and the exponent pair, such
that for $\Omega\in\{D,\T^2\}$ and coefficient arrays $Z,Z'$,
\begin{equation}\label{red:eq:uniform-gaussian-lipschitz}
\left|\mathfrak F_{\boldsymbol r}^{\Omega}(W_h^{\Omega}(Z))-
\mathfrak F_{\boldsymbol r}^{\Omega}(W_h^{\Omega}(Z'))\right|
 \le C\norm{Z-Z'}_{\ell^2}.
\end{equation}
\end{lemma}
\begin{proof}
Write $v=Z-Z'$ and $g_Q=h(v_Q-\overline v)$ for the mass
of the source difference in a grid square $Q$. For any union $A$
of grid squares, centering and Cauchy--Schwarz give
\[
 \left|\sum_{Q\in A}g_Q\right|
 \le 2h\norm v_{\ell^2}\sqrt{\min\{\#A,\#A^c\}}.
\]
The nearest-neighbor square grid satisfies the relative
edge-isoperimetric inequality
\[
 |\partial A|\ge \sqrt2\,\sqrt{\min\{\#A,\#A^c\}},
\]
where $\partial A$ counts internal edges (and also periodic seam
edges for $\Omega=\T^2$). Indeed, for the $k\times k$ grid with
$k=h^{-1}$ and $\#A\le k^2/2$, the edge-isoperimetric inequality of
Bollob\'as and Leader \cite{BollobasLeader1991} gives
$|\partial A|\ge\min\{2\sqrt{\#A},k\}\ge\sqrt2\sqrt{\#A}$; the torus
graph contains all edges of the square grid, so the same bound holds
for $\Omega=\T^2$. Consequently every cut has enough
capacity when each internal edge is assigned capacity
$C h\norm v_{\ell^2}$ in both directions. Gale's supply--demand
theorem \cite{Gale1957} (see also \cite{Hoffman1960}), a consequence
of the max-flow/min-cut theorem, then
produces a consistent signed edge flux $f_e$ with
\[
 \sum_{e\,\mathrm{out}(Q)}f_e-
 \sum_{e\,\mathrm{in}(Q)}f_e=g_Q,
 \qquad |f_e|\le C h\norm v_{\ell^2}.
\]
On $D$, take the flux on exterior faces to be zero.

It remains to realize these integrated face fluxes as a
continuum current. On $Q=[ih,(i+1)h]\times[jh,(j+1)h]$, put
$t=(x-ih)/h$ and $u=(y-jh)/h$, and use globally consistent
rightward/upward signed fluxes $f_L,f_R,f_B,f_T$. Define
\begin{equation}\label{red:eq:face-current}
 J_1(x,y)=\frac{(1-t)f_L+tf_R}{h},\qquad
 J_2(x,y)=\frac{(1-u)f_B+uf_T}{h}.
\end{equation}
The normal components match across common faces (and periodically
across torus seams). Their weak divergence on each cell is
\[
 \divv J=\frac{f_R-f_L+f_T-f_B}{h^2}
 =h^{-1}(v_Q-\overline v).
\]
This current satisfies the zero-normal-trace convention on $D$ and
$\max_i\norm{J_i}_{L^\infty}\le C\norm v_{\ell^2}$. Every
Luxemburg $\psi_r$-norm with $r\ge2$, including $r=\infty$,
is bounded by a universal multiple of $L^\infty$ on a space of
unit area. The triangle inequality for minimum-flow gauges,
\cref{red:lem:flow-lipschitz}, now gives
\eqref{red:eq:uniform-gaussian-lipschitz}.
\end{proof}

\begin{proposition}
\label{red:prop:gaussian-concentration}
Let $h=2^{-m}$, $\Omega\in\{D,\T^2\}$, and fix
$r_1,r_2\in[2,\infty]$. Suppose there are $B,p_0>0$ such that,
for all sufficiently large $m$,
\[
 \Pp\!\left(\mathfrak F_{\boldsymbol r}^{\Omega}(W_h^{\Omega})
 \le B\sqrt m\right)\ge p_0.
\]
Then there are $B',c>0$ such that, for all sufficiently large $m$,
\begin{equation}\label{red:eq:gaussian-concentration}
 \Pp\!\left(\mathfrak F_{\boldsymbol r}^{\Omega}(W_h^{\Omega})
 >B'\sqrt m\right)\le2e^{-cm}.
\end{equation}
\end{proposition}
\begin{proof}
By \cref{red:lem:uniform-gaussian-lipschitz}, the gauge is an
$L$-Lipschitz function of a standard Gaussian vector, with $L$
independent of the dimension. Gaussian concentration around a
median $M_m$ shows that
$M_m\le B\sqrt m+C(L,p_0)$. Indeed, otherwise the lower tail
below $B\sqrt m$ would have probability less than $p_0$.
The upper concentration tail beyond a fixed larger multiple
$B'\sqrt m$ is then at most $2e^{-cm}$.
\end{proof}

\begin{corollary}
\label{red:cor:positive-probability-equivalence}
In the dyadic Gaussian condition (iii) of \cref{red:thm:main}, it is
equivalent to require the existence of $B,p_0>0$ such that
\[
 \Pp\!\left(\mathfrak F_{\boldsymbol r}^{D}
 (W_{h_N}^{D})\le B\sqrt{\log N}\right)\ge p_0
\]
for all sufficiently large $N$. The same conclusion holds for a torus
flow with its periodic divergence constraint.
\end{corollary}
\begin{proof}
The high-probability statement implies the positive-probability bound.
Conversely, $h_N=2^{-m(N)}$ and $m(N)\asymp\log N$, so
\cref{red:prop:gaussian-concentration} applies with a changed
constant and yields the Gaussian condition (iii).
\end{proof}
Thus the Gaussian input to the matching equivalence need only be proved
with a fixed positive probability.

\section{From point configurations to histograms}

We show that, up to an additive error of order $h$, matching the two
point clouds is equivalent to transporting their dyadic histograms. This
replaces the atomic problem by a transport problem between densities.

For a coupling $\pi$ on $D\times D$, define
\begin{equation}
 \mathfrak T_{\boldsymbol r}(\varrho_0,\varrho_1)
 :=\inf_{\pi\in\Pi(\varrho_0\dd x,\varrho_1\dd x)}
 \max_{j=1,2}\norm{x_j-y_j}_{L^{\psi_{r_j}}(\pi)}.
 \label{red:eq:aniso-transport-gauge}
\end{equation}
Let $P_h^D\mu_N$ be obtained by spreading each atomic mass uniformly on its
dyadic cell.

\begin{proposition}
\label{red:prop:atomic-hist}
There is a universal $C$ such that, for every pair of point clouds and every
dyadic $h$,
\begin{equation}
 \mathfrak T_{\boldsymbol r}(P_h^D\mu_N,P_h^D\nu_N)
 \le \mathfrak M_N^{\boldsymbol r}+Ch,
 \qquad
 \mathfrak M_N^{\boldsymbol r}
 \le C\bigl(\mathfrak T_{\boldsymbol r}(P_h^D\mu_N,P_h^D\nu_N)+h\bigr).
 \label{red:eq:atomic-hist-equivalence}
\end{equation}
The constant is uniform over the critical exponents, including an
$L^\infty$ endpoint.
\end{proposition}

\begin{proof}
Fix a permutation.  Couple the two associated cells using the same relative
coordinate.  Each coordinate displacement increases by at most $2h$, and
the first inequality follows from the Luxemburg triangle inequality.

Conversely, let $\pi$ couple the histograms.  Attach posterior atomic labels
inside every occupied cell and define
\[
 q_{ik}=N\int\lambda_i(x)\mu_k(y)\dd\pi(x,y).
\]
Then $Q=(q_{ik})$ is doubly stochastic.  If the histogram coordinate norms
are at most $r$, the fractional labelled coupling has both coordinate norms
at most $r+2h$.  Decompose $Q$ into permutation matrices.  For finite
coordinates, the average over this decomposition of the sum of the two
coordinate modulars is at most $2$, so one permutation has that sum at most
$2$; \cref{red:lem:absorb} absorbs this fixed factor.  For an $L^\infty$
coordinate, every positive entry of $Q$ obeys the corresponding displacement
bound, and a Birkhoff decomposition can be chosen inside the positive
support.  Thus the same permutation satisfies the hard coordinate bound.
\end{proof}

\section{Transport and divergence currents}

This section contains the deterministic core of the reduction. For
densities bounded above and away from zero, we prove that the
anisotropic transport gauge $\mathfrak T_{\boldsymbol r}$ is comparable
to the minimum-flow gauge $\mathfrak F_{\boldsymbol r}^D$ of their
difference, first for finite exponents and then at the endpoint
$(2,\infty)$. The passage from transport to flow requires a bound on
the densities along displacement interpolation, which we derive from
noncrossing of transport rays and a Jacobian inequality. The domain
$D$ is convex, so straight segments remain in $D$.

\subsection{Noncrossing of transport rays}

\begin{lemma}
\label{red:lem:noncrossing}
Let $L:\R^2\to\R\cup\{+\infty\}$ be strictly convex on its convex effective
domain.  Let $\Gamma\subset D\times D$ be $c$-cyclically monotone for
$c(x,y)=L(y-x)$.  For $0<t<1$, the map
\[
 (x,y)\longmapsto(1-t)x+ty
\]
is injective on $\Gamma$.
\end{lemma}

\begin{proof}
Suppose $(1-t)x_0+ty_0=(1-t)x_1+ty_1$ and put $v_i=y_i-x_i$.
Then $x_1-x_0=t(v_0-v_1)$, and the crossed displacements are
\[
 y_1-x_0=tv_0+(1-t)v_1,
 \qquad
 y_0-x_1=(1-t)v_0+tv_1.
\]
They remain in the convex effective domain.  If $v_0\ne v_1$, strict
convexity gives
\[
 L(y_1-x_0)+L(y_0-x_1)<L(v_0)+L(v_1),
\]
contradicting two-point cyclical monotonicity.  Hence $v_0=v_1$, and then
$x_0=x_1$ and $y_0=y_1$.
\end{proof}

\subsection{Density bounds for displacement interpolation}

\begin{lemma}
\label{red:lem:jacobian-package}
Let $D\subset\R^2$ be bounded and convex, let $L\in C^2(\R^2)$, and assume
that $D^2L$ is uniformly positive definite on the bounded displacement set
$D-D$.  Let $c(x,y)=L(y-x)$ and let $\varrho_0,\varrho_1$ be smooth densities on
$D$ satisfying $0<\underline\varrho\le\varrho_i\le\overline\varrho$.  Then the $c$-optimal plan is induced
by a unique map $T$.  Moreover $T$ is approximately differentiable almost
everywhere, a Kantorovich potential is twice differentiable in the Alexandrov
sense almost everywhere, and at almost every point where these differentials
exist one has the measure-theoretic Jacobian identity
\begin{equation}
 \varrho_0(x)=\varrho_1(T(x))\,|\det DT(x)|.
 \label{red:eq:jacobian-map}
\end{equation}
If, for some $t\in(0,1)$, the interpolation map
$T_t=(1-t)\Id+tT$ is injective on a full-$\varrho_0\,dx$-measure transport set,
then the area formula applies without multiplicity and
\begin{equation}
 \varrho_t(T_t(x))\,|\det DT_t(x)|=\varrho_0(x)
 \label{red:eq:jacobian-interpolation}
\end{equation}
almost everywhere on that set, where
$\varrho_t\,dx=(T_t)_\#(\varrho_0\,dx)$.
\end{lemma}

\begin{proof}
For the translation cost considered here, strict convexity gives the twist
condition and $C^2$ regularity on the bounded displacement set gives the local
semiconcavity hypotheses.  Existence and uniqueness of the optimal map are a
special case of Villani's Monge theorem; see \cite[Thm.~10.28 and
Ex.~10.35]{Villani}.  Once injectivity and approximate
differentiability are available, the Jacobian identity is precisely the
Jacobian equation of \cite[Thm.~11.1]{Villani}.  This theorem applies to
$T_t=(1-t)\Id+tT$ because approximate differentiability is
preserved by affine combinations with the identity; the noncrossing lemma
supplies injectivity for $0<t<1$.  This yields
\eqref{red:eq:jacobian-interpolation} without a multiplicity factor.  Only
these almost-everywhere properties are used; no classical $C^1$ regularity
of the optimal map is needed.
\end{proof}

\begin{lemma}
\label{red:lem:density-interpolation}
Let $L\in C^2(\R^2)$ have Hessian uniformly positive definite on $D-D$,
and let $T$ be the optimal map for the cost $L(y-x)$ between smooth densities satisfying
$0<\underline\varrho\le\varrho_0,\varrho_1\le\overline\varrho$.  Put $T_t=(1-t)\Id+tT$ and
$\varrho_t=(T_t)_\#(\varrho_0\dd x)$.  Then
\begin{equation}
 \varrho_t\le \varrho_*:=\frac{\overline\varrho^2}{\underline\varrho}
 \qquad(0\le t\le1).
 \label{red:eq:rho-bound}
\end{equation}
\end{lemma}

\begin{proof}
Choose the Kantorovich-potential convention
$\phi(x)+\chi(y)\le L(y-x)$ with equality at $y=T(x)$.  Put
$A=D^2L(T(x)-x)$.  By \cref{red:lem:jacobian-package}, the required approximate
differentials and the Jacobian identity exist almost everywhere.  At Alexandrov points, touching by
$x'\mapsto L(T(x)-x')-\chi(T(x))$ gives $D^2\phi\le A$, while
differentiating $\nabla\phi(x)=-\nabla L(T(x)-x)$ gives
\[
 D^2\phi=A(I-DT).
\]
Consequently
\[
 DT=I-A^{-1}D^2\phi
 =A^{-1/2}\bigl(I-A^{-1/2}D^2\phi A^{-1/2}\bigr)A^{1/2},
\]
so $DT$ is similar to the symmetric matrix
\[
 B=I-A^{-1/2}D^2\phi A^{-1/2}\ge0.
\]
Thus the eigenvalues $\lambda_1,\lambda_2$ of $DT$ are real and nonnegative.
In particular $\det DT\ge0$, so the absolute value in
\eqref{red:eq:jacobian-map} can now be removed.  Fix $t\in(0,1)$.  By
\cref{red:lem:noncrossing}, $T_t$ is injective on the transport graph.  Thus
\cref{red:lem:jacobian-package} applies without multiplicity.  Moreover
$\det DT_t>0$ because $(1-t)+t\lambda_k>0$, so the absolute value in
\eqref{red:eq:jacobian-interpolation} can also be removed.  We have
\[
 \det DT_t
 =\prod_{k=1}^2((1-t)+t\lambda_k)
 \ge\br{(1-t)+t\sqrt{\det DT}}^2.
\]
By \eqref{red:eq:jacobian-map}, now without absolute values,
$\det DT=\varrho_0/(\varrho_1\circ T)\ge \underline\varrho/\overline\varrho$, hence
$\det DT_t\ge \underline\varrho/\overline\varrho$.  Using \eqref{red:eq:jacobian-interpolation},
\[
 \varrho_t(T_t(x))=\frac{\varrho_0(x)}{\det DT_t(x)}
 \le\frac{\overline\varrho^2}{\underline\varrho}.
\]
At $t=0,1$ the same bound follows directly from
$\varrho_0,\varrho_1\le\overline\varrho\le \overline\varrho^2/\underline\varrho$ (the unit-mass bounds imply
$\underline\varrho\le1\le\overline\varrho$).
\end{proof}

\begin{remark}
With the opposite convention for the Kantorovich potential, all signs in the
intermediate identities reverse simultaneously.  The invariant statement is
that $DT$ is similar to a symmetric positive semidefinite matrix.
\end{remark}

\subsection{Boundary smoothing and approximation}

The previous lemma requires smooth densities, whereas the histogram densities
used later are piecewise constant. The passage between them requires a
recovery approximation, not merely lower semicontinuity. We use the
following local construction on the square.

\begin{lemma}
\label{red:lem:local-smoothing}
Let $D=[0,1]^2$.  There are Markov kernels $K_\tau(x,\dd z)$,
$0<\tau<1/4$, with smooth nonnegative densities $k_\tau(x,z)$ such that:
\begin{enumerate}[label=\textup{(\roman*)}]
\item $K_\tau$ preserves Lebesgue measure and is symmetric:
\[
 \int_D k_\tau(x,z)\,\dd z=1,
 \qquad
 \int_D k_\tau(x,z)\,\dd x=1,
 \qquad
 k_\tau(x,z)=k_\tau(z,x).
\]
\item $k_\tau(x,z)=0$ unless $|x-z|_\infty\le\tau$.
\item If $S_\tau\varrho(z)=\int_D k_\tau(x,z)\varrho(x)\,\dd x$, then
$S_\tau\varrho$ is smooth, $S_\tau\varrho\to\varrho$ in $L^1(D)$, and
\[
 \underline\varrho\le\varrho\le\overline\varrho \quad\Longrightarrow\quad
 \underline\varrho\le S_\tau\varrho\le\overline\varrho.
\]
\item If $\pi$ couples $\varrho_0\dd x$ and $\varrho_1\dd x$, then
\[
 \pi_\tau(\dd x'\dd y')
 :=\int_{D\times D}K_\tau(x,\dd x')K_\tau(y,\dd y')\,\pi(\dd x\dd y)
\]
couples $S_\tau\varrho_0\dd x$ and $S_\tau\varrho_1\dd x$.  Moreover
$|(y'-x')-(y-x)|_\infty\le2\tau$ on the support of the lifted coupling.
Consequently, for every bounded continuous displacement cost $c(y-x)$,
\[
 \int c(y'-x')\,\dd\pi_\tau\longrightarrow
 \int c(y-x)\,\dd\pi.
\]
If in addition $|y_2-x_2|\le r$ $\pi$-a.s., then
$|y'_2-x'_2|\le r+2\tau$ $\pi_\tau$-a.s.
\end{enumerate}
\end{lemma}

\begin{proof}
Choose an even $C_c^\infty(\R)$ probability density $\eta_\tau$ supported
in $(-\tau,\tau)$.  In one dimension define the reflected, two-periodic
kernel on $[0,1]$ by the method of images
\[
 k^{(1)}_\tau(x,z)
 =\sum_{n\in\mathbb Z}\Bigl[
   \eta_\tau(z-x+2n)+\eta_\tau(z+x+2n)
 \Bigr],
 \qquad x,z\in[0,1].
\]
It is nonnegative, symmetric, smooth, and its integral in either variable is
one.  Equivalently, it is convolution by $\eta_\tau$ after even
$2$-periodic reflection.  Since $\tau<1/4$, every nonzero image term has
$|x-z|\le\tau$; for the reflected terms near $0$ or $1$, this follows from
$x+z\le\tau$ or $2-x-z\le\tau$.  Take the product of the two
one-dimensional kernels to obtain $k_\tau$ on $D$.

The Markov and bistochastic identities imply preservation of mass and the
pointwise bounds $\underline\varrho,\overline\varrho$.  Standard approximate-identity theory for the even
periodic extension gives $L^1$ convergence and smoothness.  The formula for
$\pi_\tau$ gives the asserted marginals by Fubini.  Local support gives
$|x'-x|_\infty,|y'-y|_\infty\le\tau$, hence the displacement and hard-support
bounds.  Uniform continuity of a bounded continuous cost on the compact set
$D-D$ gives the final convergence.
\end{proof}

\begin{corollary}
\label{red:cor:transport-recovery}
Assume $0<\underline\varrho\le\varrho_i\le\overline\varrho$ and let $\varrho_i^\tau=S_\tau\varrho_i$.
If $r>\mathfrak T_{\boldsymbol r}(\varrho_0,\varrho_1)$ and
$r'>r$, then for all sufficiently small $\tau$ there is a coupling of
$\varrho_0^\tau$ and $\varrho_1^\tau$ whose finite-coordinate modulars at scale
$r'$ are bounded by a fixed constant independent of $\tau$.  At a hard
coordinate, a bound $|y_j-x_j|\le r$ is replaced by
$|y_j'-x_j'|\le r'$.  The fixed modular constant may be absorbed into a
fixed further enlargement of $r'$ by \cref{red:lem:absorb}.
\end{corollary}

\begin{proof}
Start from a coupling witnessing the gauge at scale $r$ and apply
\cref{red:lem:local-smoothing}.  For finite coordinates the costs at the larger
scale $r'$ are bounded continuous functions on $D-D$ and their integrals
converge to those of the original coupling, which are at most one at scale
$r$ and hence at most one at scale $r'$.  For a hard coordinate choose
$\tau<(r'-r)/2$.
\end{proof}

\subsection{Finite critical exponents}

For finite $r_1,r_2$, set
\[
 L_r(z)=\sum_{j=1}^2\psi_{r_j}(|z_j|/r).
\]
The generalized Benamou--Brenier formula for this convex Lagrangian is standard; see, for example, \cite{Santambrogio,Villani}:
\begin{equation}
 \inf_\pi\int L_r(y-x)\dd\pi
 =\inf_{\partial_t\varrho_t+\divv J_t=0}
 \int_0^1\int_D\varrho_tL_r(J_t/\varrho_t)\dd x\dd t.
 \label{red:eq:BB}
\end{equation}

\begin{theorem}
\label{red:thm:transport-flow-finite}
Assume $2<r_1,r_2<\infty$ and $1/r_1+1/r_2=1/2$, and let
$0<\underline\varrho\le\varrho_0,\varrho_1\le\overline\varrho$. Then
\begin{equation}
 C_{\underline\varrho,\overline\varrho}^{-1}\mathfrak F_{\boldsymbol r}^{D}(\varrho_0-\varrho_1)
 \le \mathfrak T_{\boldsymbol r}(\varrho_0,\varrho_1)
 \le C_{\underline\varrho,\overline\varrho}\mathfrak F_{\boldsymbol r}^{D}(\varrho_0-\varrho_1),
 \label{red:eq:finite-transport-flow}
\end{equation}
with a constant depending only on $\underline\varrho,\overline\varrho$ and the fixed exponents.
\end{theorem}

\begin{proof}
\emph{Flow to transport.}
Let $\divv J=\varrho_0-\varrho_1$, put
$\varrho_t=(1-t)\varrho_0+t\varrho_1$ and $J_t=J$, and let
$a>\norm J_{\boldsymbol r}$.  Since $\varrho_t\ge\underline\varrho$ and the perspective
$s\mapsto s\psi_{r_j}(q/s)$ is decreasing,
\[
 \int_0^1\!\int\varrho_t
 \psi_{r_j}\pa{\frac{|J_j|}{\varrho_t(a/\underline\varrho)}}
 \le\underline\varrho\int\psi_{r_j}(|J_j|/a)\le\underline\varrho.
\]
Thus the total $L_{a/\underline\varrho}$ action is at most $2\underline\varrho\le2$.  By
\eqref{red:eq:BB} and \cref{red:lem:absorb}, the two displacement norms are at most
$2a/\underline\varrho$.

\emph{Transport to flow.}
Take $r>\mathfrak T_{\boldsymbol r}(\varrho_0,\varrho_1)$ and then fix $r'>r$.
Let $\varrho_i^\tau=S_\tau\varrho_i$.  By
\cref{red:cor:transport-recovery}, for small $\tau$ there is a coupling of the
smooth densities $\varrho_0^\tau,\varrho_1^\tau$ whose two coordinate modulars
at scale $r'$ are bounded by a universal fixed constant.  After a fixed
enlargement of $r'$ using \cref{red:lem:absorb}, we may assume that their
sum is at most $2+o_\tau(1)$.

For $\varepsilon>0$, minimize
$L_{r'}(z)+\varepsilon|z|^2$ between the smoothed marginals.  The minimizer is
induced by a map and \cref{red:lem:density-interpolation} applies.  Its
interpolation densities obey
$\varrho_t^{\tau,\varepsilon}\le\varrho_*=\overline\varrho^2/\underline\varrho$, uniformly in both
$\tau$ and $\varepsilon$.  If
$v_t(T_t(x))=T(x)-x$ and $J_t=\varrho_tv_t$, set
$J^{\tau,\varepsilon}=\int_0^1J_t\dd t$ and
$a(x)=\int_0^1\varrho_t(x)\dd t\le\varrho_*$.  Coordinatewise Jensen gives
\[
 \varrho_*\psi_{r_j}\pa{\frac{|J_j^{\tau,\varepsilon}|}
 {\varrho_* r'}}
 \le\int_0^1\varrho_t\psi_{r_j}(|v_{t,j}|/r')\dd t.
\]
Hence the currents are bounded, uniformly in $\tau,\varepsilon$, in the
corresponding Orlicz spaces at scale $C_{\underline\varrho,\overline\varrho}r'$.

First let $\varepsilon\downarrow0$.  Then let $\tau\downarrow0$.  The
superlinear Young functions give uniform integrability of the current
coordinates, hence weak compactness in $L^1$; the divergence constraint
passes distributionally because
$\varrho_i^\tau\to\varrho_i$ in $L^1$.  Lower semicontinuity of convex integral
functionals gives a limiting current $J$ with
$\divv J=\varrho_0-\varrho_1$ and
$\norm{J_j}_{L^{\psi_{r_j}}}\le C_{\underline\varrho,\overline\varrho}r'$.  Finally let $r'\downarrow r$ and
$r\downarrow\mathfrak T_{\boldsymbol r}(\varrho_0,\varrho_1)$.

The smoothing step is where \cref{red:cor:transport-recovery} enters: it gives
an upper bound for the smoothed transport problem, whereas lower
semicontinuity alone gives the inequality in the wrong direction.
\end{proof}

\subsection{The critical endpoint $(2,\infty)$}

Assume, after exchanging coordinates if necessary, that
$\boldsymbol r=(2,\infty)$.

\begin{theorem}
\label{red:thm:transport-flow-endpoint}
If $0<\underline\varrho\le\varrho_0,\varrho_1\le\overline\varrho$, then
\begin{equation}
 \mathfrak T_{(2,\infty)}(\varrho_0,\varrho_1)
 \asymp_{\underline\varrho,\overline\varrho}
 \mathfrak F_{(2,\infty)}^{D}(\varrho_0-\varrho_1).
 \label{red:eq:endpoint-transport-flow}
\end{equation}
\end{theorem}

\begin{proof}
\emph{Flow to transport.}
Let $a>\max(\norm{J_1}_{\exp L^2},\norm{J_2}_\infty)$ and use the linear
density path.  Then $|J_2/\varrho_t|\le a/\underline\varrho$.  By the weak admissibility identity
\eqref{red:eq:weak-div-no-flux}, extending $\varrho_t$ and $J$ by zero outside $D$
produces a distributional continuity equation on $\R^2$ with no boundary
source.  The standard superposition principle for continuity equations, see
\cite{AGS}, then
produces a dynamic plan concentrated on trajectories in $\overline D$ whose
endpoints satisfy
$|y_2-x_2|\le a/\underline\varrho$.  Jensen along trajectories and the perspective estimate
control the first coordinate exponential-square modular at the same scale,
up to a fixed factor.

\emph{Transport to flow.}
Let $r$ exceed the endpoint transport gauge and fix $r'>r$.  Choose a
coupling $\pi_0$ with
\[
 |y_2-x_2|\le r\quad\pi_0\text{-a.s.},
 \qquad
 \int\psi_{r_1}(|y_1-x_1|/r)\dd\pi_0\le1.
\]
Smooth the two marginals and the coupling by
\cref{red:lem:local-smoothing}.  For all sufficiently small $\tau$ the
resulting coupling $\pi_0^\tau$ has smooth marginals
$\varrho_i^\tau=S_\tau\varrho_i$, still satisfying $\underline\varrho\le\varrho_i^\tau\le\overline\varrho$,
and
\[
 |y_2-x_2|\le r'\quad\pi_0^\tau\text{-a.s.},
 \qquad
 \int\psi_{r_1}(|y_1-x_1|/r')\dd\pi_0^\tau\le 1+o_\tau(1).
\]
Thus it is enough to obtain estimates uniform in $\tau$ for these smooth
marginals.

Use the explicit smooth convex barriers
\[
 b_k(s)=k\,(|s|-1)_+^4+k^{-1}s^2.
\]
They are $C^2$, strictly convex, satisfy
$\sup_{|s|\le1}b_k(s)\le k^{-1}$, and diverge uniformly on
$\{|s|\ge1+\eta\}$ for every $\eta>0$.  For fixed $\tau$, minimize the
strictly convex smooth costs
\[
 L_{r',k}(z)=\psi_{r_1}(|z_1|/r')+b_k(z_2/r').
\]
The smoothed benchmark gives a cost bounded uniformly in $k$ and $\tau$.
For every $k$, \cref{red:lem:density-interpolation} applies to the smooth
marginals and gives
$\varrho_t^{\tau,k}\le\varrho_*=\overline\varrho^2/\underline\varrho$.  Passing first $k\to\infty$ yields
a coupling $\pi^\tau$ satisfying
\[
 |y_2-x_2|\le r'\quad\pi^\tau\text{-a.s.},
 \qquad
 \int\psi_{r_1}(|y_1-x_1|/r')\dd\pi^\tau\le C,
\]
and the interpolation density bound passes to the limit by weak convergence
against nonnegative test functions.

Let $\varrho_t^\tau=(e_t)_\#\pi^\tau$ and
$J_t^\tau=(e_t)_\#((y-x)\pi^\tau)$.  Disintegrating over $e_t$ shows
\[
 |J_{t,2}^\tau|\le r'\varrho_t^\tau\le r'\varrho_*
 \quad\text{a.e.}
\]
Hence $J^\tau=\int_0^1J_t^\tau\dd t$ satisfies
$\norm{J_2^\tau}_\infty\le r'\varrho_*$, while the first component is
uniformly bounded in $\exp L^2$ by the same perspective/Jensen estimate as in
\cref{red:thm:transport-flow-finite}.  Extracting a subsequence as
$\tau\downarrow0$, use weak compactness in $L^1$ for the first component
and weak-* compactness in $L^\infty$ for the second.  The divergence equation
passes to the limit and convex lower semicontinuity yields a current $J$ with
\[
 \divv J=\varrho_0-\varrho_1,
 \qquad
 \norm{J_1}_{\exp L^2}+\norm{J_2}_\infty\le C_{\underline\varrho,\overline\varrho}r'.
\]
Finally let $r'\downarrow r$ and then $r\downarrow\mathfrak T_{(2,\infty)}$.
\end{proof}

The hard constraint enters this proof only in the limit $k\to\infty$; no
uniform $L^\infty$ bound is used for the approximating problems before the
convex barriers become infinite.

\section{Bounds on the histogram densities}

The transport--flow comparison of the previous section requires densities
bounded above and below. We check that the histograms satisfy this with
probability tending to one.

\begin{lemma}
\label{red:lem:hist-ellipticity}
For every fixed $\varepsilon\in(0,1)$, if $A$ in \eqref{red:eq:h-choice} is
sufficiently large, then
\[
 \Pp\pa{1-\varepsilon\le P_{h_N}^D\mu_N,
 P_{h_N}^D\nu_N\le1+\varepsilon\text{ on }D}\longrightarrow1.
\]
\end{lemma}

\begin{proof}
A terminal cell count has mean
$\lambda_N=Nh_N^2\asymp A^2\log N$. Chernoff's inequality and a union bound
over the $h_N^{-2}=N/\lambda_N$ cells give
\[
 \Pp\pa{\max_Q|N\mu_N(Q)-\lambda_N|>
 \varepsilon\lambda_N}
 \le\frac{2N}{\lambda_N}e^{-c\varepsilon^2\lambda_N},
\]
and similarly for $\nu_N$.
\end{proof}

Combining \cref{red:prop:atomic-hist,red:thm:transport-flow-finite,red:thm:transport-flow-endpoint,red:lem:hist-ellipticity} proves the equivalence of
statements (i) and (ii) in \cref{red:thm:main}.

\section{Coupling the histograms with dyadic Gaussian noise}

It remains to compare the histogram source with the projected Gaussian
source. We use the one-dimensional empirical-process approximation of
Koml\'os, Major and Tusn\'ady \cite{KMT1975}, in the nonasymptotic
form of Bretagnolle and Massart \cite[Theorem~1]{BretagnolleMassart1989};
see also \cite[Theorem~1.8 and equation~(1.1)]{DudleyUCLT}.
For the empirical distribution function $F_N$ of independent uniform
variables on $[0,1]$, it provides a coupled standard Brownian bridge
$B$ such that, with
\[
 \Delta_N=\sup_{0\le t\le1}|F_N(t)-t-N^{-1/2}B(t)|,
\]
one has $\Pp\{\Delta_N>(12\log N+x)/N\}\le2e^{-x/6}$ for all $x\ge0$.
In particular there are absolute positive constants $C,c$ such that
\begin{equation}\label{red:eq:KMT-cdf}
 \Pp\{\Delta_N>C(\log N+x)/N\}\le Ce^{-cx},
 \qquad N\ge2,\quad x>0.
\end{equation}
We encode the fixed planar partition by intervals and apply this
theorem directly. Since we only compare bounds with fixed constants, a
correction of size $O(s_N)$ with probability tending to one suffices.

\paragraph{A deterministic divergence inverse.}
Resolve the fixed $h=2^{-m}$ square grid by alternating vertical
and horizontal equal binary splits, for $2m$ generations. Each
level-$k$ rectangle has area $2^{-k}$ and aspect ratio between
$1/2$ and $2$. For an internal rectangle $Q$ with children
$Q^0,Q^1$, put
\[
 \mathsf h_Q=|Q|^{-1/2}(\one_{Q^0}-\one_{Q^1}).
\]
For a vertical split $Q=I\times K$, let $v_I$ vanish at both
endpoints of $I$ and satisfy $v_I'=\one_{I^0}-\one_{I^1}$.
Set $V_Q=|Q|^{-1/2}v_I(x)\one_K(y)e_1$, extended by zero;
for horizontal splits use the analogous $e_2$ field. Integration
by parts in the active coordinate gives
\begin{equation}\label{red:eq:corrector}
 \int_D V_Q\cdot\nabla\phi=-\int_D\mathsf h_Q\phi,
 \qquad \norm{V_Q}_\infty\le C,
 \qquad \phi\in C^1(\overline D).
\end{equation}
There is no boundary contribution: on faces normal to the active
component the primitive vanishes, and on the other faces it is tangential.

If $g$ is constant on the terminal squares and $\int_Dg=0$, set
$d_Q=\int_{Q^0}g-\int_{Q^1}g$. The tree Haar basis gives
\[
 g=\sum_Q\frac{d_Q}{\sqrt{|Q|}}\mathsf h_Q,
 \qquad R=\sum_Q\frac{d_Q}{\sqrt{|Q|}}V_Q.
\]
Thus $\divv R=g$ with zero normal flux in the weak sense of
\eqref{red:eq:corrector}. At almost every point only one rectangle
per generation contributes. Consequently
\begin{equation}\label{red:eq:tree-discrepancy-inverse}
 \norm R_\infty\le
 C\max_Q|d_Q|\sum_{k=0}^{2m-1}2^{k/2}
 \le \frac Ch\max_Q|d_Q|.
\end{equation}
All sums are finite. For $m=0$ the mean-zero grid-constant source is zero.

\paragraph{The one-cloud coupling.}
Order the terminal squares recursively, placing all descendants
of $Q^0$ before those of $Q^1$, and assign them consecutive intervals
of length $h^2$ in $[0,1]$. Every tree rectangle $Q$ then corresponds
to a single interval $I_Q$ of length $|Q|$.
Take the uniform sample in \eqref{red:eq:KMT-cdf}. For each sample
variable in a leaf interval $I_S$, independently choose a uniform
point in $S$. These points are independent and uniform in $D$,
and their counts on tree rectangles equal the interval counts.
The construction does not condition on the event where the KMT
bound holds. Fixed interval endpoints are avoided almost surely.

For $I_S=(a_S,b_S]$, define $W_h^X$ on $S$ by
\[
 \int_S W_h^X=B(b_S)-B(a_S).
\]
These masses are centered jointly Gaussian and satisfy
\[
 \operatorname{Cov}\left(\int_S W_h^X,\int_T W_h^X\right)
 =|S|\one_{S=T}-|S||T|.
\]
Hence $W_h^X$ has exactly the law of $P_h^DW^D$, including the
mean-zero constraint. For
$g=P_h^D\mu_N-1-N^{-1/2}W_h^X$, interval increments of the
error in \eqref{red:eq:KMT-cdf} give
$|\int_Qg|\le2\Delta_N$ and $|d_Q|\le4\Delta_N$.
The deterministic inverse therefore supplies a current $R_h^X$ with
\begin{equation}\label{red:eq:one-cloud-KMT}
 \divv R_h^X=P_h^D\mu_N-1-N^{-1/2}W_h^X,
 \qquad
 \Pp\!\left\{\norm{R_h^X}_\infty>
 C\frac{\log N+x}{Nh}\right\}\le Ce^{-cx}.
\end{equation}
Its normal flux is zero. The joint construction may depend on
$N$ and $h$; this is enough for the probability comparisons below.

\begin{theorem}\label{red:thm:KMT}
If $Nh^2\asymp\log N$, there is a coupling of two independent
uniform point clouds and a field $W_h^D$ having exactly the law
of $P_h^DW^D$, together with a zero-normal-flux current $Q_h$, such that
\begin{equation}\label{red:eq:KMT-div}
 \divv Q_h=P_h^D(\mu_N-\nu_N)-\sqrt{2/N}\,W_h^D.
\end{equation}
For a deterministic $C<\infty$,
\begin{equation}\label{red:eq:KMT-infty}
 \Pp\!\left\{\norm{Q_h}_\infty\le
 C\frac{\log N}{Nh}\right\}\longrightarrow1.
\end{equation}
In particular $\norm{Q_h}_{\boldsymbol r}\le Cs_N$ with probability
tending to one for every critical exponent pair, including the endpoints.
\end{theorem}

\begin{proof}
Use independent copies of \eqref{red:eq:one-cloud-KMT} for the two
clouds and put
\[
 W_h^D=(W_h^X-W_h^Y)/\sqrt2,
 \qquad Q_h=R_h^X-R_h^Y.
\]
The Gaussian law and the divergence identity follow directly.
Take $x$ to be a sufficiently large fixed multiple of $\log N$
in \eqref{red:eq:one-cloud-KMT} and use a union bound.
Since $Nh^2\asymp\log N$, we have $\log N/(Nh)\asymp s_N$.
Finally, on the unit-area square,
$\norm f_{\psi_r}\le(\log2)^{-1/r}\norm f_\infty$
for $2\le r<\infty$, and $\norm f_{\psi_\infty}=\norm f_\infty$.
\end{proof}
The correction $Q_h$ and the Gaussian source $W_h^D$ are in general
dependent; the comparison below does not use independence.


\section{Proof of the reduction}\label{sec:probability-transfer}

We now assemble the previous sections into the proof of
\cref{red:thm:main}, and then derive the one-cloud version of the
matching estimate.

\begin{proof}[Proof of \cref{red:thm:main}]
By \cref{red:prop:atomic-hist}, atomic matching and histogram transport are
equivalent at scale $s_N$, because $h_N\asymp s_N$. On the event of
\cref{red:lem:hist-ellipticity}, the transport--flow comparison
\cref{red:thm:transport-flow-finite,red:thm:transport-flow-endpoint}
identifies histogram transport with the corresponding static minimum-flow
gauge, for finite and endpoint exponents respectively. This proves
(i)$\Leftrightarrow$(ii).

By \cref{red:lem:flow-lipschitz,red:thm:KMT} and the homogeneity of
the minimum-flow gauge,
\[
 \left|\mathfrak F_{\boldsymbol r}^{D}
       (P_{h_N}^D(\mu_N-\nu_N))
       -\sqrt{\frac2N}\,\mathfrak F_{\boldsymbol r}^{D}(W_{h_N}^{D})
 \right|\le\norm{Q_{h_N}}_{\boldsymbol r}.
\]
By \eqref{red:eq:KMT-infty}, the right-hand side is at most $Cs_N$ with
probability tending to one for every critical pair. Intersecting
this event with a high-probability bound for either flow gauge
transfers that bound to the other after changing its constant;
no independence is required. Since
$\sqrt{2/N}\,\sqrt{\log N}=\sqrt2\,s_N$,
this proves (ii)$\Leftrightarrow$(iii).
\end{proof}


\subsection{One cloud and Talagrand's evenly spread references}

The same argument gives the original one-cloud formulation of the
Ultimate Matching Conjecture, and not only its two-sample analogue.
It uses no additional Gaussian input.

\begin{corollary}
\label{red:cor:evenly-spread}
Fix a critical exponent pair $\boldsymbol r$ as in
\eqref{red:eq:critical-pair} and let $h_N$ satisfy
\eqref{red:eq:h-choice}. Suppose that for some deterministic $C$,
\begin{equation}\label{red:eq:one-cloud-input}
 \Pp\!\left(\mathfrak F_{\boldsymbol r}^D(W_{h_N}^D)
       \le C\sqrt{\log N}\right)\longrightarrow1.
\end{equation}
Let $Y_1,\ldots,Y_N$ be any deterministic evenly spread family in
$D$: there are rectangles $R_i$ partitioning $D$ up to null sets,
with $|R_i|=1/N$, $Y_i\in R_i$, and
$R_i\subset B(Y_i,10/\sqrt N)$. Then for some deterministic
$C'<\infty$, independent of the chosen family,
\begin{equation}\label{red:eq:one-cloud-matching}
 \Pp\!\left(\min_{\sigma\in S_N}\max_{j=1,2}
 \bigl\|(X_i^j-Y_{\sigma(i)}^j)_{i=1}^N\bigr\|_{\psi_{r_j},N}
 \le C'\sqrt{\frac{\log N}{N}}\right)\longrightarrow1.
\end{equation}
\end{corollary}

The conclusion includes a hard coordinate endpoint whenever the
corresponding Gaussian hypothesis \eqref{red:eq:one-cloud-input}
holds; the corollary does not supply that hypothesis.

\begin{proof}
Let $\varrho_h=P_{h_N}^D\mu_N$. The one-cloud coupling
\eqref{red:eq:one-cloud-KMT} gives
\[
 \varrho_h-1=N^{-1/2}W_{h_N}^X+\divv R_{h_N}^X,
\]
where $W_{h_N}^X$ has the law of $W_{h_N}^D$ and
$\norm{R_{h_N}^X}_\infty\le Cs_N$ with probability tending to one.
The same bound holds in the coordinate Orlicz gauge. Therefore
\eqref{red:eq:one-cloud-input} and the flow-gauge triangle inequality give
\[
 \mathfrak F_{\boldsymbol r}^D(\varrho_h-1)
 \le N^{-1/2}\mathfrak F_{\boldsymbol r}^D(W_{h_N}^X)
       +\norm{R_{h_N}^X}_{\boldsymbol r}
 \le C_1s_N
\]
with probability tending to one. By
\cref{red:lem:hist-ellipticity} applied to the first cloud,
$1/2\le\varrho_h\le3/2$ on another event of probability tending
to one. Hence the transport--flow comparison,
\cref{red:thm:transport-flow-finite,red:thm:transport-flow-endpoint},
yields a coupling from $\varrho_h\dd x$ to $\dd x$ whose two
coordinate Luxemburg norms are bounded by $C_2s_N$.

There is also a coupling from $\mu_N$ to $\varrho_h\dd x$ which
moves each atom only inside its $h_N$-cell, with coordinate
displacement at most $h_N$. Finally, the deterministic map which
sends $x\in R_i$ to $Y_i$ couples Lebesgue measure to
$N^{-1}\sum_i\delta_{Y_i}$, moving every point by at most
$10/\sqrt N$. Gluing these three couplings and using the
Luxemburg triangle inequality gives a coupling $\gamma$ between
the two equal-weight atomic measures for which
\[
 \max_j\|x_j-y_j\|_{L^{\psi_{r_j}}(\gamma)}
 \le C_2s_N+h_N+10/\sqrt N\le C_3s_N.
\]

It remains to replace $\gamma$ by a permutation. The coupling $\gamma$
determines a doubly stochastic matrix, so by Birkhoff's theorem it is a
convex combination of matching permutations. At a scale slightly larger
than $C_3s_N$, the mean over these permutations of the sum of the two
coordinate Young-function modulars is at most $2$. Therefore some
single permutation has this sum at most $2$, and a fixed further
enlargement of scale makes each coordinate modular at most $1$.
For $r_j=\infty$, all positive entries of the doubly stochastic
matrix already obey the hard coordinate bound; we select a
permutation in that support while controlling the remaining
finite-coordinate modular. This yields \eqref{red:eq:one-cloud-matching}
with a constant uniform over the evenly spread family.
\end{proof}

\section{From the torus to the square}
\label{sec:fold}
The reduction requires a Gaussian source on the square, whereas the
flow construction of Part~\ref{umc:partI} yields a periodic current on
$\T^2=(\mathbb R/\mathbb Z)^2$. In this section we show that folding the
torus onto the square relates the two projected sources exactly, so that
a torus bound implies the corresponding square bound
(\cref{prop:torus-implies-square}).

\subsection{Folding of currents}
On $\T^1=[0,1)$ define
\[
 F_1(u)=2\min\{u,1-u\},\qquad
 F=(F_1,F_1):\T^2\longrightarrow D=[0,1]^2.
\]
The notation $F=(F_1,F_1)$ means $F(u_1,u_2)=(F_1(u_1),F_1(u_2))$.
Outside the folds $u_i=0,1/2$, the derivative is
$DF(u)=2R_u$, where $R_u$ is diagonal with entries $\pm1$. Each
$y\in(0,1)^2$ has exactly four preimages.
For integrable scalar functions and vector fields define
\begin{equation}
 (\mathscr S_0 f)(y)=\frac14\sum_{F(u)=y}f(u),\qquad
 (\mathscr S_1 J)(y)=\frac12\sum_{F(u)=y}R_u J(u).
 \label{eq:unit-fold}
\end{equation}
These formulas are the pushforward against normalized Lebesgue measure;
in particular $F_\#\mathrm du=\mathrm dy$.
For mean-zero distributions the scalar pushforward is defined by
$\langle\mathscr S_0 f,\varphi\rangle
=\langle f,\varphi\circ F\rangle$.

\begin{lemma}
\label{lem:unit-fold-current}
Suppose $J\in L^1(\T^2;\mathbb R^2)$ and
$\operatorname{div}_{\T}J=f$ in distributions. Then
$J^D=\mathscr S_1J$ satisfies
\[
 -\int_D J^D\cdot\nabla\varphi\,\dd y
 =\langle\mathscr S_0f,\varphi\rangle
 \quad\text{for every }\varphi\in C^1(\overline D).
\]
In particular, $J^D$ has zero normal trace on $\partial D$.
For every Young function $\Theta$ and $j=1,2$,
\begin{equation}
 \int_D\Theta\left(\frac{|J_j^D|}{2a}\right)\dd y
 \le\int_{\T^2}\Theta\left(\frac{|J_j|}{a}\right)\dd u,
 \qquad
 \|J_j^D\|_{L^\Theta(D)}\le2\|J_j\|_{L^\Theta(\T^2)}.
 \label{eq:unit-fold-norm}
\end{equation}
The same inequality holds when $L^\Theta=L^\infty$.
\end{lemma}
\begin{proof}
On every inverse sheet $\dd u=\dd y/4$, and the chain rule gives
$\nabla(\varphi\circ F)(u)=2R_u\nabla\varphi(F(u))$. Thus
\[
 -\int_{\T^2}J\cdot\nabla(\varphi\circ F)\,\dd u
 =-\int_D\left(\frac12\sum_{F(u)=y}R_uJ(u)\right)
 \cdot\nabla\varphi(y)\,\dd y.
\]
For an $L^1$ current, testing with the periodic Lipschitz function
$\varphi\circ F$ is justified by smooth approximation; in our application
the source is a finite-dimensional Gaussian cell function. The
distributional divergence identity on the torus gives the first claim.
Since the test functions need not vanish on $\partial D$, this identity
also shows that the fold creates no boundary flux. Finally
$J_j^D=2\cdot(1/4)\sum\varepsilon_{u,j}J_j(u)$, and Jensen's inequality
together with $F_\#\mathrm du=\mathrm dy$ proves \eqref{eq:unit-fold-norm}.
\end{proof}

\subsection{Compatibility with dyadic projections}
Recall that $P_h^{\T}$ and $P_h^D$ are the fixed-grid conditional
expectations. Restricted to $L_0^2$, their ranges are the mean-zero
cellwise-constant subspaces. Accordingly the projected noises
$W_h^{\T},W_h^D$ have exactly the meaning specified in
\eqref{red:eq:compatible-source}.
If $h=2^{-m}$ is a square mesh, take the torus mesh
\[
 \widetilde h=h/2=2^{-(m+1)}.
\]
The four folding lines coincide with grid lines at this mesh. Thus
every torus $\widetilde h$-square is mapped bijectively onto a square
$h$-square, and the following deterministic relation holds up to
cell-boundary null sets:
\begin{equation}
 \mathscr S_0P_{\widetilde h}^{\T}f
 =P_h^{D}\mathscr S_0f.
 \label{eq:commute-projections}
\end{equation}

\begin{lemma}
\label{lem:noise-fold}
Let $W^{\T}$ be centered mean-zero isonormal white noise over
$L_0^2(\T^2)$. For $h=2^{-m}$ and $\widetilde h=h/2$,
\begin{equation}
 \mathscr S_0W_{\widetilde h}^{\T}
 \ \stackrel{\mathrm{law}}=\ P_h^{D}W^D,
 \label{eq:exact-fold-law}
\end{equation}
where $W^D$ is standard mean-zero white noise on $D$.
\end{lemma}
\begin{proof}
For deterministic $\phi,\psi\in L_0^2(D)$, the isonormal
covariance of the folded field is
\begin{align*}
 \mathbb E\bigl[\langle\mathscr S_0W_{\widetilde h}^{\T},
 \phi\rangle\langle\mathscr S_0W_{\widetilde h}^{\T},
 \psi\rangle\bigr]
 &=\langle P_{\widetilde h}^{\T}(\phi\circ F),
 P_{\widetilde h}^{\T}(\psi\circ F)\rangle_{L^2(\T^2)}\\
 &=\langle P_h^{D}\phi,P_h^{D}\psi\rangle_{L^2(D)}.
\end{align*}
The second equality holds because $F$ is measure preserving and the
reflected dyadic mesh is exactly aligned. Both fields are centered
Gaussian on the same finite-dimensional cell space, so equality of
covariances proves equality in law.
\end{proof}

\begin{proposition}
\label{prop:torus-implies-square}
For every pair $\boldsymbol r$ for which the coordinate Orlicz
spaces are defined, and every $m\ge1$,
\[
 \mathbb P\left\{\mathfrak F_{\boldsymbol r}^{D}
 (P_{2^{-m}}^{D}W^D)\le2t\right\}
 \ge
 \mathbb P\left\{\mathfrak F_{\boldsymbol r}^{\T}
 (P_{2^{-(m+1)}}^{\T}W^{\T})\le t\right\},
 \qquad t>0.
\]
Consequently, a bound $O(\sqrt m)$ with high probability for the
unit-torus Gaussian flow implies the identical exact-cell-source
square Gaussian estimate, with at most a universal change of constant.
\end{proposition}
\begin{proof}
For any torus sample and any admissible periodic current, use
\cref{lem:unit-fold-current} to obtain a no-flux square current for the
folded source, at coordinate cost at most twice the torus cost.
Take infima (or approximating sequences) and then use
\cref{lem:noise-fold} to identify the source law.
\end{proof}

\begin{remark}
A general no-flux square current cannot be pulled back to the torus at
unchanged Gaussian source law simply by reversing the four-sheeted
averaging: the fold destroys information from three independent
complementary sheet modes. The matching--square-flow equivalence is
symmetric, but the folding result is used only in the
torus-to-square direction.
\end{remark}

\section{Choice of dyadic mesh}
\label{sec:alignment}
The torus Gaussian bound is formulated along the meshes $h_m=2^{-m}$ and
expressed in $\sqrt m$, whereas the matching reduction uses the mesh $h_N$
of \eqref{red:eq:h-choice}. We check that these two parametrizations
describe the same family of dyadic projections, and then combine the
previous sections into a matching criterion stated on the torus.

In the square matching reduction one chooses a dyadic $h_N$ with
$A\sqrt{\log N/N}\le h_N<2A\sqrt{\log N/N}$. Given
$N$, write $h_N=2^{-m(N)}$, so that
\[
 m(N)=\frac{\log N}{2\log2}+O(\log\log N),\qquad
 \sqrt{m(N)}\asymp\sqrt{\log N}.
\]
The square-source estimate follows from the torus theorem at the single
adjacent depth $m(N)+1$. Conversely, a square-source high-probability
bound formulated at every $h_N$ gives a bound along every $m$ by choosing
integers $N_m$ with
\[
 N_m\asymp A^2m4^m,
 \qquad
 A\sqrt{\frac{\log N_m}{N_m}}\le2^{-m}
 <2A\sqrt{\frac{\log N_m}{N_m}}.
\]
For instance, one may take
$N_m=\lceil4A^2(\log 2)m4^m\rceil$: then
$A\sqrt{(\log N_m)/N_m}/2^{-m}\to1/\sqrt2$,
which lies strictly between $1/2$ and $1$.
This reverse calibration concerns the square
Gaussian gauges; it does not reverse the torus fold.

\begin{corollary}
\label{cor:big-picture}
Fix $\boldsymbol r$ satisfying \eqref{red:eq:critical-pair}.
Suppose that there exist constants $B,p_0>0$ such that, for every
sufficiently large $m$,
\[
 \mathbb P\left\{\mathfrak F_{\boldsymbol r}^{\T}
 (P_{2^{-m}}^{\T}W^{\T})\le B\sqrt m\right\}
 \ge p_0.
\]
Then Talagrand's square matching estimate follows:
\[
 \mathbb P\left\{\mathfrak M_N^{\boldsymbol r}
 \le C'\sqrt{\log N/N}\right\}\longrightarrow1
\]
for a deterministic $C'<\infty$.
\end{corollary}


\begin{proof}
\Cref{red:prop:gaussian-concentration}, applied on $\T^2$, upgrades
the assumed lower probability bound to an exponentially small failure
probability at a possibly larger constant. Choose the square mesh
$h_N$ as in \cref{sec:alignment} and apply
\cref{prop:torus-implies-square} at the adjacent torus mesh $h_N/2$.
This gives the square Gaussian gauge estimate with high probability.
The three-way matching--histogram--square-Gaussian equivalence
\cref{red:thm:main} finishes the proof, since
$\sqrt{2/N}\,\sqrt{\log N}=\sqrt2s_N$.
\end{proof}

Thus the fixed-positive-probability estimate of \cref{thm:gauss}
suffices for the finite critical matching estimates.
Part~\ref{umc:partIII} shows that this torus criterion cannot hold at
the hard endpoint $(2,\infty)$, since its square consequence is false.

\section{Consequences and the hard endpoint}
\label{red:sec:consequences}

We conclude Part~\ref{umc:partII} by recording what the reduction gives
for the quartic pair and for the hard endpoint, which is the subject of
Part~\ref{umc:partIII}.
The reduction of \cref{red:thm:main} applies to every critical exponent pair,
including both hard-coordinate endpoints. The different outcomes for finite
and infinite exponents concern the Gaussian minimum-flow estimate itself,
not the validity of the matching--Gaussian comparison.

\paragraph{The quartic case.}
For $\boldsymbol r=(4,4)$, the isotropic comparison
\cref{red:lem:isotropic-quartic} identifies the matching statement, up to
universal changes of constants, with the square Gaussian estimate
\[
 \mathbb P\!\left\{
 \inf_{\substack{\operatorname{div}J=W_{h_N}^{D}\\J\cdot n=0}}
    \|J\|_{\exp L^4(D)}\le C\sqrt{\log N}\right\}
 \longrightarrow 1.
\]
Thus \cref{thm:gauss}, combined with \cref{cor:big-picture}, resolves the
finite critical pairs, including the symmetric quartic case.

\paragraph{The hard endpoint.}
For $\boldsymbol r=(2,\infty)$ the corresponding square gauge is
\[
 \mathfrak F_{(2,\infty)}^D(W_{h_N}^{D})
 =\inf_{\substack{\operatorname{div}J=W_{h_N}^{D}\\J\cdot n=0}}
 \max\{\|J_1\|_{\exp L^2(D)},\|J_2\|_{\infty}\}.
\]
The comparison theorem remains applicable, including the $L^\infty$
coordinate: the endpoint KMT correction is of the target order
and can be absorbed by changing the fixed constant. However,
\cref{neg:thm:Gaussian} in Part~\ref{umc:partIII} shows that, on the square,
\[
  \frac{\mathfrak F_{(2,\infty)}^D(W_{2^{-m}}^D)}{\sqrt m}
       \xrightarrow{\mathbb P}+\infty.
\]
Consequently the hard-endpoint matching assertion is false
(\cref{neg:thm:main}). The one-way torus-to-square implication of
\cref{prop:torus-implies-square} also rules out a
fixed-positive-probability torus upper bound at this endpoint.

\paragraph{Dual tests and the cutoff.}
The complementary Young function to $\psi_r(t)=e^{|t|^r}-1$ grows,
for large arguments, as $q(\log q)^{1/r}$; near zero it grows as
$q^{r/(r-1)}$. At $r=2$, the admissible dual gradient cost is therefore
comparable to $|q|\sqrt{\log(2+|q|)}$, while the associate coordinate
for $L^\infty$ is $L^1$. The lower-bound argument of Part~\ref{umc:partIII}
only needs a weak Orlicz--H\"older inequality for smooth potentials
(\cref{neg:lem:dual}); it does not use an identification of the full
Banach dual of an exponential Orlicz space. Such an identification would
require care, because $\psi_r$ does not satisfy the $\Delta_2$ condition
at infinity.

The dyadic projection is necessary: unregularized planar white noise
is almost surely not in $H^{-1}$, so it has no static $L^2$ divergence
current. The specific matching mesh $h_N\asymp\sqrt{\log N/N}$,
however, is fixed by the occupancy and histogram-regularization
requirements of this part, not by the ultraviolet divergence alone.


\clearpage
\part{Failure of the hard endpoint}\label{umc:partIII}

In this part we prove that at the hard coordinate endpoint $(2,\infty)$,
and by symmetry at $(\infty,2)$, the optimal matching cost divided by
$\sqrt{\log N/N}$ diverges in probability (\cref{neg:thm:main}).
The obstruction we show here is qualitative, but a determination of the
exact divergence rate of the cost can be found in the GitHub
repository~\cite{umc:repository} associated to this file.

\section{Statement, mechanism, and Gaussian reduction}
\label{neg:sec:introduction}

This section states the result, explains the mechanism, and reduces
the theorem to a statement about the projected Gaussian source on the
square. This part uses its own notation; where an object
coincides with one of Part~\ref{umc:partII}, we say so.

Let $D=[0,1]^2$, let $(X_i)_{i=1}^N$ and $(Y_i)_{i=1}^N$ be
independent uniform samples in $D$, and write
\[
 s_N=\sqrt{\frac{\log N}{N}},\qquad
 \psi_2(t)=e^{t^2}-1.
\]
For a permutation $\pi$, denote by $\Delta_j^\pi$
the $N$-vector of $j$th coordinate displacements
$X_i^j-Y_{\pi(i)}^j$, with the normalized Luxemburg norm
$\|\cdot\|_{\psi_2,N}$. Put
\begin{equation}\label{neg:eq:matching-cost}
 M_N=\min_{\pi\in S_N}
 \max\left\{\|\Delta_1^\pi\|_{\psi_2,N},
              \max_i|\Delta_2^\pi(i)|\right\}.
\end{equation}

\begin{theorem}
\label{neg:thm:main}
For the two independent point clouds,
\begin{equation}\label{neg:eq:negative-result}
  \frac{M_N}{s_N}\xrightarrow[N\to\infty]{\mathbb P}+\infty.
\end{equation}
The assertion holds after exchanging the coordinate axes. In fact,
the following one-cloud conclusion is uniform over \emph{all}
deterministic reference configurations $z^{(N)}=(z_i^{(N)})_{i=1}^N\in D^N$,
including the evenly spread configurations of
\cref{red:cor:evenly-spread}. For every fixed
$A<\infty$,
\begin{equation}\label{neg:eq:reference-negative}
 \sup_{z^{(N)}\in D^N}\mathbb P\!\left\{
   \min_{\pi\in S_N}\max\left(
     \|(X_i^1-z_{\pi(i)}^{(N),1})_i\|_{\psi_2,N},
     \max_i|X_i^2-z_{\pi(i)}^{(N),2}|
   \right)\le A s_N\right\}\longrightarrow0.
\end{equation}
\end{theorem}

\paragraph{Idea of the proof.}
Part~\ref{umc:partI} establishes \emph{upper} Gaussian-flow bounds for
finite exponents, which Part~\ref{umc:partII} transfers to matching. Here
we work in the opposite direction: we build random test potentials of
bounded dual cost whose expected pairing with the projected noise is
too large for any fixed endpoint bound.
Three probability spaces are involved. Brownian motion supplies
a model martingale; an auxiliary uniform point in the square samples
the gradients of a deterministic tree of triangles; finally, Gaussian
white noise selects the signs of the bubbles on that tree. The
spatial-branch law holds for \emph{every} assignment of signs,
whereas Gaussian freshness produces a positive expectation over those
signs. The two conditionings must be kept separate.

\paragraph{Why the endpoint is different.}
The complementary Young function to $\exp(t^2)-1$ grows like
$|v|\sqrt{\log(2+|v|)}$. Thus a dual gradient potential $u$ is charged
according to
\begin{equation}\label{neg:eq:PhiPsi}
 \Phi(a)=\int_0^{|a|}\sqrt{\log(2+r)}\,\mathrm dr,
 \qquad \Psi(a,b)=\Phi(a)+|b|.
\end{equation}
We seek an upper-triangular predictable control for the Brownian martingale
$\int_0^1L_t\,\mathrm dB_t$ with bounded terminal $\Psi$-cost but large
\begin{equation}\label{neg:eq:detreward}
   \mathbb E\int_0^1\sqrt{\det L_t}\,\mathrm dt.
\end{equation}
In an upper-triangular parametrization, $Y_t>0$ defines a change of
probability measure, and the ratio $M_t=X_t/Y_t$ is a centered
martingale under the $Y_t$-biased law. A sequence of short
local-time stages generates a large accumulated reward while keeping
$|M_t|$ small when the logarithmic weight on $Y_t$ is large. The
terminal cost therefore remains uniformly bounded even as the reward
increases. The obstruction is thus structural; it is not the failure
of a particular candidate flow.

Spatially, a determinant-one chart converts $\sqrt{\det L_t}$ into the
$L^2$ size of a bubble. Choosing its sign from fresh Gaussian noise
turns this size into a positive expected pairing. Unlike a deterministic
test, the adaptive potentials exploit this reward without increasing
the convex terminal gradient cost.

\paragraph{Plan of the proof.}
The argument has four steps: Brownian reward (\cref{neg:lem:Brownian}),
spatial realization (\cref{neg:lem:spatial}), a weak dual lower bound
(\cref{neg:lem:dual}), and concentration together with the matching
transfer (\cref{neg:thm:Gaussian}). The Brownian and spatial steps are
new. From Part~\ref{umc:partII} we use the square-flow Lipschitz
estimate \cref{red:lem:uniform-gaussian-lipschitz}, the triangle
inequality of \cref{red:lem:flow-lipschitz}, the mesh choice
\eqref{red:eq:h-choice}, and the probability-level comparisons of \cref{red:prop:atomic-hist},
\cref{red:thm:transport-flow-endpoint}, \cref{red:lem:hist-ellipticity}
and \cref{red:thm:KMT}.
The argument takes place directly on the square: the torus-to-square
folding of \cref{sec:fold} goes in one direction only, and it is not
reversed here.

\paragraph{The square Gaussian source.}
Let $\mathcal W$ be isonormal white noise on $L^2(D)$,
with mean-zero restriction $\mathcal W_0$. Let $P_h^D$ be cellwise
averaging on the fixed dyadic square grid of mesh $h=2^{-m}$, and put
$W_h=P_h^D\mathcal W_0$. Thus $W_h$ is a random \emph{function},
not a continuum white-noise distribution; its pairing with a test
$u$ equals $\mathcal W(P_h^Du)$ even for a random $u$ (interpreted
through the finite cell coefficients).
For a mean-zero distribution $f$, define the square, no-flux
minimum-flow gauge
\begin{equation}\label{neg:eq:flow-gauge}
 F_h(f)=\inf_{\substack{J=(J_1,J_2)\\
             -\int_DJ\cdot\nabla\varphi=\langle f,\varphi\rangle\ \forall\varphi\in C^\infty(\overline D)}}
       \max\{\|J_1\|_{L^{\psi_2}(D)},\|J_2\|_{L^\infty(D)}\},
 \qquad \mathcal F_m=F_{2^{-m}}(W_{2^{-m}}).
\end{equation}
The subscript $h$ merely indicates the mesh on which $f$ is supported;
admissible currents are continuum fields. Thus $F_h$ is the gauge
$\mathfrak F_{(2,\infty)}^{D}$ of \eqref{red:eq:aniso-flow}, and $W_h$
is the square source $W_h^D$ of \eqref{red:eq:compatible-source}.
The Gaussian coefficient representation is
\begin{equation}\label{neg:eq:coefficients}
  W_h|_Q=h^{-1}(Z_Q-\overline Z),\qquad
  \overline Z=h^2\sum_Q Z_Q,
\end{equation}
with independent standard normal $Z_Q$, as in
\eqref{red:eq:gaussian-coefficients}. By
\cref{red:lem:uniform-gaussian-lipschitz}, the coefficient-space gauge
$Z\mapsto\mathcal F_m(Z)$ is $C_{\mathrm{Lip}}$-Lipschitz with
$C_{\mathrm{Lip}}$ independent of $m$. We shall prove
\begin{equation}\label{neg:eq:gaussian-negative}
  \frac{\mathcal F_m}{\sqrt m}
     \xrightarrow[m\to\infty]{\mathbb P}+\infty
\end{equation}
for this square source, and then deduce \cref{neg:thm:main} in
\cref{neg:sec:conclusion} from the eventwise comparisons of
Part~\ref{umc:partII}.

\section{A pathwise dual inequality}
\label{neg:sec:duality}

The test potentials of \cref{neg:sec:spatial} are chosen from the
noise itself. The following inequality bounds their pairing with the
source by the flow gauge, pathwise, so that no independence between
test and noise is needed.

\begin{lemma}\label{neg:lem:dual}
There is an absolute constant $C_D$ such that, for every dyadic $h$,
for every mean-zero Lipschitz $u:D\to\mathbb R$, and for every realized
finite-dimensional source $W_h$,
\begin{equation}\label{neg:eq:dual-bound}
 |\mathcal W(P_h^Du)|
 \le C_D\,F_h(W_h)
   \left(1+\int_D\Psi(\partial_1u,\partial_2u)\,\mathrm dx\right).
\end{equation}
The assertion holds pathwise even when the test $u$ is selected
measurably from the same noise $\mathcal W$.
\end{lemma}
\begin{proof}
The complementary Young bound, with an absolute $C$, is
\begin{equation}\label{neg:eq:Young}
  |ab|\le e^{a^2}-1+C\Phi(b)\qquad (a,b\in\mathbb R).
\end{equation}
Indeed, if $|a|\le2\sqrt{\log(2+|b|)}$, the term
$C\Phi(b)$ pays for $|ab|$. Otherwise $|b|\le e^{a^2/4}$ and
$|a|\ge2\sqrt{\log2}$, so $|a|e^{a^2/4}\le C(e^{a^2}-1)$; the lower
bound on $|a|$ justifies this estimate for small arguments as well.

Fix $r>F_h(W_h)$ and an admissible current with
$\|J_1\|_{\psi_2},\|J_2\|_\infty\le r$. The weak no-flux identity,
first for smooth tests and then by Lipschitz approximation, gives
\[
 |\mathcal W(P_h^Du)|=\left|\int_DJ\cdot\nabla u\right|
 \le r\left(1+C\int_D\Phi(\partial_1u)
                   +\int_D|\partial_2u|\right).
\]
Here $\int_D(e^{(J_1/r)^2}-1)\le1$ because $D$ has unit area.
Taking the infimum in $r$ proves the claim. The inequality is
deterministic for each current and each test, so it does not use any
independence between $u$ and the noise. Since $P_h^Du$ lies in a fixed
finite cell space, its pairing with the realized $W_h$ is unambiguous.
\end{proof}

\section{Brownian local time and arbitrarily large reward}
\label{neg:sec:brownian}

This section constructs the model martingale: a bounded, predictable,
upper-triangular control whose terminal $\Psi$-cost is bounded by an
absolute constant while its determinant reward \eqref{neg:eq:detreward}
is as large as we wish. It is the only place where the endpoint
structure of $\Psi$ enters.

\begin{lemma}
\label{neg:lem:Brownian}
For every $D_0>0$ there are deterministic times
$0=t_0<t_1<\cdots<t_K=1$ and bounded continuous functions
\[
 L_j:(\mathbb R^2)^{j-1}\longrightarrow
       \left\{\begin{pmatrix}a&b\\0&d\end{pmatrix}:a,d>0\right\},
 \qquad j=1,\dots,K,
\]
such that $L_t=L_j(B_{t_1}-B_{t_0},\dots,B_{t_{j-1}}-B_{t_{j-2}})$
on $[t_{j-1},t_j)$ is predictable, takes values in a fixed compact
set of upper-triangular matrices with positive diagonals, and
\begin{equation}\label{neg:eq:Brownian-claim}
 \mathbb E\,\Psi\!\left(\int_0^1L_t\,\mathrm dB_t\right)
 \le C_*,\qquad
 \mathbb E\int_0^1\sqrt{\det L_t}\,\mathrm dt\ge D_0,
\end{equation}
where $C_*<\infty$ is absolute, independent of $D_0$.
\end{lemma}

\paragraph{Heuristics.}
The following formal computation explains why bounded cost and
unbounded reward are compatible; it is not used in the proof. In the
construction below, the control satisfies $\sigma_t\beta_t=R$ during
an active stage. The reward rate $\sqrt{\det\mathcal L_t}
=Y_t\sqrt R$ therefore has conditional mean $\sqrt R$, and the reward is
$\sqrt R$ times the active real duration. Stage $j$ lasts $R_j/R$
in real time, so the reward is of order $\sqrt R$ as soon as
$\sum_jR_j\ge R$. Each stage contributes a resource of order
$q\sqrt{K_q}$, with $q=\gamma/J$ and $K_q\approx2\log(J/\gamma)$.
Summing over $J$ stages gives
$\sum_jR_j\approx\gamma\sqrt{2\log(J/\gamma)}$, hence a reward of order
$(\log J)^{1/4}$, which is unbounded in $J$.

The cost stays bounded because the two coordinates are tied together.
By stage $j$, the multiplicative coordinate $Y$ has accumulated
logarithmic variance $s\le2S_j$, which inflates the logarithmic weight
in $\Phi$ by a factor $\sqrt{1+s}\lesssim\sqrt{S_j}$. Meanwhile the
local-time stopping confines the ratio $M=X/Y$ to $|M|\le b_j=S_j^{-1/2}$.
The product $|M|\sqrt{1+s}$ is therefore bounded at every stage; this
is the content of \eqref{neg:eq:terminal-cost}. In this way the
$\sqrt{\log}$ growth of the complementary Young function is exactly
compensated by the confinement $|M|\le S_j^{-1/2}$.

\begin{proof}
\emph{Step 1: independent operational stages.}
Each successful stage gains about $q\sqrt{\log(1/q)}$, at the price of
a failure probability $O(q)$. With $J$ stages and $q=\gamma/J$, the
total resource $\sum_jR_j$ defined below grows like $\sqrt{\log J}$, while the probability of some
failure stays small.
Fix a small absolute $\gamma>0$ and an integer $J\gg1$. Put
\[
 \begin{gathered}
 q=\frac\gamma J,\qquad K_q=\log(1+q^{-2}),\qquad
 S_j=2^j,\qquad b_j=S_j^{-1/2},\qquad \ell_j=q b_j,\\
 w_j(v)=\frac{S_j}{K_q(v+\ell_j^2)},\qquad 0\le v\le b_j^2.
 \end{gathered}
\]
In stage $j$ let $U$ be a fresh standard real Brownian motion started
at $0$, and stop it at
\[
 T_j=\tau_{\ell_j}\wedge\tau_{\{\pm b_j\}}\wedge b_j^2,
\]
where $\tau_{\ell}$ is the first time that Tanaka-normalized local time
at $0$ reaches $\ell$. A stage succeeds when $\tau_{\ell_j}$ is the
first stopping event; then $U_{T_j}=0$. Put
\[
 R_j=\int_0^{T_j}\sqrt{w_j(v)}\,\mathrm dv.
\]
Tanaka's formula at the bounded stop yields
$\mathbb E|U_{T_j}|=\mathbb EL^0_{T_j}\le\ell_j$, so
$\mathbb P(\text{exit at }\pm b_j)\le q$. L\'evy's local-time identity
and Brownian reflection imply
\begin{equation}\label{neg:eq:local-tail}
 c\min\{1,\ell/\sqrt v\}
 \le\mathbb P(\tau_\ell>v)
 \le C\min\{1,\ell/\sqrt v\}
 \qquad(\ell,v>0).
\end{equation}
The probability of reaching the time cap without success is therefore
at most $Cq$. For a sufficiently small absolute $c_0$ and
$\ell_j^2\le v\le c_0b_j^2$, subtracting the $O(q)$ total failure
probability from the left side of \eqref{neg:eq:local-tail} gives
\begin{equation}\label{neg:eq:successful-tail}
 \mathbb P(T_j>v,\ \text{stage $j$ succeeds})
      \ge c\frac{\ell_j}{\sqrt v}.
\end{equation}
Consequently, by the tail integral formula,
\begin{equation}\label{neg:eq:resource-bounds}
 \begin{aligned}
  \mathbb E[R_j;\,\text{success}]&\ge
  \frac{c\sqrt{S_j}\,\ell_j}{\sqrt{K_q}}
       \int_{\ell_j^2}^{c_0b_j^2}\frac{\mathrm dv}{v}
       \ge c q\sqrt{K_q},\\
  0\le R_j&\le\frac{\sqrt{S_j}}{\sqrt{K_q}}
      \int_0^{b_j^2}\frac{\mathrm dv}{\sqrt{v+\ell_j^2}}
       \le\frac2{\sqrt{K_q}}.
 \end{aligned}
\end{equation}
Each operational stage succeeds with probability at least $1-Cq$.
Thus the probability that all $J$ stages succeed is at least
$1-CJq\ge1-C\gamma$. Conditioned on this event, the resources
remain independent. Their sum $S_R$ satisfies
\[
 \mu_J:=\mathbb E[S_R\mid\text{all succeed}]
       \ge cJq\sqrt{K_q}=c\gamma\sqrt{K_q}\longrightarrow\infty,
 \qquad
 \operatorname{Var}(S_R\mid\text{all succeed})
       \le\frac{2}{\sqrt{K_q}}\mu_J.
\]
Fix any prescribed real-time resource $R>0$. Once $\gamma$ is chosen
sufficiently small and $J$ is large enough that $\mu_J\ge2R$,
Chebyshev's inequality gives
$\mathbb P(S_R<R\mid\text{all succeed})=o_J(1)$. Hence
\begin{equation}\label{neg:eq:reach-resource}
 \mathbb P\left(\text{all stages succeed and }
                   \sum_{j=1}^J R_j\ge R\right)\ge\frac12.
\end{equation}

\emph{Step 2: embedding in ordinary time.}
Run the stages successively using the first coordinate $B^1$ of a
planar Brownian motion; each completed successful stage restarts at
$M=0$. Within stage $j$, transform operational time $v$ into real time
by
\begin{equation}\label{neg:eq:clock}
 \begin{aligned}
  \theta_j(v)&=\frac1R\int_0^v\sqrt{w_j(w)}\,\mathrm dw,
  \qquad v=\theta_j^{-1}(t-t_j^{\mathrm{start}}),\\
  \sigma_t&=(\theta_j'(v))^{-1/2},
  \qquad\beta_t=\sqrt{w_j(v)}\,(\theta_j'(v))^{-1/2}.
 \end{aligned}
\end{equation}
Within a stage, $\sigma_t^2\,\mathrm dt=\mathrm dv$.
Consequently the increment of $M_t=\int_0^t\sigma_r\,\mathrm dB_r^1$,
reparametrized by $v$, is a standard Brownian motion. The strong
Markov property at the stage starting times makes successive
successful stages independent. The stopping rules depend only on
$B^1$, and $\sigma,\beta$ are predictable in that filtration.
Stop at the first failed stage, the end of stage $J$, or time $1$, and
put $\sigma=\beta=0$ afterward. During an active stage,
\begin{equation}\label{neg:eq:clock-identities}
  \sigma_t\beta_t=R,
  \qquad \int_{\mathrm{stage}\ j}\beta_t^2\,\mathrm dt
       \le\int_0^{b_j^2}w_j(v)\,\mathrm dv
       =\frac{S_j}{K_q}\log(1+q^{-2})=S_j.
\end{equation}
The real duration of stage $j$ is $R_j/R$. On the event in
\eqref{neg:eq:reach-resource}, every stage reached before time $1$
succeeds and the total available duration is at least $1$.
The process is therefore active for almost every real time
$t\in[0,1]$ on that event, so
\begin{equation}\label{neg:eq:sigbeta-gain}
 \mathbb E\int_0^1\sqrt{\sigma_t\beta_t}\,\mathrm dt
     \ge\frac12\sqrt R.
\end{equation}
For fixed $J,R$ the coefficients are bounded, though their bounds may
depend on $J,R$.

\emph{Step 3: a triangular two-dimensional martingale.}
Set
\[
 s_t=\int_0^t\beta_r^2\,\mathrm dr,
 \qquad Y_t=\exp\!\left(\int_0^t\beta_r\,\mathrm dB_r^2-\frac12s_t\right),
 \qquad X_t=Y_tM_t.
\]
The clock depends only on $B^1$, so conditionally on its entire path
$\mathbb E(Y_t\mid B^1)=1$. The product rule, with zero quadratic
covariation between the $B^1$ and $B^2$ integrals, gives
\begin{equation}\label{neg:eq:L-martingale}
 \mathrm d\begin{pmatrix}X_t\\Y_t\end{pmatrix}
 =\underbrace{\begin{pmatrix}
  Y_t\sigma_t&M_tY_t\beta_t\\0&Y_t\beta_t
 \end{pmatrix}}_{\mathcal L_t}
 \mathrm d\begin{pmatrix}B_t^1\\B_t^2\end{pmatrix},
 \qquad \sqrt{\det\mathcal L_t}
       =Y_t\sqrt{\sigma_t\beta_t}.
\end{equation}
It follows from \eqref{neg:eq:sigbeta-gain} that
\begin{equation}\label{neg:eq:reward-lower}
 \mathbb E\int_0^1\sqrt{\det\mathcal L_t}\,\mathrm dt
 \ge\frac12\sqrt R.
\end{equation}

We now bound the terminal cost. If $j$ is the last stage visited, then
$|M_1|\le b_j$ and $s_1\le\sum_{k\le j}S_k\le2S_j$. Conditionally on
$B^1$ and under weighting by $Y_1$, the random variable $\log Y_1$ is
Gaussian with mean $s_1/2$ and variance $s_1$. Since
$\Phi(x)\le C|x|\sqrt{\log(2+|x|)}$ and $|M_1|\le1$,
\begin{equation}\label{neg:eq:terminal-cost}
 \begin{aligned}
 \mathbb E[\Phi(X_1)\mid B^1]
 &\le C|M_1|\,\mathbb E\!
    \left[Y_1\sqrt{\log(2+Y_1|M_1|)}\mid B^1\right]\\
 &\le C|M_1|\sqrt{1+s_1}\le C.
 \end{aligned}
\end{equation}
Also $\mathbb E|Y_1-1|\le2$. The convex cost concerns $(X_1,Y_1-1)$,
not $(X_1,Y_1)$: its second coordinate is centered. Since
$(X_0,Y_0-1)=(0,0)$, we have proved
\[
 \mathbb E\Psi(X_1,Y_1-1)\le C_*
 \quad\text{with $C_*$ independent of $J,R$.}
\]

\emph{Step 4: bounded continuous simple controls.}
First stop the pair $(X,Y-1)$ at $Y=H$, or at time $1$. Conditional
Jensen and optional sampling preserve the terminal $\Psi$ bound.
Before this stopping time, $Y_t\le H$, $|M_t|\le1$,
and the coefficients $\sigma,\beta$ are bounded for fixed
$J,R$; hence the stopped integrand is bounded. As $H\to\infty$,
monotone convergence recovers the reward \eqref{neg:eq:reward-lower}.
Choose $R$ so large that this reward exceeds $D_0+2$, and then choose
$H$ so that the stopped reward exceeds $D_0+1$.

It remains to pass to simple controls with continuous coefficients.
Approximate the bounded predictable, upper-triangular
nonnegative-diagonal integrand in $L^2(\mathbb P\otimes\mathrm dt)$ by
simple predictable processes on deterministic time intervals.
Approximate all finite-dimensional coefficient functions by bounded
continuous functions of finitely many earlier Brownian increments.
Refine the partition to include observation times, clip each matrix to
a common compact box of upper-triangular matrices with nonnegative
diagonals, and add $\varepsilon I$. These operations preserve $L^2$
convergence as $\varepsilon\downarrow0$. The It\^o isometry gives
convergence of terminal integrals in $L^2$, hence convergence of their
$\Psi$ costs in $L^1$, since $\Psi(z)\le C(1+|z|^p)$ for some $p<2$ and
the approximants are uniformly square integrable. Continuity of
$\sqrt{\det L}$ on the common compact box gives convergence of the
rewards in $L^1(\mathbb P\otimes\mathrm dt)$. Fix one sufficiently
accurate approximation. The diagonal lower bound $\varepsilon>0$,
the deterministic partition and the terminal estimate are then as
claimed, with an absolute $C_*$ after increasing it once.
\end{proof}

\section{Realization by spatially adapted Gaussian bubbles}
\label{neg:sec:spatial}

This section transfers the Brownian model of \cref{neg:lem:Brownian}
to test potentials on $D$ (\cref{neg:lem:spatial}). The model is
realized without sampling Brownian motion in physical space: the
Brownian increments are replaced by bounded innovations, obtained by
following the children of a uniform spatial point through a tree of
triangles. The family of bubbles is chosen with the help of white
noise, but the resulting spatial estimates hold for every fixed
realization of the chosen signs.

\paragraph{Reference geometry.}
Let $T=\{(x_1,x_2):x_1,x_2\ge0,\ x_1+x_2\le1\}$. Choose four
disjoint small axis-parallel dyadic squares strictly inside $T$,
whose centers form a square, and put a translated tent
$(r-\|x-c\|_\infty)_+$ in each, with checkerboard signs. Their sum
$b$ is compactly supported in $\operatorname{int}T$, piecewise affine,
and satisfies
\begin{equation}\label{neg:eq:affine-moments}
  \int_T b(x)\,\ell(x)\,\mathrm dx=0
  \quad\text{for every affine }\ell.
\end{equation}
The checkerboard arrangement cancels both the total mass and the first
moments; this is what will give Gaussian orthogonality between
ancestors and descendants. The law of $\nabla b(x)$ for uniform
$x\in T$ is supported on
$0,\pm e_1,\pm e_2$ and symmetric under signed coordinate
permutations. Choose $c_b>0$ so that
$|T|^{-1}\int_T\nabla b\otimes\nabla b=c_b^2 I$, and denote by
$\nu$ the law of $c_b^{-1}\nabla b(U)$ for uniform $U\in T$.
Then $\nu$ is bounded, centered, has covariance $I$, and is invariant
under signed coordinate permutations. Fix a sufficiently fine finite
dyadic triangulation into congruent right-isosceles triangles
of leg $2^{-k_0}$ on which $b$ is affine. The integer $k_0$ and
template $b$ are absolute.

\paragraph{Two conditional laws.}
The next two lemmas serve different purposes. In the first,
we condition on earlier Gaussian observations to justify
freshness of the next selected bubble. In the second, we freeze all
bubble signs and condition on the spatial branch to obtain the
i.i.d.\ innovation law and Jensen's stopping inequality.

\begin{lemma}
\label{neg:lem:fresh}
Suppose that a finite sequence of nonzero Hilbert-space vectors
is chosen predictably from an isonormal Gaussian process: at step $k$, the vector $b_k$ is
measurable with respect to previously revealed pairings
$\mathcal W(b_1),\dots,\mathcal W(b_{k-1})$, takes one of finitely
many values on each finite history, and is orthogonal to every
previous vector on that history. Let $\mathcal G_{k-1}$ denote that
history. Then, conditionally on $\mathcal G_{k-1}$,
$\mathcal W(b_k)$ is centered Gaussian with variance $\|b_k\|_2^2$.
For every finite-valued $\mathcal G_{k-1}$-measurable Hilbert-space
vector $v_k$,
\begin{equation}\label{neg:eq:conditional-sign}
 \mathbb E\left[\operatorname{sgn}\mathcal W(b_k)
                    \;\mathcal W(v_k)\mid\mathcal G_{k-1}\right]
 =\sqrt{\frac2\pi}\,
     \frac{\langle v_k,b_k\rangle}{\|b_k\|_2}.
\end{equation}
\end{lemma}

Here $v_k$ may itself depend on the previous observations and need
not be orthogonal to previous directions. The finite-valued assumption
is the one used below, and it makes all random isonormal pairings
unambiguous.

\begin{proof}
We proceed inductively through the finite decision tree. On any
branch determined by previously observed pairings, the selected
directions are fixed vectors, and the event selecting that branch is
measurable with respect to their Gaussian coordinates. Conditional on
the actual observed values, the restriction of white noise to the
orthogonal complement of their span remains a fresh isonormal
Gaussian process. The next direction belongs to this
complement, so its normalized pairing is a fresh standard Gaussian.
Decompose $v_k=\langle v_k,b_k\rangle b_k/\|b_k\|_2^2+r_k$.
Although $\mathcal W(r_k)$ may have a past-measurable component, its
remaining centered component is conditionally independent of
$\mathcal W(b_k)$, and the past-measurable component has zero pairing
with its centered sign. Since $\mathbb E|G|=\sqrt{2/\pi}$ for
$G\sim N(0,1)$, the formula follows. An induction over the finite
histories justifies the conditioning.
\end{proof}

\begin{lemma}
\label{neg:lem:branch}
Fix a bounded continuous simple control $L$ as in
\cref{neg:lem:Brownian} and a deterministic assignment
$Q\mapsto s(Q)\in\{\pm1\}$ of signs to all potential parents.
Let $U$ be uniform on $D$, independent of an auxiliary sequence of
independent $\nu$-distributed vectors, and let
$(\mathcal S_k)_{k=0}^n$ be a filtration with the following properties.
\begin{enumerate}[label=(\alph*),leftmargin=2.5em]
\item \emph{Complete congruent subdivisions.} Before microstep $k$,
  the active triangle containing $U$ is $\mathcal S_{k-1}$-measurable
  and has the form $Q=x_Q+d_QAO_QT$, where $A$ is the current chart,
  $O_Q$ is a signed coordinate permutation, and $x_Q,d_Q,A,O_Q$ are
  $\mathcal S_{k-1}$-measurable. Its children are the images of the
  reference triangulation under the same affine map.
\item \emph{Signs fixed before revelation.} The increment recorded at
  microstep $k$ is
  \[
   \xi_k=s(Q)\,O_Q\,c_b^{-1}\nabla b(\widehat U),
   \qquad \widehat U=d_Q^{-1}O_Q^TA^{-1}(U-x_Q)\in T,
  \]
  and
  $\mathcal S_k$ is generated by $\mathcal S_{k-1}$ and the child of
  $Q$ containing $U$.
\item \emph{Predictable retiling.} At the end of a control block, each
  active triangle is retiled into complete triangles of the form in
  (a) for the new chart, plus a discarded remainder. Which retained
  triangle contains $U$, or whether $U$ lies in the discarded
  remainder, is $\mathcal S$-measurable before the next microstep.
  After $U$ is discarded, the increments $\xi_k$ are the subsequent
  auxiliary $\nu$-vectors, which are added to $\mathcal S_k$.
\item \emph{No further information.} Apart from (a)--(c),
  $\mathcal S_k$ carries no information on the position of $U$
  inside its current triangle.
\end{enumerate}
The bubble construction of \cref{neg:lem:spatial}, with all signs
frozen after the white noise is fixed, satisfies (a)--(d).
Then the ideal increments $\xi_1,\dots,\xi_n$
are independent with law $\nu$. The actual gradient path is the
ideal vector martingale stopped at the first discarded region. In
particular, for the convex $\Psi$ in \eqref{neg:eq:PhiPsi},
\begin{equation}\label{neg:eq:stopped-Jensen}
 \mathbb E_{U\sim\mathrm{Unif}(D)}
      \Psi(G_{\tau}(U))
 \le\mathbb E_{U,\mathrm{aux}}\Psi(G_n),
\end{equation}
where $G_k$ is the corresponding ideal gradient martingale,
$\tau$ is the geometric stopping index, and the expectation on
the right also includes the independent innovations appended after
loss.
\end{lemma}

Once all signs are frozen, \cref{neg:lem:branch} is a statement about
the uniform point $U$ and the auxiliary innovations only; the white
noise plays no role in it.

\begin{proof}
Fix all bubble signs. With $U$ uniform on $D$, the initial
current triangle is one of two congruent halves of $D$. Conditional
on its identity, $U$ is uniform in that triangle. At a microstep,
the complete affine child triangulation partitions the parent into
congruent equal-area triangles. Thus the normalized reference
increment has law $\nu$, irrespective of the parent, its signed
permutation $O$, and the already fixed sign. Conditional on any
retained new-chart triangle after a retiling, $U$ is again uniform
in that triangle. The portion discarded at a retiling is identified
\emph{before} the next increment; thereafter fresh independent
artificial $\nu$-increments supply the ideal continuation.

By (b)--(d), the filtration records the active parent, the
selected child and its innovation, each retiling/survival decision,
and the appended innovations after loss; it does not disclose the
exact location of $U$ inside the current triangle. Hence,
conditional on $\mathcal S_{k-1}$, the point $\widehat U$ is uniform
in $T$ on survival, and by the invariance of $\nu$ under signed
coordinate permutations the next ideal increment always has law $\nu$.
Iterated conditioning proves the claimed i.i.d.\ law and the martingale
property of $G_k$. Discarding is a stopping time for this filtration,
and conditional Jensen for $G_n$ at the stopped $\sigma$-algebra yields
\eqref{neg:eq:stopped-Jensen}. The geometric decisions may depend on
previous innovations through the chosen charts; predictability is
sufficient.
\end{proof}

\begin{lemma}
\label{neg:lem:spatial}
Fix a simple control $L$ supplied by \cref{neg:lem:Brownian} and put
$a_t=\sqrt{\det L_t}$. There are mean-zero random piecewise affine
Lipschitz potentials $u_m$ on $D$, built from $\mathcal W$ and
$h=2^{-m}$, such that
\begin{align}
 \int_D\Psi(\partial_1u_m,\partial_2u_m)\,\mathrm dx
   &\le\mathbb E\Psi\!\left(\int_0^1L_t\,\mathrm dB_t\right)
          +o_L(1)\quad\text{almost surely},
       \label{neg:eq:spatial-cost}\\
 \liminf_{m\to\infty}\frac{\mathbb E\,\mathcal W(P_h^Du_m)}{\sqrt m}
   &\ge c_*\,\mathbb E\int_0^1a_t\,\mathrm dt,
       \label{neg:eq:spatial-reward}
\end{align}
where $c_*>0$ is absolute. The error and the minimum admissible
$m$ may depend on the fixed control, but $c_*$ does not.
\end{lemma}
\begin{proof}
The potential $u_m$ is a sum of signed bubbles on an adaptive tree of
triangles. Steps~1 and~2 build the tree and prove the cost bound
\eqref{neg:eq:spatial-cost} for every choice of signs; Steps~3 and~4
average over the Gaussian signs and prove the reward bound
\eqref{neg:eq:spatial-reward}.

\emph{Step 1: charts and depths.}
Let $0=t_0<\cdots<t_K=1$ be the fixed control partition and put
\[
 n=\left\lfloor\frac{m}{4k_0}\right\rfloor,
 \qquad n_j=\lfloor nt_j\rfloor-\lfloor nt_{j-1}\rfloor.
\]
The $j$th block has $n_j$ microsteps. Along any active branch in that
block, the matrix $L_j$ depends only on completed previous block
increments. Write $a_j=\sqrt{\det L_j}$ and choose the chart $A_j$
by
\begin{equation}\label{neg:eq:chart}
 A_j^{-T}=\frac{L_j}{a_j},\qquad \det A_j=1.
\end{equation}
The chart and its inverse have operator norms bounded by a
constant depending only on the fixed control. An active
triangle can be written
$Q=x_Q+d_QA_jO_QT$ with $O_Q$ a signed coordinate permutation.
On each microstep it is subdivided using the image of the fixed
reference triangulation.

At the beginning of each new control block, freeze the old
gradient (which is affine on each current triangle) and retile that
triangle on a new-chart square lattice of spacing
$2^{-\lfloor\sqrt m\rfloor}d_Q$. Retain only lattice squares wholly
inside the triangle, and split each square into two. Since the old and
new charts have uniformly bounded operator norms for the fixed control,
the omitted boundary strip has relative area at most
$C_L2^{-\lfloor\sqrt m\rfloor}$, uniformly in the parent.
The initial two triangles are treated similarly. Because there are
only $K+1$ retilings, the total discarded area is at most
\begin{equation}\label{neg:eq:area-loss}
 \varepsilon_m\le C_L(K+1)2^{-\lfloor\sqrt m\rfloor}=o_L(1).
\end{equation}
For every active triangle,
\begin{equation}\label{neg:eq:leg-lower}
 d_Q\ge2^{-k_0n-(K+1)\lfloor\sqrt m\rfloor}
       \ge2^{-m/3}
\end{equation}
for all sufficiently large $m$ depending on $L$.

\emph{Step 2: bubbles and spatial Brownian limit.}
For each active parent $Q$ in block $j$, define, extended by zero
off $Q$,
\begin{equation}\label{neg:eq:bubble}
 b_Q(x)=\frac{d_Qa_j}{c_b\sqrt n}\,
 b\!\left(d_Q^{-1}O_Q^TA_j^{-1}(x-x_Q)\right).
\end{equation}
Reveal $\mathcal W(b_Q)$, select
$s_Q=\operatorname{sgn}\mathcal W(b_Q)$, and add $s_Qb_Q$ to $u_m$.
The total potential is continuous, piecewise affine, and Lipschitz,
because each bubble is supported strictly inside its parent and the
tree is finite. Moreover $\int_Db_Q=0$ by \eqref{neg:eq:affine-moments},
so $\int_Du_m=0$.

On a child of $Q$, the new gradient increment is exactly
\begin{equation}\label{neg:eq:grad-increment}
 \nabla(s_Qb_Q)=\frac{L_j\xi_Q}{\sqrt n},
 \qquad \xi_Q=s_Qc_b^{-1}O_Q\nabla b.
\end{equation}
At each block end, store the sum of the normalized innovations
$\xi_Q/\sqrt n$ along the branch. The next block's $L_{j+1}$ is
evaluated from the completed earlier block sums, before any new
bubble is revealed. The spatial branch lemma \cref{neg:lem:branch}
applies even though the signs were chosen from $\mathcal W$, because we
first freeze all those signs. Thus the spatial random point is the only
randomness used in the martingale law; the Gaussian noise is averaged
only in Step~3.

The ideal block sums converge jointly to independent Brownian
increments with covariances $(t_j-t_{j-1})I$ by the multivariate
central limit theorem. Since all coefficient maps $L_j$ are bounded
and continuous, the ideal terminal gradient converges in law to
$\int_0^1L_t\,\mathrm dB_t$. Its increments have magnitude at most
$C_L/\sqrt n$, so conditional Hoeffding--Azuma bounds give
uniformly sub-Gaussian tails for the ideal gradient. The function
$\Psi$ has subquadratic growth. Therefore its terminal expectations
converge, and the stopped-martingale Jensen bound
\eqref{neg:eq:stopped-Jensen} gives
\eqref{neg:eq:spatial-cost} uniformly over every possible
assignment of bubble signs.

\emph{Step 3: freshness of Gaussian rewards.}
Reveal the bubbles breadth first, with a deterministic order within
each level. When $b_Q$ is chosen, all previous Gaussian pairings
and their signs are already known. Every previous bubble is either
disjoint from $Q$ or affine on $Q$. Hence \eqref{neg:eq:affine-moments}
gives
\begin{equation}\label{neg:eq:orthogonality}
 \langle b_Q,b_{Q'}\rangle_{L^2(D)}=0
 \quad\text{whenever $b_{Q'}$ was revealed earlier}.
\end{equation}
For a fixed $m$, every past innovation belongs to a finite set, and
the simple control sees only finitely many completed block sums.
Hence only finitely many chart and bubble choices can arise, and
\cref{neg:lem:fresh} applies to $v_Q=P_h^Db_Q$. Since $P_h^D$ is an
orthogonal projection,
\begin{equation}\label{neg:eq:fresh-gain}
 \mathbb E\left[s_Q\mathcal W(P_h^Db_Q)\mid\mathcal G_{Q^-}\right]
 =\sqrt{\frac2\pi}\,
    \frac{\|P_h^Db_Q\|_2^2}{\|b_Q\|_2}.
\end{equation}
The square-cell Poincar\'e inequality and affine scaling give
\begin{equation}\label{neg:eq:projection-error}
 \|P_h^Db_Q-b_Q\|_2\le Ch\|\nabla b_Q\|_2
 \le C_L\frac{h}{d_Q}\|b_Q\|_2
 =o_L(1)\|b_Q\|_2
\end{equation}
uniformly in active $Q$, by \eqref{neg:eq:leg-lower}.
Writing $|T|=1/2$, the determinant-one change of variables in
\eqref{neg:eq:bubble} also shows
\begin{equation}\label{neg:eq:bubble-norm}
 \|b_Q\|_2=\frac{\|b\|_{L^2(T)}}{c_b|T|}\,
              \frac{a_j|Q|}{\sqrt n}.
\end{equation}
Thus for all large $m$, with a universal positive $c$,
\begin{equation}\label{neg:eq:gain-per-parent}
 \mathbb E[s_Q\mathcal W(P_h^Db_Q)\mid\mathcal G_{Q^-}]
       \ge \frac{c\,a_j|Q|}{\sqrt n}.
\end{equation}

\emph{Step 4: sum across the tree.}
At any microstep, the sum of $a_j|Q|$ over active parents equals the
spatial expectation of the coefficient $a_j$ along a surviving
branch. By \eqref{neg:eq:area-loss}, replacing surviving by ideal
branches changes this expectation by at most $C_L\varepsilon_m$.
For each fixed sign realization the ideal law is the i.i.d.\ law
of \cref{neg:lem:branch}, so the ideal expectation does not depend
on the signs. Taking Gaussian expectations, summing
\eqref{neg:eq:gain-per-parent} over the $n$ microsteps, and using
the blockwise central limit theorem gives
\[
 \liminf_{m\to\infty}
   \frac{\mathbb E\mathcal W(P_h^Du_m)}{\sqrt m}
 \ge c\lim_{m\to\infty}\sqrt{\frac nm}
 \frac1n\sum_{\text{microsteps}}
              \mathbb E_{\rm ideal} a_{j(\text{step})}
 =\frac{c}{\sqrt{4k_0}}\,
    \mathbb E\int_0^1\sqrt{\det L_t}\,\mathrm dt.
\]
The final constant depends only on the fixed reference template,
not on the control or the prescribed reward. This is
\eqref{neg:eq:spatial-reward}.
\end{proof}

\section{Divergence of the Gaussian gauge and failure of matching}
\label{neg:sec:conclusion}

We combine the three previous sections to prove the Gaussian
divergence \eqref{neg:eq:gaussian-negative}, first in expectation and
then in probability, and transfer it to matching with the comparisons
of Part~\ref{umc:partII}.

\begin{theorem}
\label{neg:thm:Gaussian}
For the projected \emph{square} white noise in
\eqref{neg:eq:coefficients},
\begin{equation}\label{neg:eq:Gaussian-divergence}
 \frac{F_{2^{-m}}(W_{2^{-m}})}{\sqrt m}
  \xrightarrow[m\to\infty]{\mathbb P}+\infty.
\end{equation}
\end{theorem}
\begin{proof}
Fix an arbitrary $D_0>0$ and choose once and for all the
simple control of \cref{neg:lem:Brownian}. By
\cref{neg:lem:spatial,neg:lem:dual}, for all sufficiently large $m$,
\[
 |\mathcal W(P_h^Du_m)|\le
 C_D\mathcal F_m(C_*+2)
 \quad\text{almost surely},
\]
while
\[
 \liminf_{m\to\infty}
      \frac{\mathbb E\mathcal W(P_h^Du_m)}{\sqrt m}
 \ge c_* D_0.
\]
Consequently
\begin{equation}\label{neg:eq:expectation-divergence}
 \liminf_{m\to\infty}\frac{\mathbb E\mathcal F_m}{\sqrt m}
 \ge\frac{c_*D_0}{C_D(C_*+2)}.
\end{equation}
The denominator and $c_*$ are independent of $D_0$. We first send
$m\to\infty$ for this fixed control, and then let $D_0$ range over
arbitrarily large fixed values. This proves
$\mathbb E\mathcal F_m/\sqrt m\to\infty$.

The mesh-independent Lipschitz estimate
\cref{red:lem:uniform-gaussian-lipschitz} and the Gaussian
concentration inequality imply, for every $t>0$,
\[
 \mathbb P\{\mathcal F_m\le\mathbb E\mathcal F_m-t\}
       \le\exp\left(-\frac{t^2}{2C_{\rm Lip}^2}\right).
\]
For any fixed $A$, eventually $\mathbb E\mathcal F_m\ge2A\sqrt m$.
Substituting $t=A\sqrt m$ gives
\eqref{neg:eq:Gaussian-divergence}.
\end{proof}

The equivalence of \cref{red:thm:main} alone would only show that no
fixed constant bounds $M_N/s_N$ with probability tending to one. To
obtain divergence in probability, we transfer the estimates at the
level of probabilities.

\begin{proof}[Proof of \cref{neg:thm:main}]
Choose the dyadic mesh $h_N$ of \eqref{red:eq:h-choice}, with
$A_0s_N\le h_N<2A_0s_N$ and
$m(N)=-\log_2h_N\asymp\log N$. On the high-probability histogram
ellipticity event of \cref{red:lem:hist-ellipticity}, the
atomic-to-histogram comparison \cref{red:prop:atomic-hist} and the
endpoint transport--flow theorem \cref{red:thm:transport-flow-endpoint}
give, with constants independent of $N$,
\begin{equation}\label{neg:eq:eventwise-matching}
 F_{h_N}\big(P_{h_N}^D(\mu_N-\nu_N)\big)
 \le C\big(M_N+h_N\big).
\end{equation}
In the coupling of \cref{red:thm:KMT} there is an exact divergence
corrector $Q_{h_N}$ with
\begin{equation}\label{neg:eq:eventwise-KMT}
 P_{h_N}^D(\mu_N-\nu_N)
   =\sqrt{\frac2N}\,W_{h_N}+\operatorname{div}Q_{h_N},
 \qquad\max\{\|Q_{h_N,1}\|_{\psi_2},\|Q_{h_N,2}\|_\infty\}
       \le C s_N
\end{equation}
on an event of probability $1-o(1)$ (the finite horizontal
coordinate has an even smaller $o_{\mathbb P}(s_N)$ correction).
The triangle inequality for the minimum-flow gauge
(\cref{red:lem:flow-lipschitz}) then yields, pathwise on this event,
\begin{equation}\label{neg:eq:transfer}
 \sqrt{\frac2N}\,F_{h_N}(W_{h_N})
 \le C\big(M_N+h_N+s_N\big).
\end{equation}
For each fixed $A<\infty$, the event $\{M_N\le As_N\}$ therefore
implies $\{\mathcal F_{m(N)}\le C_A\sqrt{m(N)}\}$ on the good
coupling event. By \cref{neg:thm:Gaussian},
\[
 \mathbb P(M_N\le As_N)
 \le\mathbb P\{\mathcal F_{m(N)}\le C_A\sqrt{m(N)}\}+o(1)
 \longrightarrow0.
\]
This proves \eqref{neg:eq:negative-result}.

For the one-cloud statement, fix any deterministic
reference $z^{(N)}\in D^N$ and let $p_N(A;z^{(N)})$ be the probability
of matching one uniform sample to it with both coordinate costs at
most $As_N$.
For two independent samples, the corresponding events are
independent and each has probability $p_N(A;z^{(N)})$. Composing the two
permutations via their common reference, and using the triangle
inequality for the discrete Luxemburg norm and for the supremum,
produces a two-sample matching of cost at most $2As_N$. Thus
\[
 p_N(A;z^{(N)})^2\le\mathbb P(M_N\le2As_N)\longrightarrow0.
\]
The right side is independent of the reference; taking the
supremum over $z^{(N)}\in D^N$ proves
\eqref{neg:eq:reference-negative}. Finally, the symmetry of the
sample law under exchange of axes proves the other hard endpoint.
\end{proof}
\clearpage
\phantomsection
\addcontentsline{toc}{part}{Appendix}
\appendix
\section{Geometric estimates for the adaptive construction}\label{app:geometry}

This appendix collects the deterministic geometric estimates behind
the adaptive construction of Part~\ref{umc:partI}. Throughout, $K>1$
is the distortion bound between consecutive metrics, $c_0=c_0(K)$ the
inball ratio, $s$ the terminal scale, and $\theta,D_0,\vartheta$ the
Whitney parameters of Section~\ref{sec:geom}. Appendix~\ref{app:birthchart}
compares a cell's sizes in its birth metric and in later metrics;
these contraction bounds feed the ancestor-decay factor
$\vartheta\sqrt K$ used in Sections~\ref{sec:bootstrap}
and~\ref{sec:phase}. Appendix~\ref{app:kernelchain} proves the bounds
\eqref{eq:singular} and \eqref{eq:support} on the ray kernel, which
are the input of Lemma~\ref{lem:pk}. Appendix~\ref{app:allbranches}
gives the finite Whitney partition with boundary receivers used in
Lemma~\ref{lem:whitney}, fixes the order of the constants, and records
the indexing of actual refinements along a revealed ancestral chain.

\subsection{Chart contraction and terminal cells}\label{app:birthchart}
We use the metric convention
\[
 d_G(x,z)=|G^{-1/2}(x-z)|,\qquad
 \operatorname{diam}_G E=\sup_{x,z\in E}d_G(x,z).
\]
\begin{lemma}\label{lem:birthchart}
Fix $K>1$. At an active generation let $Q$ be a square of side $\ell$
in its birth chart $G$, where $\det G=1$. Let $G^+$ be the
predictably selected next metric, with
\[
 K^{-1}G\preceq G^+\preceq KG,\qquad \det G^+=1.
\]
Let $D=\operatorname{diam}_{G^+}Q$. Suppose that $D>D_0s$
and that $Q$ is split by the balanced Whitney construction of
Lemma~\ref{lem:whitney}. Write $C_A$ for a uniform bound such that
every terminal union created by this split has $G^+$-diameter at
most $C_As$. Then the following claims hold simultaneously:
\begin{align}
 \sqrt{2/K}\,\ell&\le D\le \sqrt{2K}\,\ell,\label{eq:appD}\\
 \ell(S)&\le \theta\sqrt{2K}\,\ell
 \quad\text{for every nonterminal child }S,\label{eq:appside}\\
 \operatorname{diam}_{G^+}S
 &\le (C_A\sqrt{2K}/D_0)\,\ell
 \quad\text{for every terminal union }S.\label{eq:appterm}
\end{align}
In particular, for any prescribed $\vartheta>0$ one can choose
$\theta>0$ sufficiently small, and subsequently $D_0$ sufficiently
large, so that (i) every new nonterminal side obeys
$\ell(S)\le\vartheta\ell$ and (ii) every newly terminal
union obeys $\operatorname{diam}_{G^+}S\le\vartheta\ell$.
The choices depend only on the predetermined shape and clipping
parameters, not on $s$, the depth, the parent location, or the
eccentricity of $G$ in fixed physical coordinates.
A parent declared terminal without splitting has
$\operatorname{diam}_G Q\le\sqrt K D_0s$.
\end{lemma}
\begin{proof}
The transformation from birth coordinates to new coordinates is
\[
 A=(G^+)^{-1/2}G^{1/2}.
\]
Since
$K^{-1}G\preceq G^+\preceq KG$, its singular values lie in
$[K^{-1/2},K^{1/2}]$. The diameter of a Euclidean square of
side $\ell$ is $\sqrt2\ell$; this proves \eqref{eq:appD}.
The Whitney upper cap in the new chart is $\ell(S)\le\theta D$,
which combined with \eqref{eq:appD} proves \eqref{eq:appside}.
For a terminal child $\operatorname{diam}_{G^+}S\le C_As$.
Since the parent was split, $D>D_0s$, and by
$D\le\sqrt{2K}\ell$ we have
$s<\sqrt{2K}\ell/D_0$. This proves \eqref{eq:appterm}.
Take $\theta<\vartheta/\sqrt{2K}$ and
$D_0>C_A\sqrt{2K}/\vartheta$, together with the other finite lower
bounds of Lemma~\ref{lem:exactcover}. The constant $C_A$ of the
receiver construction is determined before $D_0$ is chosen and is
independent of $D_0$. Finally, a no-split parent with $D\le D_0s$ satisfies
\[\operatorname{diam}_G Q=\sqrt2\ell\le\sqrt K D\le\sqrt K D_0s.\]
\end{proof}

\begin{lemma}
\label{lem:ancestorindexed}
Suppose the sequence of active ancestors of a cell $Q_k$ is
$Q_j\supset Q_{j+1}\supset\cdots\supset Q_k$, with each
$Q_p$ a nonterminal square of side $\ell_p$ in $G_p$ for
$j\le p\le k$, and
\[
 \ell_{p+1}\le\vartheta\ell_p,\qquad
 K^{-1}G_p\preceq G_{p+1}\preceq KG_p.
\]
Then
\begin{equation}\label{eq:appancestor}
 \frac{\operatorname{diam}_{G_j}Q_k}{\ell_j}
 \le\sqrt2(\vartheta\sqrt K)^{k-j}.
\end{equation}
If instead $Q_k$ is a terminal union born from the last active
parent $Q_{k-1}$, the same inequality holds with $\sqrt2$
replaced by an absolute fixed constant.
\end{lemma}

A terminal leaf is never refined again, so the lemma concerns only
actual refinements.
\begin{proof}
By multiplication of the consecutive Loewner comparisons,
$G_j^{-1}\preceq K^{k-j}G_k^{-1}$, whence
$\operatorname{diam}_{G_j}Q_k\le
 K^{(k-j)/2}\operatorname{diam}_{G_k}Q_k$.
For a square $\operatorname{diam}_{G_k}Q_k=\sqrt2\ell_k$
and $\ell_k\le\vartheta^{k-j}\ell_j$.
For a terminal union use the terminal version of
Lemma~\ref{lem:birthchart} at the last step, giving the same
product with a fixed prefactor (Figure~\ref{fig:metric-refinement}). Once terminal the cell is frozen.
\end{proof}

\begin{figure}[H]
\centering
\begin{tikzpicture}[>=Latex,font=\footnotesize,line join=round]
  % square in birth chart
  \draw[line width=.8pt] (0.1,.55) rectangle (2.2,2.65);
  \node[align=center] at (1.15,3.05) {birth chart $G_j$};
  \node at (1.15,1.6) {$Q_j$};
  \draw[<->,thin] (.1,.28)--(2.2,.28);
  \node[anchor=north] at (1.15,.23) {$\ell_j$};
  \draw[->,line width=.8pt] (2.45,1.6)--(4.10,1.6);
  \node[align=center] at (3.26,2.12) {predictable\\metric update};
  \node[align=center] at (3.26,1.18) {$G_j\longmapsto G_j^+$};
  % same parent in updated chart, sheared box
  \draw[line width=.8pt] (4.42,.56)--(6.83,.56)--(7.35,2.64)--(4.94,2.64)--cycle;
  \draw[fill=black!16,line width=.75pt] (5.1,1.0) rectangle (5.79,1.69);
  \draw[fill=black!30,line width=.75pt] (6.14,.83)--(6.77,.83)--(6.67,1.36)--(6.14,1.36)--cycle;
  \node[align=center] at (5.85,3.05) {updated chart $G_j^+$};
  \node[align=center] at (5.85,.18) {$D=\diam_{G_j^+}Q_j$};
  % branches
  \draw[->,line width=.75pt] (7.60,1.95)--(9.02,2.50);
  \draw[->,line width=.75pt] (7.60,1.26)--(9.02,.69);
  \draw[fill=black!12,line width=.8pt] (9.25,2.19) rectangle (10.22,3.16);
  \node[anchor=west] at (10.35,2.7) {active $Q_{j+1}$};
  \node[anchor=west] at (10.35,2.28) {$\ell_{j+1}\le\vartheta\ell_j$};
  \draw[fill=black!30,line width=.8pt] (9.26,.38)--(9.76,.42)--(9.97,.77)--(9.39,.96)--cycle;
  \node[anchor=west] at (10.35,.73) {terminal $S$};
  \node[anchor=west] at (10.35,.32) {final reveal only};
\end{tikzpicture}
\caption{One actual refinement step. The parent is a square in its birth metric $G_j$ but is generally oblique in the updated chart $G_j^+$. A newly active child receives birth metric $G_{j+1}=G_j^+$; a terminal union receives its last increment and is not refined further. Comparing diameters in earlier birth metrics also incurs the factor $\sqrt K$ per actual chart transition.}
\label{fig:metric-refinement}
\end{figure}

\paragraph{Finite depth and area floor.}
All the square sides and diameter contractions in
Lemmas~\ref{lem:birthchart} and~\ref{lem:ancestorindexed} are measured
in the specified birth charts. A split-created terminal union enters
only through the estimate for its final transition; it is not a
square in a subsequent chart and is never split again. The factor
$\sqrt K$ in \eqref{eq:appancestor} comes in addition to the factor
$\sqrt{2K}$ in the newborn-side bound \eqref{eq:appside}, and each
is counted exactly once.
Every active square is born with side at least $2s$, by construction;
its physical area is the square of that side because $\det G=1$.
Since $\ell_p\le\vartheta^p\ell_0$ along a chain of active
squares, the number of strict refinements is at most
\[
 J_0\le 2+\frac{\log(\ell_0/(2s))}{|\log\vartheta|}
 =O(\log(1/s))=O(m).
\]
Every terminal union contains an entire uncovered $s$-square;
therefore all cells, at every revealed stage, have area at least $s^2$.

\subsection{Ray kernel estimates}\label{app:kernelchain}
We work on a convex parent in the intrinsic chart of
Section~\ref{sec:geom} and, as in the proof of Lemma~\ref{lem:pk},
drop the hats: $Q$, $x,y$, $K_Q$ and $\omega_Q$ stand for
$\widehat Q$, $z,w$, $\widehat K_Q$ and $\widehat\omega_Q$. Thus
$D=\diam Q$, $B(z_0,c_0D)\subset Q$, and
$\omega_Q(z)=D^{-2}\omega((z-z_0)/D)$ is the rescaled bump of
\eqref{eq:ray}, with $\omega\ge0$ smooth, of integral one and
supported in $B(0,c_0/4)$. Then $\int_Q\omega_Q=1$ and
$\|\nabla^j\omega_Q\|_\infty\le C_jD^{-2-j}$ for $j=0,1$.
All constants below may depend on $c_0$ and the fixed bump, but not
on $D/s$, the chart eccentricity, or the random realization.

\begin{lemma}\label{lem:app-ray}
For $x,y\in Q$, $x\ne y$, the vector kernel
\[
 K_Q(x,y)=(x-y)\int_1^\infty t\omega_Q(y+t(x-y))\,dt
\]
satisfies
\[
 |K_Q(x,y)|\le C|x-y|^{-1},\qquad
 |\nabla_xK_Q(x,y)|+|\nabla_yK_Q(x,y)|
 \le C|x-y|^{-2}.
\]
It vanishes when $|x-y|>C\operatorname{dist}(x,\partial Q)$.
\end{lemma}
\begin{proof}
Write $r=|x-y|$. On the support of the integrand,
$z=y+t(x-y)\in\operatorname{supp}\omega_Q\subset Q$,
so $tr=|z-y|\le D$. The integral therefore ends by $D/r$.
The bounds on the bump give
\[
 |K_Q(x,y)|\le CrD^{-2}\int_1^{D/r}t\,dt\le C/r.
\]
Differentiating the prefactor contributes at most
$CD^{-2}\int_1^{D/r}t\,dt\le C/r^2$.
Differentiating the bump in $x$ or $y$ contributes at most
$CrD^{-3}\int_1^{D/r}t^2\,dt\le C/r^2$.
These are the asserted gradient bounds.

For the support assertion, choose a supporting line $P$ for the convex domain which realizes
$\operatorname{dist}(x,\partial Q)$, with nonnegative inward signed
distance $\delta_P$. A point $z$ in the bump support satisfies
$\delta_P(z)\ge cD$ for fixed $c>0$, whereas $\delta_P(y)\ge0$.
As $x=(1-t^{-1})y+t^{-1}z$, one gets
$\operatorname{dist}(x,\partial Q)=\delta_P(x)\ge cD/t$.
Combining with $r\le D/t$ proves the support assertion.
\end{proof}

In the notation of Section~\ref{sec:geom}, Lemma~\ref{lem:app-ray}
is exactly \eqref{eq:singular} and \eqref{eq:support}.

\subsection{Finite Whitney geometry and actual-reveal indexing}
\label{app:allbranches}
We record the finite partition of a clipped parallelogram and the
indexing of actual refinements. The stopped scalar process itself is
treated in Proposition~\ref{prop:stopped}.
All assertions about disjointness of cells, pointwise values, and
Gaussian evaluations are understood up to Lebesgue-null cell
boundaries, with a fixed half-open dyadic convention.  Continuous
versions of the local Gaussian fields are used on the closures of
individual children.  The constants below never depend on the
physical eccentricity of the old metric.

\begin{lemma}
\label{lem:exactcover}
Let $Q\subset\mathbb R^2$ be convex with
$B(z_0,c_0D)\subset Q$, $D=\diam Q$, and let $\mathcal D_s$
be a translated, fixed, nested dyadic lattice of $s$-squares.
Fix $A>8/c_0$ and $\gamma>(A+\sqrt2)/2+2$.
For any $\theta\in(0,1/8)$ there is a constant
$C_A=C(c_0,A,\gamma,\theta)$ such that, if
$D>D_0s$ with $D_0>\max\{2A,4/\theta\}$, the following
algorithm gives a finite partition, modulo null boundaries.
\begin{enumerate}[label=(\roman*),leftmargin=2.5em]
\item Select the inclusion-maximal dyadic squares $B\subset Q$
whose side $2^js$ satisfies $2s\le 2^js\le\theta D$ and
$\dist(B,\partial Q)\ge\gamma 2^js$.  Keep them as individual
children.
\item Among the remaining $s$-squares, retain as \emph{receivers}
those which are wholly contained in $Q$.  All such squares
lie within $C_A s$ of $\partial Q$.
\item For every remaining $s$-square $R$ meeting both the interior
and the complement of $Q$, take the measurable clipped piece
$R\cap Q$ and attach it to one receiver within $C_A s$,
using deterministic lexicographic tie breaking.
\item Each receiver, together with all pieces assigned to it, is a
single terminal child.  A receiver with no assigned pieces is
also a terminal child.
\end{enumerate}
The selected large squares and resulting terminal children are
pairwise disjoint modulo null sets and cover $Q$ modulo null sets.
Every terminal child $S$ satisfies
\[
   s^2\le |S|\le C_A s^2,\qquad
   \diam S\le C_A s,\qquad
   \sup_{z\in S}\dist(z,\partial Q)\le C_A s.
\]
If $B$ is a selected large square, its side $\ell_B$ satisfies
$\ell_B\le\theta D$ and
$\gamma\ell_B\le\dist(z,\partial Q)\le C_A\ell_B$
for $z\in B$.  In particular the partition has the pointwise
Whitney comparisons needed for the projected kernel in
Lemma~\ref{lem:pk}.
\end{lemma}
\begin{proof}
Any $s$-square meeting $Q$ is either contained in a selected
large square, is a full uncovered square, or has nonempty
intersection with $Q$ and with its complement.  Half-open
boundaries make these three assignments disjoint modulo null
sets.  Inclusion-maximal squares in the same dyadic lattice
are disjoint or nested, so distinct selected squares are
disjoint.  The dyadic hierarchy is finite because its sides
are between $s$ and $\theta D$.

If a full uncovered $s$-square $R$ had an admissible $2s$
ancestor, that ancestor would lie in some selected square;
this contradicts uncoveredness.  Since $2s<\theta D$,
the $2s$ ancestor fails either containment in $Q$ or the
boundary-distance inequality.  In either case its points,
and therefore all of $R$, are within
$(2\gamma+3\sqrt2)s$ of $\partial Q$.  Thus there are no
uncovered squares far from the boundary, regardless of the
orientation or translation of the lattice.

For a clipped $R$, take $z\in R\cap Q$, so
$\dist(z,\partial Q)\le\sqrt2s$.  Set
$z_t=(1-t)z+tz_0$, $t=As/D<1/2$.  By the inball and
convexity, and then by the Lipschitz property of boundary
distance,
\[
 Ac_0s\le\dist(z_t,\partial Q)
 \le(A+\sqrt2)s,\qquad |z_t-z|\le As.
\]
Since $Ac_0>8$, the unique half-open $s$-square $R'$ with
$z_t\in R'$ is contained in $Q$.  Were $R'$ covered by a
selected dyadic square of side $L\ge2s$, then
$\dist(z_t,\partial Q)\ge\gamma L\ge2\gamma s$,
contrary to $2\gamma>A+\sqrt2$.  Thus $R'$ is uncovered
and is an eligible receiver.  Its distance from $R$ is at
most $(A+2\sqrt2)s$.  This construction works at an acute
corner even when $R$ meets two boundary sides; no side is
chosen, only the inball center.

Every piece assigned to a given receiver originates in one
of the at most $C(A)$ lattice squares in a fixed-radius
$O_A(s)$ neighborhood.  The receiver supplies area at least
$s^2$, and the union of at most $C(A)$ clipped pieces has
area at most $C_A s^2$ and diameter at most $C_A s$.
The boundary-distance assertion follows either from the
uncovered-full-square estimate or the clipped-square estimate.

Finally, for a selected $B$ of side $\ell$, its parent $B^\uparrow$
of side $2\ell$ is not admissible.  If it fails containment,
the distance to $\partial Q$ is at most $2\sqrt2\ell$ on $B$;
if it fails clearance, the upper bound is
$(2\gamma+2\sqrt2)\ell$; if it fails the cap, then
$D<2\ell/\theta$ and every point of $Q$ is within $D$
of the boundary.  The lower bound on $B$ is its
admissibility.  These are the stated uniform comparisons.
The proof is independent of the relative translation of the
lattice and of all angles except through the fixed inball
ratio $c_0$.
\end{proof}

\paragraph{Choice of constants.}
Fix the clipping constant $K>1$.  A square in its birth
metric transformed to the updated chart has singular values
between $K^{-1/2}$ and $K^{1/2}$; thus it has an inball
ratio at least $c_0=c_0(K)>0$.  Fix $A>8/c_0$,
then $\gamma>(A+\sqrt2)/2+2$.  These determine the terminal
union constant $C_A$, independently of $s$ and of
$D_0$.  Fix the contraction target $\vartheta$; then choose
$\theta<\vartheta/\sqrt{2K}$ and
\[
 D_0>\max\{2A,4/\theta,C_A\sqrt{2K}/\vartheta\}.
\]
Increasing $D_0$ never changes $C_A$ or the selected-square
comparisons.  Consequently Lemma~\ref{lem:birthchart}
applies uniformly to the exact finite partition, and by
reducing $\vartheta$ we enforce
$\vartheta\sqrt K<1$ and subsequently
$(\vartheta\sqrt K)^{2\beta}(1-\eta)^{-2a_*}<\rho_1$.

\begin{lemma}
\label{lem:actualreveal}
Let $Q_0\supset Q_1\supset\cdots\supset Q_b$
be the cells of the \emph{actual} revealed ancestral chain of
a point, where $Q_j$ is born with metric $G_j$, and each
revealed increment $V_j$ is routed in $G_j$.  Suppose every
strict transition from $Q_{p-1}$ to $Q_p$ satisfies the
comparison and side/diameter bounds of
Lemma~\ref{lem:birthchart}.  If $Q_b$ is a square or a
newly terminal union, then for every $1\le j\le b$,
\[
 \frac{\diam_{G_j}(Q_b)}{\ell_j}
  \le C(\vartheta\sqrt K)^{b-j},
\]
where $\ell_j$ is the side of $Q_j$ in its birth metric
when $Q_j$ is a square; for a newly terminal $Q_j$ its
last intrinsic diameter serves as the length at $j$.
If $Q_b$ is stopped directly, without a further reveal,
its $G_b$-diameter is at most $\sqrt K D_0s$.
At padded times $k>b$, $U_k=U_b$, $G_k=G_b$ and
$T_k=T_b$ on this leaf.  No new factor
$(\vartheta\sqrt K)^{k-b}$ is assigned to its geometry,
and no new Gaussian envelope is charged.
\end{lemma}
\begin{proof}
For a chain of square children, the intrinsic birth-side
contractions multiply to $\ell_b\le\vartheta^{b-j}\ell_j$.
As in the proof of Lemma~\ref{lem:ancestorindexed}, the Loewner
comparisons imply $G_j^{-1}\preceq K^{b-j}G_b^{-1}$ without any
commutativity assumption; this gives the asserted ancestor
diameter bound.
If the final child is a newly terminal union, apply the
terminal-diameter estimate only on the final transition and
use the ordinary side contractions earlier.  If the parent
stops without splitting, its new-chart diameter is
$\le D_0s$, while the old/new metric comparison gives
$\diam_{G_b}Q_b\le\sqrt K D_0s$.  Frozen padded times
have no strict transitions.
\end{proof}

\paragraph{Stopping convention.}
Lemma~\ref{lem:actualreveal} counts only actual refinements.  Once a
leaf is terminal, its field and memory are frozen; only the decreasing
clock continues to act on its exponential integrand.  At a failed
fresh-envelope test, the numerical modular is killed.  These are the
conventions of the stopped-supermartingale argument in
Proposition~\ref{prop:stopped}.  On successful completion, the terminal
partition and success event are measurable with respect to the revealed
history, so the unrevealed source is the conditional Gaussian bridge of
Lemma~\ref{lem:bridge}.

\section*{Use of AI tools}
\addcontentsline{toc}{section}{Use of AI tools}
The starting point of this work was the previous project \cite{umc:AGGT26} on
  the finite $p$ cost case. Intense interactions with LLMs spanning several weeks
   led to a first proof of the Ultimate matching conjecture in the subcritical case,
    leaving the critical case (as well a the endpoint cases)
     seemingly out of reach. Somewhat unexpectedly to us, a solution for the critical
      finite exponent case was finally found during another round of exploration, see the 'monolithic' file \texttt{finite-exponents.pdf} in the
       \href{https://github.com/grottofrancesco/UltimateMatchingTheorem.git}{repository}.
  After this GPT-6 quickly obtained the exact asymptotic for the endpoint cases, see the 'monolithic' file \texttt{endpoints.pdf} in the same repository.
  The content of this file was mostly generated and proof-checked by the LLMs during interactions with us aimed at understanding the 'monolithic' versions.
  We have an overall understanding of the proofs but plan to provide a fully digested version shortly.

\begin{thebibliography}{99}

\bibitem{umc:AKT84}
M.~Ajtai, J.~Koml\'os and G.~Tusn\'ady,
\emph{On optimal matchings},
Combinatorica \textbf{4} (1984), 259--264.

\bibitem{AGS}
L.~Ambrosio, N.~Gigli and G.~Savar\'e,
\emph{Gradient Flows in Metric Spaces and in the Space of Probability
Measures}, Birkh\"auser, 2008.

\bibitem{umc:AG19}
L.~Ambrosio and F.~Glaudo,
\emph{Finer estimates on the 2-dimensional matching problem},
J. \'Ec. polytech. Math. \textbf{6} (2019), 737--765.

\bibitem{umc:AGTmap19}
L.~Ambrosio, F.~Glaudo and D.~Trevisan,
\emph{On the optimal map in the 2-dimensional random matching problem},
Discrete Contin. Dyn. Syst. \textbf{39} (2019), 7291--7308.

\bibitem{umc:AGT22}
L.~Ambrosio, M.~Goldman and D.~Trevisan,
\emph{On the quadratic random matching problem in two-dimensional domains},
Electron. J. Probab. \textbf{27} (2022), paper no.~54.

\bibitem{umc:AST19}
L.~Ambrosio, F.~Stra and D.~Trevisan,
\emph{A PDE approach to a 2-dimensional matching problem},
Probab. Theory Related Fields \textbf{173} (2019), 433--477.

\bibitem{umc:AVT24}
L.~Ambrosio, F.~Vitillaro and D.~Trevisan,
\emph{Sharp PDE estimates for random two-dimensional bipartite matching
with power cost function}, Rend. Lincei Mat. Appl. \textbf{35} (2024),
323--342.

\bibitem{umc:AGGT26}
V.~Armegioiu, M.~Goldman, F.~Grotto and D.~Trevisan,
\emph{Exact asymptotics for the 2D Euclidean random matching problem},
arXiv:2609.08717 (2026).

\bibitem{BollobasLeader1991}
B.~Bollob\'as and I.~Leader,
\emph{Edge-isoperimetric inequalities in the grid},
Combinatorica \textbf{11} (1991), 299--314.

\bibitem{BretagnolleMassart1989}
J.~Bretagnolle and P.~Massart,
\emph{Hungarian constructions from the nonasymptotic viewpoint},
Ann. Probab. \textbf{17} (1989), 239--256.

\bibitem{umc:CLPS14}
S.~Caracciolo, C.~Lucibello, G.~Parisi and G.~Sicuro,
\emph{Scaling hypothesis for the Euclidean bipartite matching problem},
Phys. Rev. E \textbf{90} (2014), 012118.

\bibitem{umc:COPW26}
X.~Cheng, F.~Otto, M.~Palmieri and A.~Wachtel,
\emph{A geometrically linear approximation of min--max optimal matching},
arXiv:2609.08680 (2026).

\bibitem{DudleyUCLT}
R.~M.~Dudley, \emph{Uniform Central Limit Theorems}, 2nd ed.,
Cambridge University Press, 2014.

\bibitem{Gale1957}
D.~Gale, \emph{A theorem on flows in networks},
Pacific J. Math. \textbf{7} (1957), 1073--1082.

\bibitem{umc:GHO21}
M.~Goldman, M.~Huesmann and F.~Otto,
\emph{Quantitative linearization results for the Monge--Amp\`ere equation},
Comm. Pure Appl. Math. \textbf{74} (2021), 2483--2560.

\bibitem{umc:GHO23}
M.~Goldman, M.~Huesmann and F.~Otto,
\emph{Almost sharp rates of convergence for the average cost and
displacement in the optimal matching problem}, preprint, 2023;
Abel Symposium proceedings.

\bibitem{umc:GO20}
M.~Goldman and F.~Otto,
\emph{A variational proof of partial regularity for optimal
transportation maps}, Ann. Sci. \'Ec. Norm. Sup\'er. \textbf{53} (2020),
1209--1233.

\bibitem{GT-combinatorial}
M.~Goldman and D.~Trevisan,
\emph{Optimal transport methods for combinatorial optimization over two
random point sets}, Probab. Theory Related Fields \textbf{188} (2024),
1315--1384.

\bibitem{Hoffman1960}
A.~J.~Hoffman, \emph{Some recent applications of the theory of linear
inequalities to extremal combinatorial analysis},
in: Combinatorial Analysis, Proc. Sympos. Appl. Math. \textbf{10},
Amer. Math. Soc., Providence, 1960, 113--127.

\bibitem{KMT1975}
J.~Koml\'os, P.~Major and G.~Tusn\'ady,
\emph{An approximation of partial sums of independent RV's and the
sample DF.~I}, Z. Wahrscheinlichkeitstheorie verw. Gebiete \textbf{32}
(1975), 111--131.

\bibitem{umc:Ouimet26}
F.~Ouimet,
\emph{A solution to Problem 5.3 in Talagrand's My Favorite Problems:
simultaneous sum and bottleneck bounds for random planar matchings},
preprint (2026),
\href{https://doi.org/10.13140/RG.2.2.26527.93606/2}{doi:10.13140/RG.2.2.26527.93606/2}.

\bibitem{Santambrogio}
F.~Santambrogio, \emph{Optimal Transport for Applied Mathematicians},
Birkh\"auser, 2015.

\bibitem{umc:Talagrand92}
M.~Talagrand,
\emph{Matching random samples in many dimensions},
Ann. Appl. Probab. \textbf{2} (1992), 846--856.

\bibitem{umc:Talagrand2005}
M.~Talagrand,
\emph{The Generic Chaining: Upper and Lower Bounds of Stochastic
Processes}, Springer, 2005.

\bibitem{umc:Talagrand2021}
M.~Talagrand,
\emph{Upper and Lower Bounds for Stochastic Processes: Decomposition
Theorems}, 2nd ed., Springer, 2021.

\bibitem{TalagrandProblems}
M.~Talagrand, \emph{My favorite problems}, manuscript, 2026.

\bibitem{Villani}
C.~Villani, \emph{Optimal Transport: Old and New}, Springer, 2009.

\bibitem{umc:repository}
\url{https://github.com/grottofrancesco/UltimateMatchingTheorem}.

\end{thebibliography}


\end{document}
